% !TEX program = lualatex
% ============================================================
% Manual de Fundamentos de Matemáticas
% Esqueleto inicial para FdM1 y FdM2
% Filosofía del curso:
%   Todo símbolo relevante debe definirse antes de utilizarse.
% ============================================================

\documentclass[11pt,a4paper,openany]{book}

% ------------------------------------------------------------
% Codificación, idioma y tipografía
% ------------------------------------------------------------
\usepackage{fontspec}
\usepackage[spanish,es-nodecimaldot,es-tabla]{babel}
\usepackage{microtype}
\usepackage{csquotes}

\setmainfont{Latin Modern Roman}
\setsansfont{Latin Modern Sans}
\setmonofont{Latin Modern Mono}

% ------------------------------------------------------------
% Geometría y diseño general
% ------------------------------------------------------------
\usepackage[a4paper,margin=3cm,headheight=15pt]{geometry}
\usepackage{setspace}
\onehalfspacing
\usepackage{parskip}
\usepackage{fancyhdr}
\usepackage{xcolor}
\usepackage{titlesec}
\usepackage{enumitem}
\usepackage{booktabs}
\usepackage{array}
\usepackage{longtable}
\usepackage{tabularx}

% Ajuste del ancho reservado para numeros de capitulo en el indice.
\makeatletter
\renewcommand*{\l@chapter}[2]{%
  \ifnum \c@tocdepth >\m@ne
    \addpenalty{-\@highpenalty}%
    \vskip 1.0em \@plus\p@
    \setlength\@tempdima{3em}%
    \begingroup
      \parindent \z@ \rightskip \@pnumwidth
      \parfillskip -\@pnumwidth
      \leavevmode \bfseries
      \advance\leftskip\@tempdima
      \hskip -\leftskip
      #1\nobreak\hfil \nobreak\hb@xt@\@pnumwidth{\hss #2}\par
      \penalty\@highpenalty
    \endgroup
  \fi}
\renewcommand*{\l@section}{\@dottedtocline{1}{1.5em}{3em}}
\renewcommand*{\l@subsection}{\@dottedtocline{2}{4.5em}{3.8em}}
\makeatother

% ------------------------------------------------------------
% Matemáticas
% ------------------------------------------------------------
\usepackage{amsmath,amssymb,amsthm,mathtools}
\usepackage{thmtools}
\usepackage{bm}
\usepackage{mathrsfs}

% ------------------------------------------------------------
% Hiperenlaces y referencias
% ------------------------------------------------------------
\usepackage{hyperref}
\usepackage[nameinlink,capitalise,noabbrev]{cleveref}

\hypersetup{
  colorlinks=true,
  linkcolor=blue!50!black,
  citecolor=blue!50!black,
  urlcolor=blue!50!black,
  pdftitle={Fundamentos de Matemáticas},
  pdfauthor={Antonio Falcó Montesinos},
  pdfsubject={Manual de Fundamentos de Matemáticas},
  pdfkeywords={lógica, conjuntos, relaciones, estructuras algebraicas, números}
}

% ------------------------------------------------------------
% Colores institucionales discretos
% ------------------------------------------------------------
\definecolor{FdMBlue}{RGB}{20,55,100}
\definecolor{FdMGray}{RGB}{245,245,245}
\definecolor{FdMRed}{RGB}{130,30,30}
\definecolor{FdMGreen}{RGB}{25,100,60}

% ------------------------------------------------------------
% Encabezados
% ------------------------------------------------------------
\pagestyle{fancy}
\fancyhf{}
\fancyhead[LE,RO]{\thepage}
\fancyhead[LO]{\nouppercase{\rightmark}}
\fancyhead[RE]{\nouppercase{\leftmark}}

% ------------------------------------------------------------
% Formato de títulos
% ------------------------------------------------------------
\titleformat{\chapter}[display]
  {\normalfont\huge\bfseries\color{FdMBlue}}
  {\chaptertitlename\ \thechapter}
  {20pt}
  {\Huge}

\titleformat{\section}
  {\normalfont\Large\bfseries\color{FdMBlue}}
  {\thesection}{1em}{}

\titleformat{\subsection}
  {\normalfont\large\bfseries\color{FdMBlue}}
  {\thesubsection}{1em}{}

% ------------------------------------------------------------
% Entornos de teoremas
% ------------------------------------------------------------
% ------------------------------------------------------------
% Entornos de teoremas
% ------------------------------------------------------------
\declaretheoremstyle[
  headfont=\bfseries\color{FdMBlue},
  bodyfont=\normalfont,
  spaceabove=10pt,
  spacebelow=10pt,
  mdframed={
    nobreak=true,
    shadow=true,
    shadowsize=3pt,
    shadowcolor=black!35,
    backgroundcolor=FdMGray,
    linecolor=FdMBlue,
    linewidth=0.8pt,
    roundcorner=4pt,
    skipabove=10pt,
    skipbelow=10pt,
    innertopmargin=7pt,
    innerbottommargin=7pt,
    innerleftmargin=8pt,
    innerrightmargin=8pt
  }
]{fdmdefstyle}

\declaretheoremstyle[
  headfont=\bfseries\color{FdMRed},
  bodyfont=\normalfont,
  spaceabove=10pt,
  spacebelow=10pt,
  mdframed={
    nobreak=true,
    shadow=true,
    shadowsize=3pt,
    shadowcolor=black!35,
    backgroundcolor=FdMRed!4!white,
    linecolor=FdMRed,
    linewidth=0.8pt,
    roundcorner=4pt,
    skipabove=10pt,
    skipbelow=10pt,
    innertopmargin=7pt,
    innerbottommargin=7pt,
    innerleftmargin=8pt,
    innerrightmargin=8pt
  }
]{fdmthmstyle}

\declaretheoremstyle[
  headfont=\bfseries\color{FdMGreen},
  bodyfont=\normalfont,
  spaceabove=10pt,
  spacebelow=10pt,
  mdframed={
    nobreak=true,
    shadow=true,
    shadowsize=3pt,
    shadowcolor=black!35,
    backgroundcolor=FdMGreen!4!white,
    linecolor=FdMGreen,
    linewidth=0.8pt,
    roundcorner=4pt,
    skipabove=10pt,
    skipbelow=10pt,
    innertopmargin=7pt,
    innerbottommargin=7pt,
    innerleftmargin=8pt,
    innerrightmargin=8pt
  }
]{fdmexstyle}


\declaretheorem[style=fdmdefstyle,numberwithin=chapter,name=Definición]{definicion}
\declaretheorem[style=fdmthmstyle,sibling=definicion,name=Teorema]{teorema}
\declaretheorem[style=fdmthmstyle,sibling=definicion,name=Proposición]{proposicion}
\declaretheorem[style=fdmthmstyle,sibling=definicion,name=Lema]{lema}
\declaretheorem[style=fdmthmstyle,sibling=definicion,name=Corolario]{corolario}
\declaretheorem[style=fdmexstyle,sibling=definicion,name=Ejemplo]{ejemplo}
\declaretheorem[style=fdmexstyle,sibling=definicion,name=Contraejemplo]{contraejemplo}
\declaretheorem[style=fdmdefstyle,sibling=definicion,name=Método]{metodo}
\declaretheorem[style=fdmdefstyle,sibling=definicion,name=Advertencia]{advertencia}
\declaretheorem[style=fdmdefstyle,sibling=definicion,name=Convención]{convencion}

\renewcommand{\theHdefinicion}{\thechapter.\arabic{definicion}}
\renewcommand{\theHteorema}{\theHdefinicion}
\renewcommand{\theHproposicion}{\theHdefinicion}
\renewcommand{\theHlema}{\theHdefinicion}
\renewcommand{\theHcorolario}{\theHdefinicion}
\renewcommand{\theHejemplo}{\theHdefinicion}
\renewcommand{\theHcontraejemplo}{\theHdefinicion}
\renewcommand{\theHmetodo}{\theHdefinicion}
\renewcommand{\theHadvertencia}{\theHdefinicion}
\renewcommand{\theHconvencion}{\theHdefinicion}

% ------------------------------------------------------------
% Entornos pedagógicos
% ------------------------------------------------------------
\newenvironment{objetivos}
{\section*{Objetivos del capítulo}\phantomsection\addcontentsline{toc}{section}{Objetivos del capítulo}\begin{itemize}[leftmargin=2em]}
{\end{itemize}}

\newenvironment{seminario}
{\section*{Seminario asociado}\phantomsection\addcontentsline{toc}{section}{Seminario asociado}\begin{itemize}[leftmargin=2em]}
{\end{itemize}}

\newenvironment{erroresfrecuentes}
{\section*{Errores frecuentes}\phantomsection\addcontentsline{toc}{section}{Errores frecuentes}\begin{itemize}[leftmargin=2em]}
{\end{itemize}}

\newenvironment{ejerciciospropuestos}
{\section*{Ejercicios propuestos}\phantomsection\addcontentsline{toc}{section}{Ejercicios propuestos}\begin{enumerate}[leftmargin=2em,label=\textbf{\thechapter.\arabic*.}]}
{\end{enumerate}}

% Citas textuales destacadas dentro del texto.
\newenvironment{cita}
{\begin{quote}\small\itshape}
{\end{quote}}

\newcommand{\resumen}[1]{%
  \section*{Resumen del capítulo}%
  \phantomsection%
  \addcontentsline{toc}{section}{Resumen del capítulo}%
  #1
}

% ------------------------------------------------------------
% Comandos de conjuntos numéricos
% Cada símbolo se introduce explícitamente en el texto antes de su uso.
% Estos comandos solo fijan una escritura uniforme.
% ------------------------------------------------------------
\newcommand{\N}{\mathbb{N}}
\newcommand{\Z}{\mathbb{Z}}
\newcommand{\Q}{\mathbb{Q}}
\newcommand{\R}{\mathbb{R}}
\newcommand{\C}{\mathbb{C}}

\newcommand{\Pow}{\mathcal{P}}
\newcommand{\id}{\operatorname{id}}
\newcommand{\mcd}{\operatorname{mcd}}
\newcommand{\mcm}{\operatorname{mcm}}
\newcommand{\im}{\operatorname{Im}}

% ------------------------------------------------------------
% Datos del documento
% ------------------------------------------------------------
\title{\Huge Fundamentos de Matemáticas\\[0.4cm]
\Large Manual para FdM1 y FdM2}
\author{Antonio Falcó Montesinos\\Universidad CEU Cardenal Herrera}
\date{Curso académico 2026--2027}

\begin{document}

% ------------------------------------------------------------
% Portada
% ------------------------------------------------------------
\frontmatter
\maketitle

\chapter*{Principio rector del manual}
\addcontentsline{toc}{chapter}{Principio rector del manual}

La filosofía fundamental de este manual es la siguiente:

\begin{center}
\fbox{\parbox{0.85\textwidth}{\centering\Large
Todo símbolo relevante debe definirse antes de utilizarse.}}
\end{center}

Este principio no es una mera regla tipográfica. Es una norma de rigor matemático y de responsabilidad expositiva. En particular, a lo largo del texto se seguirá la siguiente disciplina:

\begin{enumerate}[label=\textbf{R\arabic*.},leftmargin=2.5em]
  \item Antes de emplear un símbolo matemático, se indicará qué objeto representa.
  \item Antes de utilizar una operación, se precisará sobre qué conjunto está definida.
  \item Antes de enunciar una propiedad, se especificará el contexto en el que dicha propiedad tiene sentido.
  \item Antes de usar una abreviatura o notación especializada, se explicará su significado.
  \item Los ejemplos se presentarán después de la definición correspondiente, no antes, salvo que tengan una función motivadora explícita.
  \item Toda demostración deberá indicar claramente cuáles son las hipótesis, cuál es la conclusión y qué método de razonamiento se está empleando.
\end{enumerate}

Esta norma será especialmente importante en los primeros capítulos, donde el estudiante debe aprender no solo nuevos contenidos, sino también una forma rigurosa de escribir y leer matemáticas.

\chapter*{Prólogo}
\addcontentsline{toc}{chapter}{Prólogo}

Estas notas constituyen una introducción a los fundamentos del razonamiento matemático y a la construcción de los sistemas numéricos básicos. Su objetivo principal no es únicamente presentar una colección de conceptos, definiciones y resultados, sino ayudar al estudiante a adquirir una forma de pensar propia de la matemática universitaria.

A lo largo del curso se insistirá en tres ideas fundamentales. La primera es que la matemática posee un lenguaje propio, preciso y estructurado. La segunda es que los objetos matemáticos no se estudian de manera aislada, sino a través de sus propiedades, relaciones y estructuras. La tercera es que una afirmación matemática requiere una justificación rigurosa: ejemplos, cálculos y figuras pueden orientar la intuición, pero no sustituyen a una demostración.

El curso se divide en dos partes. En Fundamentos de Matemáticas I se introducen la lógica elemental, el razonamiento matemático, la teoría básica de conjuntos, las aplicaciones, las relaciones binarias, las primeras estructuras algebraicas y los números naturales. En Fundamentos de Matemáticas II se estudian los números enteros, la divisibilidad, el Teorema Fundamental de la Aritmética, la aritmética modular, los números racionales, los números reales y los números complejos.

Las clases magistrales estarán dedicadas a la presentación ordenada de los conceptos y resultados principales. Los seminarios tendrán un papel igualmente importante: en ellos se trabajará la resolución de problemas, la escritura de demostraciones, la construcción de ejemplos y contraejemplos, y la discusión de errores frecuentes.

\tableofcontents

% ------------------------------------------------------------
% Estructura académica
% ------------------------------------------------------------
\chapter*{Organización académica del curso}
\addcontentsline{toc}{chapter}{Organización académica del curso}

El curso completo de Fundamentos de Matemáticas se organiza en dos asignaturas semestrales:

\begin{center}
\begin{tabularx}{0.95\textwidth}{>{\bfseries}l X c c}
\toprule
Asignatura & Contenido general & Clases magistrales & Seminarios \\
\midrule
FdM1 & Lógica, conjuntos, relaciones, estructuras algebraicas y números naturales & 15 sesiones & 15 sesiones \\
FdM2 & Enteros, divisibilidad, aritmética modular, racionales, reales y complejos & 15 sesiones & 15 sesiones \\
\bottomrule
\end{tabularx}
\end{center}

Cada capítulo del manual se estructurará, en la medida de lo posible, mediante los siguientes elementos:

\begin{enumerate}[leftmargin=2.5em]
  \item objetivos del capítulo;
  \item motivación;
  \item definiciones fundamentales;
  \item ejemplos y contraejemplos;
  \item resultados principales;
  \item demostraciones modelo;
  \item errores frecuentes;
  \item seminario asociado;
  \item ejercicios propuestos;
  \item resumen del capítulo.
\end{enumerate}

% ============================================================
% FdM1
% ============================================================
\mainmatter
\part{Fundamentos de Matemáticas I}

% ------------------------------------------------------------
% ============================================================================
% Capítulo 1
% Fundamentos de Matemáticas I
% Fichero independiente para incluir con: % ============================================================================
% Capítulo 1
% Fundamentos de Matemáticas I
% Fichero independiente para incluir con: % ============================================================================
% Capítulo 1
% Fundamentos de Matemáticas I
% Fichero independiente para incluir con: \input{capitulo1.tex}
% ============================================================================

\chapter{El lenguaje de las matemáticas}
\label{chap:lenguaje-matematicas}

\begin{objetivos}
  \item Comprender la diferencia entre lenguaje natural y lenguaje matemático.
  \item Identificar definiciones, enunciados, ejemplos, contraejemplos y demostraciones.
  \item Reconocer que cada símbolo relevante debe ser definido antes de utilizarse.
  \item Distinguir entre calcular, justificar, argumentar y demostrar.
\end{objetivos}

\section{Motivación}
\label{sec:motivacion-lenguaje}

Este capítulo tiene como finalidad introducir el modo en que se escribe, se lee y se razona en matemáticas. No se trata todavía de desarrollar una teoría específica, sino de fijar las reglas básicas del lenguaje matemático que se utilizarán durante todo el curso.

En matemáticas no basta con tener una intuición correcta. También es necesario expresar esa intuición mediante un lenguaje preciso. Una misma frase del lenguaje natural puede admitir varias interpretaciones, mientras que una afirmación matemática debe tener un significado claramente determinado dentro del contexto en el que aparece.

Por ejemplo, la frase ``un número grande'' no tiene un significado matemático preciso si no se especifica qué se entiende por número, respecto de qué comparación se considera grande y en qué conjunto se está trabajando. En cambio, la afirmación
\[
  n > 10,
\]
tiene un significado preciso si previamente se ha indicado que \(n\) es un número natural, que \(10\) es el número natural diez y que el símbolo \(>\) denota la relación usual de orden en el conjunto de los números naturales.

Esta observación conduce a una regla fundamental del curso:

\begin{convencion}[Principio de definición previa]
Todo símbolo relevante debe definirse antes de utilizarse.
\end{convencion}

Este principio no es una exigencia meramente formal. Es una condición necesaria para que un texto matemático sea comprensible, verificable y riguroso. Durante el curso se insistirá de forma sistemática en la identificación de símbolos no definidos, en la escritura correcta de definiciones y en la formulación precisa de enunciados.

\section{Lenguaje natural y lenguaje matemático}
\label{sec:lenguaje-natural-matematico}

El lenguaje natural es el lenguaje que utilizamos habitualmente para comunicarnos. Es flexible, admite matices, depende del contexto y puede contener ambigüedades. Esta flexibilidad es útil en la comunicación ordinaria, pero puede ser problemática en matemáticas.

El lenguaje matemático, en cambio, busca reducir la ambigüedad. Para ello utiliza símbolos, definiciones, reglas lógicas y convenciones de notación. Su finalidad no es sustituir por completo al lenguaje natural, sino complementarlo con una estructura precisa.

\begin{definicion}[Lenguaje matemático]
Llamaremos \emph{lenguaje matemático} al sistema formado por palabras, símbolos, reglas de escritura y reglas de interpretación que se emplean para definir objetos matemáticos, formular enunciados y construir demostraciones.
\end{definicion}

El lenguaje matemático combina dos componentes:

\begin{itemize}
  \item una parte verbal, escrita en lenguaje natural;
  \item una parte simbólica, formada por signos matemáticos como \(=\), \(<\), \(\in\), \(\subseteq\), \(+\), \(\forall\), \(\exists\), entre otros.
\end{itemize}

Ninguna de estas partes debe aparecer de manera aislada. Un texto matemático correcto debe explicar qué significan los símbolos que utiliza y debe emplear las palabras con un sentido compatible con las definiciones dadas.

\begin{ejemplo}[Ambigüedad en lenguaje natural]
La frase ``todos los números tienen un siguiente'' puede interpretarse de varias formas. Para que sea una afirmación matemática precisa, debemos indicar el conjunto de números al que nos referimos.

Si trabajamos en el conjunto \(\mathbb{N}\) de los números naturales, y si por ``siguiente'' entendemos el número obtenido al sumar \(1\), entonces la afirmación puede escribirse como:
\[
  \text{para todo } n\in \mathbb{N}, \text{ existe } m\in \mathbb{N} \text{ tal que } m=n+1.
\]
En esta formulación aparecen símbolos como \(\mathbb{N}\), \(\in\), \(m\), \(n\), \(+\) y \(=\), cuyo significado debe haber sido definido o aceptado previamente como parte de la notación del curso.
\end{ejemplo}

\begin{advertencia}
Una expresión simbólica no es necesariamente más rigurosa por el simple hecho de contener símbolos. Una fórmula con símbolos no definidos puede ser menos clara que una frase escrita cuidadosamente en lenguaje natural.
\end{advertencia}

% --------------------------------------------------------------
%  Bloque nuevo: Perspectiva histórica y formal del lenguaje
% --------------------------------------------------------------

\section{Perspectiva histórica y formal del lenguaje}
\label{sec:perspectiva-bourbaki}

Bourbaki, en su obra \emph{Éléments de mathématique}~\cite{bourbaki1970},
remarca que la matemática, desde los griegos, se caracteriza por la
\textbf{demostración rigurosa}.  Una de sus observaciones más importantes es:

\begin{quote}
  « … el armazón lógico de cada capítulo está constituido por definiciones,
  axiomas y teoremas.  El lenguaje debe ser rigurosamente correcto;
  los abusos de lenguaje o de notación, sin los cuales cualquier texto
  matemático se vuelve pedante e ilegible, han sido señalados. »
  \hfill \textit{(Éléments de mathématique, p.~6)}
\end{quote}

A partir de esta reflexión, Bourbaki extrae dos principios que encajan
perfectamente con los que hemos adoptado en el curso:

\begin{convencion}[Principio de definición previa (Bourbaki)]
  Todo símbolo relevante debe definirse antes de utilizarse.
  Este principio asegura que el lector pueda reconstruir el sentido de
  cada paso sin ambigüedades.
\end{convencion}

\begin{convencion}[Método axiomático]
  La escritura de un texto matemático debe seguir el esquema
  ``definiciones → axiomas → teoremas''; de esta forma el texto es
  \emph{fácil de formalizar}.
\end{convencion}

\begin{cita}
  « La verificación de un texto formalizado sólo requiere una atención en
  cierto modo mecánica; las únicas fuentes de error provienen de la
  longitud o la complejidad del texto. » 
  \hfill \textit{(Éléments de mathématique, p.~8)}
\end{cita}

Estos principios refuerzan la \hyperref[sec:lenguaje-natural-matematico]{sección~\ref*{sec:lenguaje-natural-matematico}}:
el lenguaje natural sirve de base, pero la exactitud sólo se consigue
cuando cada símbolo y cada regla están claramente especificados.


\section{Objetos, símbolos y significado}
\label{sec:objetos-simbolos-significado}

La matemática estudia objetos. Estos objetos pueden ser números, puntos, conjuntos, funciones, relaciones, matrices, espacios, grafos u otras estructuras. Para referirnos a ellos utilizamos símbolos.

\begin{definicion}[Objeto matemático]
Un \emph{objeto matemático} es una entidad definida dentro de un contexto matemático. Puede ser un número, un conjunto, una función, una relación, una operación o una estructura formada por varios objetos relacionados entre sí.
\end{definicion}

\begin{definicion}[Símbolo]
Un \emph{símbolo} es un signo que se utiliza para representar un objeto, una relación, una operación o una propiedad matemática.
\end{definicion}

Por ejemplo, en la expresión
\[
  x \in A,
\]
el símbolo \(x\) suele representar un objeto, el símbolo \(A\) suele representar un conjunto y el símbolo \(\in\) representa la relación de pertenencia. La expresión se lee: ``\(x\) pertenece a \(A\)''.

Sin embargo, esta lectura solo es legítima si antes se ha explicado qué representa \(x\), qué representa \(A\) y qué significa \(\in\). Por esta razón adoptamos la siguiente convención.

\begin{convencion}[Definición previa de los símbolos]
Siempre que se introduzca un símbolo nuevo, se indicará explícitamente qué objeto representa. Por ejemplo, antes de escribir una expresión como \(x\in A\), se deberá haber explicado qué es \(x\), qué es \(A\) y qué significa la relación \(\in\).
\end{convencion}

\begin{ejemplo}[Introducción correcta de símbolos]
Una forma correcta de introducir la expresión \(x\in A\) es la siguiente:

Sea \(A\) un conjunto y sea \(x\) un objeto. Escribiremos \(x\in A\) para indicar que \(x\) pertenece al conjunto \(A\).
\end{ejemplo}

\begin{ejemplo}[Uso incorrecto de símbolos]
La frase siguiente no es adecuada al comienzo de un texto:
\[
  x\in A \Rightarrow x\in B.
\]
No se ha indicado qué son \(A\), \(B\), \(x\), ni qué significa el símbolo \(\Rightarrow\). Aunque un lector experto pueda adivinar la intención, el texto no cumple el principio de definición previa.
\end{ejemplo}

\subsection{Variables y constantes}
\label{subsec:variables-constantes}

En matemáticas es importante distinguir entre símbolos que representan objetos fijos y símbolos que pueden representar objetos variables.

\begin{definicion}[Variable]
Una \emph{variable} es un símbolo que puede representar distintos objetos dentro de un conjunto o contexto previamente especificado.
\end{definicion}

\begin{definicion}[Constante]
Una \emph{constante} es un símbolo que representa un objeto fijo dentro del contexto considerado.
\end{definicion}

Por ejemplo, si escribimos ``sea \(n\in \mathbb{N}\)'', el símbolo \(n\) se utiliza como variable que representa un número natural arbitrario. En cambio, el símbolo \(0\), una vez definido, representa un número natural concreto.

\begin{advertencia}
El mismo símbolo puede tener significados distintos en textos distintos. Por ejemplo, \(n\) puede representar un número natural en un problema, una dimensión en otro y un índice en otro. Por eso es necesario declarar su significado cada vez que se inicia un nuevo contexto.
\end{advertencia}

\section{Definiciones matemáticas}
\label{sec:definiciones-matematicas}

Una definición introduce un objeto, una propiedad o una relación. En matemáticas, definir no significa describir de manera aproximada, sino fijar con precisión el significado de un término o símbolo.

\begin{definicion}[Definición matemática]
Una \emph{definición matemática} es una declaración que asigna un significado preciso a un término, símbolo, objeto, propiedad o relación dentro de un contexto determinado.
\end{definicion}

Las definiciones son acuerdos de significado. No son verdaderas ni falsas en el mismo sentido que un teorema. Una vez aceptada una definición, podemos utilizarla para formular enunciados y demostrar resultados.

\begin{ejemplo}[Definición de número par]
Sea \(n\) un número entero. Diremos que \(n\) es \emph{par} si existe un número entero \(k\) tal que
\[
  n = 2k.
\]
En esta definición se debe haber indicado previamente qué es un número entero y qué significan los símbolos \(=\), \(2\) y el producto \(2k\).
\end{ejemplo}

\begin{ejemplo}[Definición de subconjunto]
Sean \(A\) y \(B\) dos conjuntos. Diremos que \(A\) es un \emph{subconjunto} de \(B\), y escribiremos
\[
  A\subseteq B,
\]
si todo elemento de \(A\) es también un elemento de \(B\).
\end{ejemplo}

Una buena definición debe cumplir varias condiciones:

\begin{itemize}
  \item debe indicar el contexto en el que se trabaja;
  \item debe introducir claramente los símbolos nuevos;
  \item debe evitar circularidades;
  \item debe permitir decidir, al menos en los casos básicos, si un objeto satisface o no la propiedad definida;
  \item debe ser compatible con el uso posterior que se hará del concepto.
\end{itemize}

\begin{advertencia}[Definiciones circulares]
Una definición es circular cuando utiliza el término que pretende definir. Por ejemplo, decir que ``un conjunto finito es un conjunto que tiene finitos elementos'' no es una definición satisfactoria si no se ha definido antes qué significa ``finito''.
\end{advertencia}

\section{Enunciados matemáticos}
\label{sec:enunciados-matematicos}

No toda frase que contiene palabras matemáticas es un enunciado matemático. Para que una frase sea un enunciado matemático debe afirmar algo susceptible de ser verdadero o falso dentro de un contexto determinado.

\begin{definicion}[Enunciado matemático]
Un \emph{enunciado matemático} es una afirmación formulada con precisión que, dentro de un contexto matemático dado, tiene un valor de verdad: verdadero o falso.
\end{definicion}

\begin{ejemplo}[Enunciados matemáticos]
Las siguientes frases son enunciados matemáticos, siempre que se hayan definido los símbolos implicados:
\begin{enumerate}
  \item El número \(2\) es par.
  \item Para todo número natural \(n\), se cumple \(n+1>n\).
  \item Existe un número entero \(k\) tal que \(k^2=4\).
\end{enumerate}
\end{ejemplo}

\begin{ejemplo}[Frases que no son enunciados matemáticos]
Las siguientes frases no son enunciados matemáticos en sentido estricto:
\begin{enumerate}
  \item ``Calcula \(2+3\)''. Es una orden, no una afirmación.
  \item ``\(x+1\)''. Es una expresión, no una afirmación.
  \item ``Los números son interesantes''. Es una valoración, no una afirmación matemática precisa.
\end{enumerate}
\end{ejemplo}

Los enunciados matemáticos pueden adoptar distintas formas. Algunas de las más frecuentes son:

\begin{itemize}
  \item afirmaciones universales: ``para todo objeto se cumple una propiedad'';
  \item afirmaciones existenciales: ``existe al menos un objeto que cumple una propiedad'';
  \item implicaciones: ``si se cumple una hipótesis, entonces se cumple una conclusión'';
  \item equivalencias: ``una propiedad se cumple si y solo si se cumple otra''.
\end{itemize}

\begin{advertencia}
En capítulos posteriores se estudiarán con detalle los conectores lógicos y los cuantificadores. En este capítulo basta con comprender que la forma lógica de un enunciado influye en el tipo de demostración que se debe realizar.
\end{advertencia}

\section{Ejemplos y contraejemplos}
\label{sec:ejemplos-contraejemplos}

Los ejemplos y contraejemplos son herramientas esenciales para comprender una definición y para analizar la validez de un enunciado.

\begin{definicion}[Ejemplo]
Un \emph{ejemplo} de una definición o propiedad es un objeto concreto que satisface dicha definición o propiedad.
\end{definicion}

\begin{definicion}[Contraejemplo]
Un \emph{contraejemplo} a un enunciado universal es un objeto concreto que cumple las hipótesis del enunciado pero no cumple su conclusión.
\end{definicion}

\begin{ejemplo}[Ejemplo de número par]
El número \(8\) es par, porque existe un número entero \(k\), concretamente \(k=4\), tal que
\[
  8 = 2k.
\]
\end{ejemplo}

\begin{ejemplo}[Contraejemplo]
Considérese la afirmación:

``Todo número entero positivo es par''.

Esta afirmación es falsa. Un contraejemplo es el número \(3\), porque \(3\) es un número entero positivo y no existe ningún número entero \(k\) tal que \(3=2k\).
\end{ejemplo}

Un solo contraejemplo basta para demostrar que un enunciado universal es falso. En cambio, muchos ejemplos no bastan para demostrar que un enunciado universal es verdadero.

\begin{advertencia}
Comprobar muchos casos particulares puede sugerir que una afirmación es verdadera, pero no constituye una demostración general. La matemática exige justificar por qué el resultado se cumple en todos los casos contemplados por el enunciado.
\end{advertencia}

\section{Demostrar y calcular}
\label{sec:demostrar-calcular}

Una parte importante del aprendizaje matemático consiste en distinguir entre calcular, justificar, argumentar y demostrar. Estas actividades están relacionadas, pero no son equivalentes.

\begin{definicion}[Cálculo]
Un \emph{cálculo} es una manipulación de expresiones mediante reglas previamente aceptadas, con el fin de obtener un resultado concreto.
\end{definicion}

\begin{definicion}[Justificación]
Una \emph{justificación} es una explicación que indica por qué un paso, una afirmación o un cálculo es válido.
\end{definicion}

\begin{definicion}[Argumento]
Un \emph{argumento} es una cadena de razones que pretende apoyar una conclusión.
\end{definicion}

\begin{definicion}[Demostración]
Una \emph{demostración} es un argumento matemático completo y riguroso que establece la verdad de un enunciado a partir de definiciones, resultados previamente demostrados y reglas lógicas aceptadas.
\end{definicion}

\begin{ejemplo}[Cálculo sin demostración]
La igualdad
\[
  2+3=5
\]
es un cálculo elemental. Pero si el objetivo fuera demostrar una propiedad general, por ejemplo
\[
  n+0=n \quad \text{para todo } n\in\mathbb{N},
\]
no bastaría con comprobar algunos casos particulares como \(1+0=1\), \(2+0=2\) y \(3+0=3\). Sería necesario dar una demostración general.
\end{ejemplo}

\begin{ejemplo}[Demostración sencilla]
Demostremos que la suma de dos números pares es par.

Sean \(a\) y \(b\) dos números enteros pares. Por definición de número par, existen dos números enteros \(r\) y \(s\) tales que
\[
  a=2r
  \quad \text{y} \quad
  b=2s.
\]
Entonces
\[
  a+b = 2r+2s = 2(r+s).
\]
Como \(r+s\) es un número entero, se deduce que \(a+b\) es par. Por tanto, la suma de dos números pares es par.
\end{ejemplo}

En esta demostración se observa una estructura típica:

\begin{enumerate}
  \item se declaran los objetos con los que se trabaja;
  \item se usan las definiciones pertinentes;
  \item se realiza un cálculo justificado;
  \item se interpreta el resultado obtenido;
  \item se formula la conclusión.
\end{enumerate}

\begin{convencion}[Estructura básica de una demostración]
Siempre que sea posible, una demostración deberá indicar explícitamente qué se supone, qué se quiere probar, qué definiciones se utilizan y cómo se obtiene la conclusión.
\end{convencion}

\section{Lectura y escritura de textos matemáticos}
\label{sec:lectura-escritura-textos}

Leer matemáticas no consiste únicamente en pasar la vista por una sucesión de símbolos. Es necesario reconstruir el significado de cada frase, comprobar qué objetos han sido definidos y entender qué papel desempeña cada hipótesis.

Al leer un texto matemático conviene hacerse las siguientes preguntas:

\begin{itemize}
  \item ¿Cuál es el contexto en el que se trabaja?
  \item ¿Qué objetos se han definido?
  \item ¿Qué símbolos aparecen y qué significa cada uno?
  \item ¿Cuál es la afirmación principal?
  \item ¿Qué hipótesis se utilizan?
  \item ¿Qué se quiere demostrar?
  \item ¿Qué definiciones o resultados previos se están aplicando?
\end{itemize}

Escribir matemáticas exige el mismo cuidado. Un texto matemático debe ser legible, preciso y ordenado. La claridad no está reñida con el rigor; al contrario, el rigor suele exigir claridad.

\subsection{Plantilla básica para escribir una demostración}
\label{subsec:plantilla-demostracion}

Una demostración elemental puede organizarse de acuerdo con la siguiente plantilla:

\begin{enumerate}
  \item \textbf{Declaración de objetos.} Indicar qué objetos se consideran y en qué conjunto o contexto viven.
  \item \textbf{Hipótesis.} Escribir con precisión qué se supone.
  \item \textbf{Objetivo.} Indicar qué se quiere probar.
  \item \textbf{Uso de definiciones.} Traducir las hipótesis mediante las definiciones correspondientes.
  \item \textbf{Desarrollo.} Encadenar igualdades, implicaciones o argumentos justificados.
  \item \textbf{Conclusión.} Explicar por qué lo obtenido prueba exactamente lo que se quería demostrar.
\end{enumerate}

\begin{ejemplo}[Reescritura rigurosa de un enunciado]
La frase ``si un número acaba en cero, entonces es divisible'' es ambigua. No se especifica qué tipo de número se considera ni por qué número debe ser divisible.

Una formulación más precisa sería:

Sea \(n\) un número entero escrito en base decimal. Si la última cifra decimal de \(n\) es \(0\), entonces \(n\) es divisible por \(10\).
\end{ejemplo}

\begin{advertencia}[Símbolos introducidos demasiado tarde]
No es recomendable escribir primero una fórmula y explicar después qué significan sus símbolos. En un texto bien organizado, los símbolos relevantes se presentan antes de ser utilizados.
\end{advertencia}

\section{Errores frecuentes}
\label{sec:errores-frecuentes-lenguaje}

En los primeros cursos universitarios aparecen algunos errores recurrentes relacionados con el lenguaje matemático. Conviene identificarlos desde el principio.

\begin{enumerate}
  \item \textbf{Usar símbolos sin definirlos.} Por ejemplo, escribir \(x\in A\) sin haber explicado qué son \(x\), \(A\) y \(\in\).

  \item \textbf{Confundir ejemplo con demostración.} Verificar un resultado para varios casos particulares no demuestra que sea verdadero en general.

  \item \textbf{Confundir definición con teorema.} Una definición introduce significado; un teorema afirma algo que debe demostrarse.

  \item \textbf{Omitir el contexto.} No es lo mismo trabajar en \(\mathbb{N}\), en \(\mathbb{Z}\), en \(\mathbb{Q}\), en \(\mathbb{R}\) o en \(\mathbb{C}\).

  \item \textbf{Utilizar flechas de manera imprecisa.} Símbolos como \(\Rightarrow\), \(\Leftrightarrow\) o \(\mapsto\) tienen significados concretos y no deben utilizarse como simples signos decorativos.

  \item \textbf{No distinguir entre objeto y propiedad.} Por ejemplo, un número puede tener la propiedad de ser par, pero ``ser par'' no es un número.

  \item \textbf{No cerrar la demostración.} Una demostración debe terminar indicando claramente que se ha probado la afirmación deseada.
\end{enumerate}

\begin{convencion}[Control de calidad de un texto matemático]
Antes de dar por correcto un texto matemático, comprobaremos que todos los símbolos relevantes han sido definidos, que las afirmaciones tienen sentido en el contexto fijado y que cada conclusión está justificada.
\end{convencion}

\begin{seminario}
  \item Identificación de símbolos no definidos en textos matemáticos breves.
  \item Reescritura rigurosa de enunciados ambiguos.
  \item Construcción de ejemplos y contraejemplos elementales.
  \item Distinción entre definición, enunciado, ejemplo, contraejemplo y demostración.
  \item Discusión de demostraciones incompletas y detección de pasos no justificados.
\end{seminario}

\begin{ejerciciospropuestos}
  \item Localiza en un texto matemático breve todos los símbolos que aparecen y escribe una lista con su significado.

  \item Escribe tres ejemplos de afirmaciones matemáticas y tres ejemplos de frases que no sean enunciados matemáticos.

  \item Reescribe de forma rigurosa las siguientes frases:
  \begin{enumerate}
    \item ``Todo número tiene un siguiente.''
    \item ``Si un número es grande, su cuadrado también lo es.''
    \item ``Los números negativos al multiplicarse dan positivo.''
  \end{enumerate}

  \item Indica cuáles de las siguientes frases son definiciones, cuáles son enunciados y cuáles no son afirmaciones matemáticas:
  \begin{enumerate}
    \item ``Un número entero \(n\) es par si existe un entero \(k\) tal que \(n=2k\).''
    \item ``Todo número primo mayor que \(2\) es impar.''
    \item ``Calcula \(7+5\).''
    \item ``Sea \(A\) un conjunto.''
  \end{enumerate}

  \item Da un contraejemplo para cada una de las siguientes afirmaciones falsas:
  \begin{enumerate}
    \item Todo número entero es positivo.
    \item Todo número impar es primo.
    \item Si \(a^2=b^2\), entonces \(a=b\), donde \(a\) y \(b\) son números enteros.
  \end{enumerate}

  \item Explica por qué comprobar que \(1+3\), \(3+5\) y \(5+7\) son números pares no demuestra por sí solo que la suma de dos números impares cualesquiera sea par.

  \item Escribe una demostración completa de la siguiente afirmación: la suma de dos números impares es par. Antes de empezar, define qué significa que un número entero sea impar.

  \item Revisa una solución propia de un ejercicio matemático e identifica todos los símbolos que has utilizado. Comprueba si cada uno de ellos ha sido definido antes de aparecer.
\end{ejerciciospropuestos}

\resumen{El capítulo introduce la disciplina básica de escritura matemática: ningún símbolo relevante debe aparecer sin haber sido previamente definido. Se ha distinguido entre lenguaje natural y lenguaje matemático, entre definición y enunciado, entre ejemplo y contraejemplo, y entre cálculo, justificación, argumento y demostración. Esta disciplina de precisión será la base de todos los capítulos posteriores.}

% ============================================================================
% Fin del Capítulo 1
% ============================================================================

% ============================================================================

\chapter{El lenguaje de las matemáticas}
\label{chap:lenguaje-matematicas}

\begin{objetivos}
  \item Comprender la diferencia entre lenguaje natural y lenguaje matemático.
  \item Identificar definiciones, enunciados, ejemplos, contraejemplos y demostraciones.
  \item Reconocer que cada símbolo relevante debe ser definido antes de utilizarse.
  \item Distinguir entre calcular, justificar, argumentar y demostrar.
\end{objetivos}

\section{Motivación}
\label{sec:motivacion-lenguaje}

Este capítulo tiene como finalidad introducir el modo en que se escribe, se lee y se razona en matemáticas. No se trata todavía de desarrollar una teoría específica, sino de fijar las reglas básicas del lenguaje matemático que se utilizarán durante todo el curso.

En matemáticas no basta con tener una intuición correcta. También es necesario expresar esa intuición mediante un lenguaje preciso. Una misma frase del lenguaje natural puede admitir varias interpretaciones, mientras que una afirmación matemática debe tener un significado claramente determinado dentro del contexto en el que aparece.

Por ejemplo, la frase ``un número grande'' no tiene un significado matemático preciso si no se especifica qué se entiende por número, respecto de qué comparación se considera grande y en qué conjunto se está trabajando. En cambio, la afirmación
\[
  n > 10,
\]
tiene un significado preciso si previamente se ha indicado que \(n\) es un número natural, que \(10\) es el número natural diez y que el símbolo \(>\) denota la relación usual de orden en el conjunto de los números naturales.

Esta observación conduce a una regla fundamental del curso:

\begin{convencion}[Principio de definición previa]
Todo símbolo relevante debe definirse antes de utilizarse.
\end{convencion}

Este principio no es una exigencia meramente formal. Es una condición necesaria para que un texto matemático sea comprensible, verificable y riguroso. Durante el curso se insistirá de forma sistemática en la identificación de símbolos no definidos, en la escritura correcta de definiciones y en la formulación precisa de enunciados.

\section{Lenguaje natural y lenguaje matemático}
\label{sec:lenguaje-natural-matematico}

El lenguaje natural es el lenguaje que utilizamos habitualmente para comunicarnos. Es flexible, admite matices, depende del contexto y puede contener ambigüedades. Esta flexibilidad es útil en la comunicación ordinaria, pero puede ser problemática en matemáticas.

El lenguaje matemático, en cambio, busca reducir la ambigüedad. Para ello utiliza símbolos, definiciones, reglas lógicas y convenciones de notación. Su finalidad no es sustituir por completo al lenguaje natural, sino complementarlo con una estructura precisa.

\begin{definicion}[Lenguaje matemático]
Llamaremos \emph{lenguaje matemático} al sistema formado por palabras, símbolos, reglas de escritura y reglas de interpretación que se emplean para definir objetos matemáticos, formular enunciados y construir demostraciones.
\end{definicion}

El lenguaje matemático combina dos componentes:

\begin{itemize}
  \item una parte verbal, escrita en lenguaje natural;
  \item una parte simbólica, formada por signos matemáticos como \(=\), \(<\), \(\in\), \(\subseteq\), \(+\), \(\forall\), \(\exists\), entre otros.
\end{itemize}

Ninguna de estas partes debe aparecer de manera aislada. Un texto matemático correcto debe explicar qué significan los símbolos que utiliza y debe emplear las palabras con un sentido compatible con las definiciones dadas.

\begin{ejemplo}[Ambigüedad en lenguaje natural]
La frase ``todos los números tienen un siguiente'' puede interpretarse de varias formas. Para que sea una afirmación matemática precisa, debemos indicar el conjunto de números al que nos referimos.

Si trabajamos en el conjunto \(\mathbb{N}\) de los números naturales, y si por ``siguiente'' entendemos el número obtenido al sumar \(1\), entonces la afirmación puede escribirse como:
\[
  \text{para todo } n\in \mathbb{N}, \text{ existe } m\in \mathbb{N} \text{ tal que } m=n+1.
\]
En esta formulación aparecen símbolos como \(\mathbb{N}\), \(\in\), \(m\), \(n\), \(+\) y \(=\), cuyo significado debe haber sido definido o aceptado previamente como parte de la notación del curso.
\end{ejemplo}

\begin{advertencia}
Una expresión simbólica no es necesariamente más rigurosa por el simple hecho de contener símbolos. Una fórmula con símbolos no definidos puede ser menos clara que una frase escrita cuidadosamente en lenguaje natural.
\end{advertencia}

% --------------------------------------------------------------
%  Bloque nuevo: Perspectiva histórica y formal del lenguaje
% --------------------------------------------------------------

\section{Perspectiva histórica y formal del lenguaje}
\label{sec:perspectiva-bourbaki}

Bourbaki, en su obra \emph{Éléments de mathématique}~\cite{bourbaki1970},
remarca que la matemática, desde los griegos, se caracteriza por la
\textbf{demostración rigurosa}.  Una de sus observaciones más importantes es:

\begin{quote}
  « … el armazón lógico de cada capítulo está constituido por definiciones,
  axiomas y teoremas.  El lenguaje debe ser rigurosamente correcto;
  los abusos de lenguaje o de notación, sin los cuales cualquier texto
  matemático se vuelve pedante e ilegible, han sido señalados. »
  \hfill \textit{(Éléments de mathématique, p.~6)}
\end{quote}

A partir de esta reflexión, Bourbaki extrae dos principios que encajan
perfectamente con los que hemos adoptado en el curso:

\begin{convencion}[Principio de definición previa (Bourbaki)]
  Todo símbolo relevante debe definirse antes de utilizarse.
  Este principio asegura que el lector pueda reconstruir el sentido de
  cada paso sin ambigüedades.
\end{convencion}

\begin{convencion}[Método axiomático]
  La escritura de un texto matemático debe seguir el esquema
  ``definiciones → axiomas → teoremas''; de esta forma el texto es
  \emph{fácil de formalizar}.
\end{convencion}

\begin{cita}
  « La verificación de un texto formalizado sólo requiere una atención en
  cierto modo mecánica; las únicas fuentes de error provienen de la
  longitud o la complejidad del texto. » 
  \hfill \textit{(Éléments de mathématique, p.~8)}
\end{cita}

Estos principios refuerzan la \hyperref[sec:lenguaje-natural-matematico]{sección~\ref*{sec:lenguaje-natural-matematico}}:
el lenguaje natural sirve de base, pero la exactitud sólo se consigue
cuando cada símbolo y cada regla están claramente especificados.


\section{Objetos, símbolos y significado}
\label{sec:objetos-simbolos-significado}

La matemática estudia objetos. Estos objetos pueden ser números, puntos, conjuntos, funciones, relaciones, matrices, espacios, grafos u otras estructuras. Para referirnos a ellos utilizamos símbolos.

\begin{definicion}[Objeto matemático]
Un \emph{objeto matemático} es una entidad definida dentro de un contexto matemático. Puede ser un número, un conjunto, una función, una relación, una operación o una estructura formada por varios objetos relacionados entre sí.
\end{definicion}

\begin{definicion}[Símbolo]
Un \emph{símbolo} es un signo que se utiliza para representar un objeto, una relación, una operación o una propiedad matemática.
\end{definicion}

Por ejemplo, en la expresión
\[
  x \in A,
\]
el símbolo \(x\) suele representar un objeto, el símbolo \(A\) suele representar un conjunto y el símbolo \(\in\) representa la relación de pertenencia. La expresión se lee: ``\(x\) pertenece a \(A\)''.

Sin embargo, esta lectura solo es legítima si antes se ha explicado qué representa \(x\), qué representa \(A\) y qué significa \(\in\). Por esta razón adoptamos la siguiente convención.

\begin{convencion}[Definición previa de los símbolos]
Siempre que se introduzca un símbolo nuevo, se indicará explícitamente qué objeto representa. Por ejemplo, antes de escribir una expresión como \(x\in A\), se deberá haber explicado qué es \(x\), qué es \(A\) y qué significa la relación \(\in\).
\end{convencion}

\begin{ejemplo}[Introducción correcta de símbolos]
Una forma correcta de introducir la expresión \(x\in A\) es la siguiente:

Sea \(A\) un conjunto y sea \(x\) un objeto. Escribiremos \(x\in A\) para indicar que \(x\) pertenece al conjunto \(A\).
\end{ejemplo}

\begin{ejemplo}[Uso incorrecto de símbolos]
La frase siguiente no es adecuada al comienzo de un texto:
\[
  x\in A \Rightarrow x\in B.
\]
No se ha indicado qué son \(A\), \(B\), \(x\), ni qué significa el símbolo \(\Rightarrow\). Aunque un lector experto pueda adivinar la intención, el texto no cumple el principio de definición previa.
\end{ejemplo}

\subsection{Variables y constantes}
\label{subsec:variables-constantes}

En matemáticas es importante distinguir entre símbolos que representan objetos fijos y símbolos que pueden representar objetos variables.

\begin{definicion}[Variable]
Una \emph{variable} es un símbolo que puede representar distintos objetos dentro de un conjunto o contexto previamente especificado.
\end{definicion}

\begin{definicion}[Constante]
Una \emph{constante} es un símbolo que representa un objeto fijo dentro del contexto considerado.
\end{definicion}

Por ejemplo, si escribimos ``sea \(n\in \mathbb{N}\)'', el símbolo \(n\) se utiliza como variable que representa un número natural arbitrario. En cambio, el símbolo \(0\), una vez definido, representa un número natural concreto.

\begin{advertencia}
El mismo símbolo puede tener significados distintos en textos distintos. Por ejemplo, \(n\) puede representar un número natural en un problema, una dimensión en otro y un índice en otro. Por eso es necesario declarar su significado cada vez que se inicia un nuevo contexto.
\end{advertencia}

\section{Definiciones matemáticas}
\label{sec:definiciones-matematicas}

Una definición introduce un objeto, una propiedad o una relación. En matemáticas, definir no significa describir de manera aproximada, sino fijar con precisión el significado de un término o símbolo.

\begin{definicion}[Definición matemática]
Una \emph{definición matemática} es una declaración que asigna un significado preciso a un término, símbolo, objeto, propiedad o relación dentro de un contexto determinado.
\end{definicion}

Las definiciones son acuerdos de significado. No son verdaderas ni falsas en el mismo sentido que un teorema. Una vez aceptada una definición, podemos utilizarla para formular enunciados y demostrar resultados.

\begin{ejemplo}[Definición de número par]
Sea \(n\) un número entero. Diremos que \(n\) es \emph{par} si existe un número entero \(k\) tal que
\[
  n = 2k.
\]
En esta definición se debe haber indicado previamente qué es un número entero y qué significan los símbolos \(=\), \(2\) y el producto \(2k\).
\end{ejemplo}

\begin{ejemplo}[Definición de subconjunto]
Sean \(A\) y \(B\) dos conjuntos. Diremos que \(A\) es un \emph{subconjunto} de \(B\), y escribiremos
\[
  A\subseteq B,
\]
si todo elemento de \(A\) es también un elemento de \(B\).
\end{ejemplo}

Una buena definición debe cumplir varias condiciones:

\begin{itemize}
  \item debe indicar el contexto en el que se trabaja;
  \item debe introducir claramente los símbolos nuevos;
  \item debe evitar circularidades;
  \item debe permitir decidir, al menos en los casos básicos, si un objeto satisface o no la propiedad definida;
  \item debe ser compatible con el uso posterior que se hará del concepto.
\end{itemize}

\begin{advertencia}[Definiciones circulares]
Una definición es circular cuando utiliza el término que pretende definir. Por ejemplo, decir que ``un conjunto finito es un conjunto que tiene finitos elementos'' no es una definición satisfactoria si no se ha definido antes qué significa ``finito''.
\end{advertencia}

\section{Enunciados matemáticos}
\label{sec:enunciados-matematicos}

No toda frase que contiene palabras matemáticas es un enunciado matemático. Para que una frase sea un enunciado matemático debe afirmar algo susceptible de ser verdadero o falso dentro de un contexto determinado.

\begin{definicion}[Enunciado matemático]
Un \emph{enunciado matemático} es una afirmación formulada con precisión que, dentro de un contexto matemático dado, tiene un valor de verdad: verdadero o falso.
\end{definicion}

\begin{ejemplo}[Enunciados matemáticos]
Las siguientes frases son enunciados matemáticos, siempre que se hayan definido los símbolos implicados:
\begin{enumerate}
  \item El número \(2\) es par.
  \item Para todo número natural \(n\), se cumple \(n+1>n\).
  \item Existe un número entero \(k\) tal que \(k^2=4\).
\end{enumerate}
\end{ejemplo}

\begin{ejemplo}[Frases que no son enunciados matemáticos]
Las siguientes frases no son enunciados matemáticos en sentido estricto:
\begin{enumerate}
  \item ``Calcula \(2+3\)''. Es una orden, no una afirmación.
  \item ``\(x+1\)''. Es una expresión, no una afirmación.
  \item ``Los números son interesantes''. Es una valoración, no una afirmación matemática precisa.
\end{enumerate}
\end{ejemplo}

Los enunciados matemáticos pueden adoptar distintas formas. Algunas de las más frecuentes son:

\begin{itemize}
  \item afirmaciones universales: ``para todo objeto se cumple una propiedad'';
  \item afirmaciones existenciales: ``existe al menos un objeto que cumple una propiedad'';
  \item implicaciones: ``si se cumple una hipótesis, entonces se cumple una conclusión'';
  \item equivalencias: ``una propiedad se cumple si y solo si se cumple otra''.
\end{itemize}

\begin{advertencia}
En capítulos posteriores se estudiarán con detalle los conectores lógicos y los cuantificadores. En este capítulo basta con comprender que la forma lógica de un enunciado influye en el tipo de demostración que se debe realizar.
\end{advertencia}

\section{Ejemplos y contraejemplos}
\label{sec:ejemplos-contraejemplos}

Los ejemplos y contraejemplos son herramientas esenciales para comprender una definición y para analizar la validez de un enunciado.

\begin{definicion}[Ejemplo]
Un \emph{ejemplo} de una definición o propiedad es un objeto concreto que satisface dicha definición o propiedad.
\end{definicion}

\begin{definicion}[Contraejemplo]
Un \emph{contraejemplo} a un enunciado universal es un objeto concreto que cumple las hipótesis del enunciado pero no cumple su conclusión.
\end{definicion}

\begin{ejemplo}[Ejemplo de número par]
El número \(8\) es par, porque existe un número entero \(k\), concretamente \(k=4\), tal que
\[
  8 = 2k.
\]
\end{ejemplo}

\begin{ejemplo}[Contraejemplo]
Considérese la afirmación:

``Todo número entero positivo es par''.

Esta afirmación es falsa. Un contraejemplo es el número \(3\), porque \(3\) es un número entero positivo y no existe ningún número entero \(k\) tal que \(3=2k\).
\end{ejemplo}

Un solo contraejemplo basta para demostrar que un enunciado universal es falso. En cambio, muchos ejemplos no bastan para demostrar que un enunciado universal es verdadero.

\begin{advertencia}
Comprobar muchos casos particulares puede sugerir que una afirmación es verdadera, pero no constituye una demostración general. La matemática exige justificar por qué el resultado se cumple en todos los casos contemplados por el enunciado.
\end{advertencia}

\section{Demostrar y calcular}
\label{sec:demostrar-calcular}

Una parte importante del aprendizaje matemático consiste en distinguir entre calcular, justificar, argumentar y demostrar. Estas actividades están relacionadas, pero no son equivalentes.

\begin{definicion}[Cálculo]
Un \emph{cálculo} es una manipulación de expresiones mediante reglas previamente aceptadas, con el fin de obtener un resultado concreto.
\end{definicion}

\begin{definicion}[Justificación]
Una \emph{justificación} es una explicación que indica por qué un paso, una afirmación o un cálculo es válido.
\end{definicion}

\begin{definicion}[Argumento]
Un \emph{argumento} es una cadena de razones que pretende apoyar una conclusión.
\end{definicion}

\begin{definicion}[Demostración]
Una \emph{demostración} es un argumento matemático completo y riguroso que establece la verdad de un enunciado a partir de definiciones, resultados previamente demostrados y reglas lógicas aceptadas.
\end{definicion}

\begin{ejemplo}[Cálculo sin demostración]
La igualdad
\[
  2+3=5
\]
es un cálculo elemental. Pero si el objetivo fuera demostrar una propiedad general, por ejemplo
\[
  n+0=n \quad \text{para todo } n\in\mathbb{N},
\]
no bastaría con comprobar algunos casos particulares como \(1+0=1\), \(2+0=2\) y \(3+0=3\). Sería necesario dar una demostración general.
\end{ejemplo}

\begin{ejemplo}[Demostración sencilla]
Demostremos que la suma de dos números pares es par.

Sean \(a\) y \(b\) dos números enteros pares. Por definición de número par, existen dos números enteros \(r\) y \(s\) tales que
\[
  a=2r
  \quad \text{y} \quad
  b=2s.
\]
Entonces
\[
  a+b = 2r+2s = 2(r+s).
\]
Como \(r+s\) es un número entero, se deduce que \(a+b\) es par. Por tanto, la suma de dos números pares es par.
\end{ejemplo}

En esta demostración se observa una estructura típica:

\begin{enumerate}
  \item se declaran los objetos con los que se trabaja;
  \item se usan las definiciones pertinentes;
  \item se realiza un cálculo justificado;
  \item se interpreta el resultado obtenido;
  \item se formula la conclusión.
\end{enumerate}

\begin{convencion}[Estructura básica de una demostración]
Siempre que sea posible, una demostración deberá indicar explícitamente qué se supone, qué se quiere probar, qué definiciones se utilizan y cómo se obtiene la conclusión.
\end{convencion}

\section{Lectura y escritura de textos matemáticos}
\label{sec:lectura-escritura-textos}

Leer matemáticas no consiste únicamente en pasar la vista por una sucesión de símbolos. Es necesario reconstruir el significado de cada frase, comprobar qué objetos han sido definidos y entender qué papel desempeña cada hipótesis.

Al leer un texto matemático conviene hacerse las siguientes preguntas:

\begin{itemize}
  \item ¿Cuál es el contexto en el que se trabaja?
  \item ¿Qué objetos se han definido?
  \item ¿Qué símbolos aparecen y qué significa cada uno?
  \item ¿Cuál es la afirmación principal?
  \item ¿Qué hipótesis se utilizan?
  \item ¿Qué se quiere demostrar?
  \item ¿Qué definiciones o resultados previos se están aplicando?
\end{itemize}

Escribir matemáticas exige el mismo cuidado. Un texto matemático debe ser legible, preciso y ordenado. La claridad no está reñida con el rigor; al contrario, el rigor suele exigir claridad.

\subsection{Plantilla básica para escribir una demostración}
\label{subsec:plantilla-demostracion}

Una demostración elemental puede organizarse de acuerdo con la siguiente plantilla:

\begin{enumerate}
  \item \textbf{Declaración de objetos.} Indicar qué objetos se consideran y en qué conjunto o contexto viven.
  \item \textbf{Hipótesis.} Escribir con precisión qué se supone.
  \item \textbf{Objetivo.} Indicar qué se quiere probar.
  \item \textbf{Uso de definiciones.} Traducir las hipótesis mediante las definiciones correspondientes.
  \item \textbf{Desarrollo.} Encadenar igualdades, implicaciones o argumentos justificados.
  \item \textbf{Conclusión.} Explicar por qué lo obtenido prueba exactamente lo que se quería demostrar.
\end{enumerate}

\begin{ejemplo}[Reescritura rigurosa de un enunciado]
La frase ``si un número acaba en cero, entonces es divisible'' es ambigua. No se especifica qué tipo de número se considera ni por qué número debe ser divisible.

Una formulación más precisa sería:

Sea \(n\) un número entero escrito en base decimal. Si la última cifra decimal de \(n\) es \(0\), entonces \(n\) es divisible por \(10\).
\end{ejemplo}

\begin{advertencia}[Símbolos introducidos demasiado tarde]
No es recomendable escribir primero una fórmula y explicar después qué significan sus símbolos. En un texto bien organizado, los símbolos relevantes se presentan antes de ser utilizados.
\end{advertencia}

\section{Errores frecuentes}
\label{sec:errores-frecuentes-lenguaje}

En los primeros cursos universitarios aparecen algunos errores recurrentes relacionados con el lenguaje matemático. Conviene identificarlos desde el principio.

\begin{enumerate}
  \item \textbf{Usar símbolos sin definirlos.} Por ejemplo, escribir \(x\in A\) sin haber explicado qué son \(x\), \(A\) y \(\in\).

  \item \textbf{Confundir ejemplo con demostración.} Verificar un resultado para varios casos particulares no demuestra que sea verdadero en general.

  \item \textbf{Confundir definición con teorema.} Una definición introduce significado; un teorema afirma algo que debe demostrarse.

  \item \textbf{Omitir el contexto.} No es lo mismo trabajar en \(\mathbb{N}\), en \(\mathbb{Z}\), en \(\mathbb{Q}\), en \(\mathbb{R}\) o en \(\mathbb{C}\).

  \item \textbf{Utilizar flechas de manera imprecisa.} Símbolos como \(\Rightarrow\), \(\Leftrightarrow\) o \(\mapsto\) tienen significados concretos y no deben utilizarse como simples signos decorativos.

  \item \textbf{No distinguir entre objeto y propiedad.} Por ejemplo, un número puede tener la propiedad de ser par, pero ``ser par'' no es un número.

  \item \textbf{No cerrar la demostración.} Una demostración debe terminar indicando claramente que se ha probado la afirmación deseada.
\end{enumerate}

\begin{convencion}[Control de calidad de un texto matemático]
Antes de dar por correcto un texto matemático, comprobaremos que todos los símbolos relevantes han sido definidos, que las afirmaciones tienen sentido en el contexto fijado y que cada conclusión está justificada.
\end{convencion}

\begin{seminario}
  \item Identificación de símbolos no definidos en textos matemáticos breves.
  \item Reescritura rigurosa de enunciados ambiguos.
  \item Construcción de ejemplos y contraejemplos elementales.
  \item Distinción entre definición, enunciado, ejemplo, contraejemplo y demostración.
  \item Discusión de demostraciones incompletas y detección de pasos no justificados.
\end{seminario}

\begin{ejerciciospropuestos}
  \item Localiza en un texto matemático breve todos los símbolos que aparecen y escribe una lista con su significado.

  \item Escribe tres ejemplos de afirmaciones matemáticas y tres ejemplos de frases que no sean enunciados matemáticos.

  \item Reescribe de forma rigurosa las siguientes frases:
  \begin{enumerate}
    \item ``Todo número tiene un siguiente.''
    \item ``Si un número es grande, su cuadrado también lo es.''
    \item ``Los números negativos al multiplicarse dan positivo.''
  \end{enumerate}

  \item Indica cuáles de las siguientes frases son definiciones, cuáles son enunciados y cuáles no son afirmaciones matemáticas:
  \begin{enumerate}
    \item ``Un número entero \(n\) es par si existe un entero \(k\) tal que \(n=2k\).''
    \item ``Todo número primo mayor que \(2\) es impar.''
    \item ``Calcula \(7+5\).''
    \item ``Sea \(A\) un conjunto.''
  \end{enumerate}

  \item Da un contraejemplo para cada una de las siguientes afirmaciones falsas:
  \begin{enumerate}
    \item Todo número entero es positivo.
    \item Todo número impar es primo.
    \item Si \(a^2=b^2\), entonces \(a=b\), donde \(a\) y \(b\) son números enteros.
  \end{enumerate}

  \item Explica por qué comprobar que \(1+3\), \(3+5\) y \(5+7\) son números pares no demuestra por sí solo que la suma de dos números impares cualesquiera sea par.

  \item Escribe una demostración completa de la siguiente afirmación: la suma de dos números impares es par. Antes de empezar, define qué significa que un número entero sea impar.

  \item Revisa una solución propia de un ejercicio matemático e identifica todos los símbolos que has utilizado. Comprueba si cada uno de ellos ha sido definido antes de aparecer.
\end{ejerciciospropuestos}

\resumen{El capítulo introduce la disciplina básica de escritura matemática: ningún símbolo relevante debe aparecer sin haber sido previamente definido. Se ha distinguido entre lenguaje natural y lenguaje matemático, entre definición y enunciado, entre ejemplo y contraejemplo, y entre cálculo, justificación, argumento y demostración. Esta disciplina de precisión será la base de todos los capítulos posteriores.}

% ============================================================================
% Fin del Capítulo 1
% ============================================================================

% ============================================================================

\chapter{El lenguaje de las matemáticas}
\label{chap:lenguaje-matematicas}

\begin{objetivos}
  \item Comprender la diferencia entre lenguaje natural y lenguaje matemático.
  \item Identificar definiciones, enunciados, ejemplos, contraejemplos y demostraciones.
  \item Reconocer que cada símbolo relevante debe ser definido antes de utilizarse.
  \item Distinguir entre calcular, justificar, argumentar y demostrar.
\end{objetivos}

\section{Motivación}
\label{sec:motivacion-lenguaje}

Este capítulo tiene como finalidad introducir el modo en que se escribe, se lee y se razona en matemáticas. No se trata todavía de desarrollar una teoría específica, sino de fijar las reglas básicas del lenguaje matemático que se utilizarán durante todo el curso.

En matemáticas no basta con tener una intuición correcta. También es necesario expresar esa intuición mediante un lenguaje preciso. Una misma frase del lenguaje natural puede admitir varias interpretaciones, mientras que una afirmación matemática debe tener un significado claramente determinado dentro del contexto en el que aparece.

Por ejemplo, la frase ``un número grande'' no tiene un significado matemático preciso si no se especifica qué se entiende por número, respecto de qué comparación se considera grande y en qué conjunto se está trabajando. En cambio, la afirmación
\[
  n > 10,
\]
tiene un significado preciso si previamente se ha indicado que \(n\) es un número natural, que \(10\) es el número natural diez y que el símbolo \(>\) denota la relación usual de orden en el conjunto de los números naturales.

Esta observación conduce a una regla fundamental del curso:

\begin{convencion}[Principio de definición previa]
Todo símbolo relevante debe definirse antes de utilizarse.
\end{convencion}

Este principio no es una exigencia meramente formal. Es una condición necesaria para que un texto matemático sea comprensible, verificable y riguroso. Durante el curso se insistirá de forma sistemática en la identificación de símbolos no definidos, en la escritura correcta de definiciones y en la formulación precisa de enunciados.

\section{Lenguaje natural y lenguaje matemático}
\label{sec:lenguaje-natural-matematico}

El lenguaje natural es el lenguaje que utilizamos habitualmente para comunicarnos. Es flexible, admite matices, depende del contexto y puede contener ambigüedades. Esta flexibilidad es útil en la comunicación ordinaria, pero puede ser problemática en matemáticas.

El lenguaje matemático, en cambio, busca reducir la ambigüedad. Para ello utiliza símbolos, definiciones, reglas lógicas y convenciones de notación. Su finalidad no es sustituir por completo al lenguaje natural, sino complementarlo con una estructura precisa.

\begin{definicion}[Lenguaje matemático]
Llamaremos \emph{lenguaje matemático} al sistema formado por palabras, símbolos, reglas de escritura y reglas de interpretación que se emplean para definir objetos matemáticos, formular enunciados y construir demostraciones.
\end{definicion}

El lenguaje matemático combina dos componentes:

\begin{itemize}
  \item una parte verbal, escrita en lenguaje natural;
  \item una parte simbólica, formada por signos matemáticos como \(=\), \(<\), \(\in\), \(\subseteq\), \(+\), \(\forall\), \(\exists\), entre otros.
\end{itemize}

Ninguna de estas partes debe aparecer de manera aislada. Un texto matemático correcto debe explicar qué significan los símbolos que utiliza y debe emplear las palabras con un sentido compatible con las definiciones dadas.

\begin{ejemplo}[Ambigüedad en lenguaje natural]
La frase ``todos los números tienen un siguiente'' puede interpretarse de varias formas. Para que sea una afirmación matemática precisa, debemos indicar el conjunto de números al que nos referimos.

Si trabajamos en el conjunto \(\mathbb{N}\) de los números naturales, y si por ``siguiente'' entendemos el número obtenido al sumar \(1\), entonces la afirmación puede escribirse como:
\[
  \text{para todo } n\in \mathbb{N}, \text{ existe } m\in \mathbb{N} \text{ tal que } m=n+1.
\]
En esta formulación aparecen símbolos como \(\mathbb{N}\), \(\in\), \(m\), \(n\), \(+\) y \(=\), cuyo significado debe haber sido definido o aceptado previamente como parte de la notación del curso.
\end{ejemplo}

\begin{advertencia}
Una expresión simbólica no es necesariamente más rigurosa por el simple hecho de contener símbolos. Una fórmula con símbolos no definidos puede ser menos clara que una frase escrita cuidadosamente en lenguaje natural.
\end{advertencia}

% --------------------------------------------------------------
%  Bloque nuevo: Perspectiva histórica y formal del lenguaje
% --------------------------------------------------------------

\section{Perspectiva histórica y formal del lenguaje}
\label{sec:perspectiva-bourbaki}

Bourbaki, en su obra \emph{Éléments de mathématique}~\cite{bourbaki1970},
remarca que la matemática, desde los griegos, se caracteriza por la
\textbf{demostración rigurosa}.  Una de sus observaciones más importantes es:

\begin{quote}
  « … el armazón lógico de cada capítulo está constituido por definiciones,
  axiomas y teoremas.  El lenguaje debe ser rigurosamente correcto;
  los abusos de lenguaje o de notación, sin los cuales cualquier texto
  matemático se vuelve pedante e ilegible, han sido señalados. »
  \hfill \textit{(Éléments de mathématique, p.~6)}
\end{quote}

A partir de esta reflexión, Bourbaki extrae dos principios que encajan
perfectamente con los que hemos adoptado en el curso:

\begin{convencion}[Principio de definición previa (Bourbaki)]
  Todo símbolo relevante debe definirse antes de utilizarse.
  Este principio asegura que el lector pueda reconstruir el sentido de
  cada paso sin ambigüedades.
\end{convencion}

\begin{convencion}[Método axiomático]
  La escritura de un texto matemático debe seguir el esquema
  ``definiciones → axiomas → teoremas''; de esta forma el texto es
  \emph{fácil de formalizar}.
\end{convencion}

\begin{cita}
  « La verificación de un texto formalizado sólo requiere una atención en
  cierto modo mecánica; las únicas fuentes de error provienen de la
  longitud o la complejidad del texto. » 
  \hfill \textit{(Éléments de mathématique, p.~8)}
\end{cita}

Estos principios refuerzan la \hyperref[sec:lenguaje-natural-matematico]{sección~\ref*{sec:lenguaje-natural-matematico}}:
el lenguaje natural sirve de base, pero la exactitud sólo se consigue
cuando cada símbolo y cada regla están claramente especificados.


\section{Objetos, símbolos y significado}
\label{sec:objetos-simbolos-significado}

La matemática estudia objetos. Estos objetos pueden ser números, puntos, conjuntos, funciones, relaciones, matrices, espacios, grafos u otras estructuras. Para referirnos a ellos utilizamos símbolos.

\begin{definicion}[Objeto matemático]
Un \emph{objeto matemático} es una entidad definida dentro de un contexto matemático. Puede ser un número, un conjunto, una función, una relación, una operación o una estructura formada por varios objetos relacionados entre sí.
\end{definicion}

\begin{definicion}[Símbolo]
Un \emph{símbolo} es un signo que se utiliza para representar un objeto, una relación, una operación o una propiedad matemática.
\end{definicion}

Por ejemplo, en la expresión
\[
  x \in A,
\]
el símbolo \(x\) suele representar un objeto, el símbolo \(A\) suele representar un conjunto y el símbolo \(\in\) representa la relación de pertenencia. La expresión se lee: ``\(x\) pertenece a \(A\)''.

Sin embargo, esta lectura solo es legítima si antes se ha explicado qué representa \(x\), qué representa \(A\) y qué significa \(\in\). Por esta razón adoptamos la siguiente convención.

\begin{convencion}[Definición previa de los símbolos]
Siempre que se introduzca un símbolo nuevo, se indicará explícitamente qué objeto representa. Por ejemplo, antes de escribir una expresión como \(x\in A\), se deberá haber explicado qué es \(x\), qué es \(A\) y qué significa la relación \(\in\).
\end{convencion}

\begin{ejemplo}[Introducción correcta de símbolos]
Una forma correcta de introducir la expresión \(x\in A\) es la siguiente:

Sea \(A\) un conjunto y sea \(x\) un objeto. Escribiremos \(x\in A\) para indicar que \(x\) pertenece al conjunto \(A\).
\end{ejemplo}

\begin{ejemplo}[Uso incorrecto de símbolos]
La frase siguiente no es adecuada al comienzo de un texto:
\[
  x\in A \Rightarrow x\in B.
\]
No se ha indicado qué son \(A\), \(B\), \(x\), ni qué significa el símbolo \(\Rightarrow\). Aunque un lector experto pueda adivinar la intención, el texto no cumple el principio de definición previa.
\end{ejemplo}

\subsection{Variables y constantes}
\label{subsec:variables-constantes}

En matemáticas es importante distinguir entre símbolos que representan objetos fijos y símbolos que pueden representar objetos variables.

\begin{definicion}[Variable]
Una \emph{variable} es un símbolo que puede representar distintos objetos dentro de un conjunto o contexto previamente especificado.
\end{definicion}

\begin{definicion}[Constante]
Una \emph{constante} es un símbolo que representa un objeto fijo dentro del contexto considerado.
\end{definicion}

Por ejemplo, si escribimos ``sea \(n\in \mathbb{N}\)'', el símbolo \(n\) se utiliza como variable que representa un número natural arbitrario. En cambio, el símbolo \(0\), una vez definido, representa un número natural concreto.

\begin{advertencia}
El mismo símbolo puede tener significados distintos en textos distintos. Por ejemplo, \(n\) puede representar un número natural en un problema, una dimensión en otro y un índice en otro. Por eso es necesario declarar su significado cada vez que se inicia un nuevo contexto.
\end{advertencia}

\section{Definiciones matemáticas}
\label{sec:definiciones-matematicas}

Una definición introduce un objeto, una propiedad o una relación. En matemáticas, definir no significa describir de manera aproximada, sino fijar con precisión el significado de un término o símbolo.

\begin{definicion}[Definición matemática]
Una \emph{definición matemática} es una declaración que asigna un significado preciso a un término, símbolo, objeto, propiedad o relación dentro de un contexto determinado.
\end{definicion}

Las definiciones son acuerdos de significado. No son verdaderas ni falsas en el mismo sentido que un teorema. Una vez aceptada una definición, podemos utilizarla para formular enunciados y demostrar resultados.

\begin{ejemplo}[Definición de número par]
Sea \(n\) un número entero. Diremos que \(n\) es \emph{par} si existe un número entero \(k\) tal que
\[
  n = 2k.
\]
En esta definición se debe haber indicado previamente qué es un número entero y qué significan los símbolos \(=\), \(2\) y el producto \(2k\).
\end{ejemplo}

\begin{ejemplo}[Definición de subconjunto]
Sean \(A\) y \(B\) dos conjuntos. Diremos que \(A\) es un \emph{subconjunto} de \(B\), y escribiremos
\[
  A\subseteq B,
\]
si todo elemento de \(A\) es también un elemento de \(B\).
\end{ejemplo}

Una buena definición debe cumplir varias condiciones:

\begin{itemize}
  \item debe indicar el contexto en el que se trabaja;
  \item debe introducir claramente los símbolos nuevos;
  \item debe evitar circularidades;
  \item debe permitir decidir, al menos en los casos básicos, si un objeto satisface o no la propiedad definida;
  \item debe ser compatible con el uso posterior que se hará del concepto.
\end{itemize}

\begin{advertencia}[Definiciones circulares]
Una definición es circular cuando utiliza el término que pretende definir. Por ejemplo, decir que ``un conjunto finito es un conjunto que tiene finitos elementos'' no es una definición satisfactoria si no se ha definido antes qué significa ``finito''.
\end{advertencia}

\section{Enunciados matemáticos}
\label{sec:enunciados-matematicos}

No toda frase que contiene palabras matemáticas es un enunciado matemático. Para que una frase sea un enunciado matemático debe afirmar algo susceptible de ser verdadero o falso dentro de un contexto determinado.

\begin{definicion}[Enunciado matemático]
Un \emph{enunciado matemático} es una afirmación formulada con precisión que, dentro de un contexto matemático dado, tiene un valor de verdad: verdadero o falso.
\end{definicion}

\begin{ejemplo}[Enunciados matemáticos]
Las siguientes frases son enunciados matemáticos, siempre que se hayan definido los símbolos implicados:
\begin{enumerate}
  \item El número \(2\) es par.
  \item Para todo número natural \(n\), se cumple \(n+1>n\).
  \item Existe un número entero \(k\) tal que \(k^2=4\).
\end{enumerate}
\end{ejemplo}

\begin{ejemplo}[Frases que no son enunciados matemáticos]
Las siguientes frases no son enunciados matemáticos en sentido estricto:
\begin{enumerate}
  \item ``Calcula \(2+3\)''. Es una orden, no una afirmación.
  \item ``\(x+1\)''. Es una expresión, no una afirmación.
  \item ``Los números son interesantes''. Es una valoración, no una afirmación matemática precisa.
\end{enumerate}
\end{ejemplo}

Los enunciados matemáticos pueden adoptar distintas formas. Algunas de las más frecuentes son:

\begin{itemize}
  \item afirmaciones universales: ``para todo objeto se cumple una propiedad'';
  \item afirmaciones existenciales: ``existe al menos un objeto que cumple una propiedad'';
  \item implicaciones: ``si se cumple una hipótesis, entonces se cumple una conclusión'';
  \item equivalencias: ``una propiedad se cumple si y solo si se cumple otra''.
\end{itemize}

\begin{advertencia}
En capítulos posteriores se estudiarán con detalle los conectores lógicos y los cuantificadores. En este capítulo basta con comprender que la forma lógica de un enunciado influye en el tipo de demostración que se debe realizar.
\end{advertencia}

\section{Ejemplos y contraejemplos}
\label{sec:ejemplos-contraejemplos}

Los ejemplos y contraejemplos son herramientas esenciales para comprender una definición y para analizar la validez de un enunciado.

\begin{definicion}[Ejemplo]
Un \emph{ejemplo} de una definición o propiedad es un objeto concreto que satisface dicha definición o propiedad.
\end{definicion}

\begin{definicion}[Contraejemplo]
Un \emph{contraejemplo} a un enunciado universal es un objeto concreto que cumple las hipótesis del enunciado pero no cumple su conclusión.
\end{definicion}

\begin{ejemplo}[Ejemplo de número par]
El número \(8\) es par, porque existe un número entero \(k\), concretamente \(k=4\), tal que
\[
  8 = 2k.
\]
\end{ejemplo}

\begin{ejemplo}[Contraejemplo]
Considérese la afirmación:

``Todo número entero positivo es par''.

Esta afirmación es falsa. Un contraejemplo es el número \(3\), porque \(3\) es un número entero positivo y no existe ningún número entero \(k\) tal que \(3=2k\).
\end{ejemplo}

Un solo contraejemplo basta para demostrar que un enunciado universal es falso. En cambio, muchos ejemplos no bastan para demostrar que un enunciado universal es verdadero.

\begin{advertencia}
Comprobar muchos casos particulares puede sugerir que una afirmación es verdadera, pero no constituye una demostración general. La matemática exige justificar por qué el resultado se cumple en todos los casos contemplados por el enunciado.
\end{advertencia}

\section{Demostrar y calcular}
\label{sec:demostrar-calcular}

Una parte importante del aprendizaje matemático consiste en distinguir entre calcular, justificar, argumentar y demostrar. Estas actividades están relacionadas, pero no son equivalentes.

\begin{definicion}[Cálculo]
Un \emph{cálculo} es una manipulación de expresiones mediante reglas previamente aceptadas, con el fin de obtener un resultado concreto.
\end{definicion}

\begin{definicion}[Justificación]
Una \emph{justificación} es una explicación que indica por qué un paso, una afirmación o un cálculo es válido.
\end{definicion}

\begin{definicion}[Argumento]
Un \emph{argumento} es una cadena de razones que pretende apoyar una conclusión.
\end{definicion}

\begin{definicion}[Demostración]
Una \emph{demostración} es un argumento matemático completo y riguroso que establece la verdad de un enunciado a partir de definiciones, resultados previamente demostrados y reglas lógicas aceptadas.
\end{definicion}

\begin{ejemplo}[Cálculo sin demostración]
La igualdad
\[
  2+3=5
\]
es un cálculo elemental. Pero si el objetivo fuera demostrar una propiedad general, por ejemplo
\[
  n+0=n \quad \text{para todo } n\in\mathbb{N},
\]
no bastaría con comprobar algunos casos particulares como \(1+0=1\), \(2+0=2\) y \(3+0=3\). Sería necesario dar una demostración general.
\end{ejemplo}

\begin{ejemplo}[Demostración sencilla]
Demostremos que la suma de dos números pares es par.

Sean \(a\) y \(b\) dos números enteros pares. Por definición de número par, existen dos números enteros \(r\) y \(s\) tales que
\[
  a=2r
  \quad \text{y} \quad
  b=2s.
\]
Entonces
\[
  a+b = 2r+2s = 2(r+s).
\]
Como \(r+s\) es un número entero, se deduce que \(a+b\) es par. Por tanto, la suma de dos números pares es par.
\end{ejemplo}

En esta demostración se observa una estructura típica:

\begin{enumerate}
  \item se declaran los objetos con los que se trabaja;
  \item se usan las definiciones pertinentes;
  \item se realiza un cálculo justificado;
  \item se interpreta el resultado obtenido;
  \item se formula la conclusión.
\end{enumerate}

\begin{convencion}[Estructura básica de una demostración]
Siempre que sea posible, una demostración deberá indicar explícitamente qué se supone, qué se quiere probar, qué definiciones se utilizan y cómo se obtiene la conclusión.
\end{convencion}

\section{Lectura y escritura de textos matemáticos}
\label{sec:lectura-escritura-textos}

Leer matemáticas no consiste únicamente en pasar la vista por una sucesión de símbolos. Es necesario reconstruir el significado de cada frase, comprobar qué objetos han sido definidos y entender qué papel desempeña cada hipótesis.

Al leer un texto matemático conviene hacerse las siguientes preguntas:

\begin{itemize}
  \item ¿Cuál es el contexto en el que se trabaja?
  \item ¿Qué objetos se han definido?
  \item ¿Qué símbolos aparecen y qué significa cada uno?
  \item ¿Cuál es la afirmación principal?
  \item ¿Qué hipótesis se utilizan?
  \item ¿Qué se quiere demostrar?
  \item ¿Qué definiciones o resultados previos se están aplicando?
\end{itemize}

Escribir matemáticas exige el mismo cuidado. Un texto matemático debe ser legible, preciso y ordenado. La claridad no está reñida con el rigor; al contrario, el rigor suele exigir claridad.

\subsection{Plantilla básica para escribir una demostración}
\label{subsec:plantilla-demostracion}

Una demostración elemental puede organizarse de acuerdo con la siguiente plantilla:

\begin{enumerate}
  \item \textbf{Declaración de objetos.} Indicar qué objetos se consideran y en qué conjunto o contexto viven.
  \item \textbf{Hipótesis.} Escribir con precisión qué se supone.
  \item \textbf{Objetivo.} Indicar qué se quiere probar.
  \item \textbf{Uso de definiciones.} Traducir las hipótesis mediante las definiciones correspondientes.
  \item \textbf{Desarrollo.} Encadenar igualdades, implicaciones o argumentos justificados.
  \item \textbf{Conclusión.} Explicar por qué lo obtenido prueba exactamente lo que se quería demostrar.
\end{enumerate}

\begin{ejemplo}[Reescritura rigurosa de un enunciado]
La frase ``si un número acaba en cero, entonces es divisible'' es ambigua. No se especifica qué tipo de número se considera ni por qué número debe ser divisible.

Una formulación más precisa sería:

Sea \(n\) un número entero escrito en base decimal. Si la última cifra decimal de \(n\) es \(0\), entonces \(n\) es divisible por \(10\).
\end{ejemplo}

\begin{advertencia}[Símbolos introducidos demasiado tarde]
No es recomendable escribir primero una fórmula y explicar después qué significan sus símbolos. En un texto bien organizado, los símbolos relevantes se presentan antes de ser utilizados.
\end{advertencia}

\section{Errores frecuentes}
\label{sec:errores-frecuentes-lenguaje}

En los primeros cursos universitarios aparecen algunos errores recurrentes relacionados con el lenguaje matemático. Conviene identificarlos desde el principio.

\begin{enumerate}
  \item \textbf{Usar símbolos sin definirlos.} Por ejemplo, escribir \(x\in A\) sin haber explicado qué son \(x\), \(A\) y \(\in\).

  \item \textbf{Confundir ejemplo con demostración.} Verificar un resultado para varios casos particulares no demuestra que sea verdadero en general.

  \item \textbf{Confundir definición con teorema.} Una definición introduce significado; un teorema afirma algo que debe demostrarse.

  \item \textbf{Omitir el contexto.} No es lo mismo trabajar en \(\mathbb{N}\), en \(\mathbb{Z}\), en \(\mathbb{Q}\), en \(\mathbb{R}\) o en \(\mathbb{C}\).

  \item \textbf{Utilizar flechas de manera imprecisa.} Símbolos como \(\Rightarrow\), \(\Leftrightarrow\) o \(\mapsto\) tienen significados concretos y no deben utilizarse como simples signos decorativos.

  \item \textbf{No distinguir entre objeto y propiedad.} Por ejemplo, un número puede tener la propiedad de ser par, pero ``ser par'' no es un número.

  \item \textbf{No cerrar la demostración.} Una demostración debe terminar indicando claramente que se ha probado la afirmación deseada.
\end{enumerate}

\begin{convencion}[Control de calidad de un texto matemático]
Antes de dar por correcto un texto matemático, comprobaremos que todos los símbolos relevantes han sido definidos, que las afirmaciones tienen sentido en el contexto fijado y que cada conclusión está justificada.
\end{convencion}

\begin{seminario}
  \item Identificación de símbolos no definidos en textos matemáticos breves.
  \item Reescritura rigurosa de enunciados ambiguos.
  \item Construcción de ejemplos y contraejemplos elementales.
  \item Distinción entre definición, enunciado, ejemplo, contraejemplo y demostración.
  \item Discusión de demostraciones incompletas y detección de pasos no justificados.
\end{seminario}

\begin{ejerciciospropuestos}
  \item Localiza en un texto matemático breve todos los símbolos que aparecen y escribe una lista con su significado.

  \item Escribe tres ejemplos de afirmaciones matemáticas y tres ejemplos de frases que no sean enunciados matemáticos.

  \item Reescribe de forma rigurosa las siguientes frases:
  \begin{enumerate}
    \item ``Todo número tiene un siguiente.''
    \item ``Si un número es grande, su cuadrado también lo es.''
    \item ``Los números negativos al multiplicarse dan positivo.''
  \end{enumerate}

  \item Indica cuáles de las siguientes frases son definiciones, cuáles son enunciados y cuáles no son afirmaciones matemáticas:
  \begin{enumerate}
    \item ``Un número entero \(n\) es par si existe un entero \(k\) tal que \(n=2k\).''
    \item ``Todo número primo mayor que \(2\) es impar.''
    \item ``Calcula \(7+5\).''
    \item ``Sea \(A\) un conjunto.''
  \end{enumerate}

  \item Da un contraejemplo para cada una de las siguientes afirmaciones falsas:
  \begin{enumerate}
    \item Todo número entero es positivo.
    \item Todo número impar es primo.
    \item Si \(a^2=b^2\), entonces \(a=b\), donde \(a\) y \(b\) son números enteros.
  \end{enumerate}

  \item Explica por qué comprobar que \(1+3\), \(3+5\) y \(5+7\) son números pares no demuestra por sí solo que la suma de dos números impares cualesquiera sea par.

  \item Escribe una demostración completa de la siguiente afirmación: la suma de dos números impares es par. Antes de empezar, define qué significa que un número entero sea impar.

  \item Revisa una solución propia de un ejercicio matemático e identifica todos los símbolos que has utilizado. Comprueba si cada uno de ellos ha sido definido antes de aparecer.
\end{ejerciciospropuestos}

\resumen{El capítulo introduce la disciplina básica de escritura matemática: ningún símbolo relevante debe aparecer sin haber sido previamente definido. Se ha distinguido entre lenguaje natural y lenguaje matemático, entre definición y enunciado, entre ejemplo y contraejemplo, y entre cálculo, justificación, argumento y demostración. Esta disciplina de precisión será la base de todos los capítulos posteriores.}

% ============================================================================
% Fin del Capítulo 1
% ============================================================================

% ------------------------------------------------------------
\chapter{Lógica y razonamiento matemático}

\begin{objetivos}
  \item Introducir el concepto de proposición matemática.
  \item Comprender los conectores lógicos fundamentales.
  \item Manejar cuantificadores universales y existenciales.
  \item Aprender los métodos básicos de demostración.
\end{objetivos}

\section{Motivación}

La matemática no se limita a manipular símbolos o a realizar cálculos. Una parte esencial de la actividad matemática consiste en formular afirmaciones precisas y decidir, mediante razonamientos rigurosos, si dichas afirmaciones son verdaderas o falsas. Para ello es necesario disponer de un lenguaje lógico básico.

En este capítulo se introducen las nociones elementales de lógica que se utilizarán durante todo el curso. No se pretende desarrollar una teoría formal completa de la lógica, sino proporcionar las herramientas necesarias para leer, escribir y demostrar enunciados matemáticos con precisión.

La regla metodológica que seguiremos es la misma que en el capítulo anterior: todo símbolo relevante debe definirse antes de utilizarse. Por esta razón, cada conector lógico y cada cuantificador se introducirá indicando primero su significado.

\section{Proposiciones}
\label{sec:proposiciones}

\begin{definicion}[Proposición]
Una \emph{proposición} es una oración declarativa a la que se puede asignar un valor de verdad: verdadero o falso.
\end{definicion}

En matemáticas, las proposiciones son las unidades básicas sobre las que se construyen los razonamientos. Una proposición no tiene por qué ser verdadera; basta con que tenga sentido afirmar que es verdadera o falsa.

\begin{ejemplo}
Las siguientes oraciones son proposiciones:
\begin{enumerate}
  \item El número $2$ es par.
  \item El número $7$ es múltiplo de $3$.
  \item Todo número natural mayor que $1$ es primo.
  \item Existe un número entero $x$ tal que $x^2=4$.
\end{enumerate}
La primera y la cuarta son verdaderas. La segunda y la tercera son falsas.
\end{ejemplo}

\begin{contraejemplo}
Las siguientes oraciones no son proposiciones matemáticas:
\begin{enumerate}
  \item Calcula $2+3$.
  \item ¿Es $5$ un número primo?
  \item Sea $x$ un número natural.
  \item El número $x$ es mayor que $3$.
\end{enumerate}
Las dos primeras no son oraciones declarativas. La tercera introduce una variable, pero no afirma nada. La cuarta depende del valor de $x$; antes de decidir si es verdadera o falsa hay que saber qué objeto representa $x$.
\end{contraejemplo}

\begin{convencion}[Variables y proposiciones]
Si una afirmación contiene una variable, no la consideraremos una proposición cerrada hasta que se especifique el conjunto al que pertenece la variable o se introduzca un cuantificador.
\end{convencion}

Por ejemplo, la frase
\[
  x^2 \geq 0
\]
no es todavía una proposición cerrada, porque no se ha indicado qué es $x$. En cambio, sí son proposiciones las siguientes afirmaciones:
\[
  \text{Para todo } x\in\mathbb{R},\ x^2\geq 0,
\]
y
\[
  \text{Existe } x\in\mathbb{R} \text{ tal que } x^2<0.
\]
La primera es verdadera y la segunda es falsa.

\section{Negación, conjunción y disyunción}
\label{sec:negacion-conjuncion}

En esta sección introducimos tres conectores lógicos básicos: negación, conjunción y disyunción. Para evitar ambigüedades, primero fijamos el significado de cada símbolo.

\begin{definicion}[Negación]
Si $P$ es una proposición, la \emph{negación} de $P$ se denota por $\neg P$ y se lee ``no $P$''. La proposición $\neg P$ es verdadera exactamente cuando $P$ es falsa.
\end{definicion}

\begin{ejemplo}
Sea $P$ la proposición ``$6$ es primo''. Entonces $\neg P$ es la proposición ``$6$ no es primo''. En este caso, $P$ es falsa y $\neg P$ es verdadera.
\end{ejemplo}

\begin{definicion}[Conjunción]
Si $P$ y $Q$ son proposiciones, la \emph{conjunción} de $P$ y $Q$ se denota por
\[
  P\wedge Q
\]
y se lee ``$P$ y $Q$''. La proposición $P\wedge Q$ es verdadera exactamente cuando $P$ y $Q$ son ambas verdaderas.
\end{definicion}

\begin{ejemplo}
La proposición
\[
  2 \text{ es primo } \wedge 2 \text{ es par}
\]
es verdadera, porque las dos proposiciones que la componen son verdaderas.
\end{ejemplo}

\begin{definicion}[Disyunción]
Si $P$ y $Q$ son proposiciones, la \emph{disyunción} de $P$ y $Q$ se denota por
\[
  P\vee Q
\]
y se lee ``$P$ o $Q$''. En matemáticas, salvo que se indique lo contrario, esta ``o'' es inclusiva: $P\vee Q$ es verdadera cuando al menos una de las proposiciones $P$ o $Q$ es verdadera.
\end{definicion}

\begin{ejemplo}
La proposición
\[
  4 \text{ es par } \vee 4 \text{ es primo}
\]
es verdadera, porque $4$ es par, aunque no sea primo.
\end{ejemplo}

\begin{advertencia}
En el lenguaje natural la palabra ``o'' puede ser ambigua. En matemáticas, la disyunción $P\vee Q$ se interpreta normalmente de forma inclusiva: permite que $P$ y $Q$ sean simultáneamente verdaderas.
\end{advertencia}

% --------------------------------------------------------------
%  Bloque nuevo: Razonamiento formal y método axiomático
% --------------------------------------------------------------

\section{Razonamiento formal y método axiomático}
\label{subsec:razonamiento-bourbaki}

En la introducción de los \emph{Éléments de mathématique}, Bourbaki describe
el método axiomático como el arte de redactar textos cuya
\textbf{formalización} sea sencilla:

\begin{quote}
  « El método axiomático no es otra cosa que el arte de redactar textos
  cuya formalización sea fácil de concebir.  No es una invención nueva,
  pero su empleo sistemático como instrumento de descubrimiento es uno
  de los rasgos originales de la matemática contemporánea. »
  \hfill \textit{(Éléments de mathématique, p.~9)}
\end{quote}

De aquí se desprende una guía práctica para la construcción de pruebas:

\begin{enumerate}
  \item \textbf{Definir previamente} todos los símbolos que aparecen.
  \item \textbf{Expresar la hipótesis} en forma de proposición bien
        estructurada (p. \ref{sec:proposiciones}).
  \item \textbf{Utilizar los conectores lógicos} (\(\neg,\wedge,\vee,\Rightarrow,\Leftrightarrow\))
        exactamente según sus definiciones (\ref{sec:negacion-conjuncion}, \ref{sec:implicacion-equivalencia}).
  \item \textbf{Aplicar cuantificadores} (\(\forall,\exists\)) respetando las reglas de
        negación (\ref{sec:negacion-cuantificados}).
  \item \textbf{Concluir} señalando que la cadena de inferencias corresponde
        a una \emph{implicación lógica} cuya verdad queda garantizada
        por el esquema axiomático.
\end{enumerate}

En otras palabras, una demostración es una \emph{traducción} de la
implicación lógica que queremos validar a una *cadena de pasos* que
cumple los criterios de \textbf{definición previa} y \textbf{formalidad}
establecidos en el capítulo anterior.

\begin{cita}
  « La verificación de un texto formalizado no pide más que una
  atención mecánica; cuando la formalización es posible, el
  razonamiento queda completamente determinado por la sintaxis. »
  \hfill \textit{(Éléments de mathématique, p.~8)}
\end{cita}

\section{Implicación y equivalencia}
\label{sec:implicacion-equivalencia}

Los conectores de implicación y equivalencia son fundamentales para formular teoremas y definiciones.

\begin{definicion}[Implicación]
Si $P$ y $Q$ son proposiciones, la \emph{implicación}
\[
  P\Rightarrow Q
\]
se lee ``si $P$, entonces $Q$''. La proposición $P$ se llama \emph{hipótesis} o \emph{antecedente}, y la proposición $Q$ se llama \emph{conclusión} o \emph{consecuente}.
\end{definicion}

Demostrar una implicación $P\Rightarrow Q$ significa justificar que siempre que $P$ sea verdadera, también lo será $Q$.

\begin{ejemplo}
La afirmación
\[
  n \text{ es múltiplo de }4 \Rightarrow n \text{ es par}
\]
es verdadera para todo número entero $n$. En efecto, si $n$ es múltiplo de $4$, existe un entero $k$ tal que $n=4k$. Entonces $n=2(2k)$, luego $n$ es par.
\end{ejemplo}

\begin{definicion}[Recíproca]
Dada una implicación $P\Rightarrow Q$, su \emph{recíproca} es la implicación
\[
  Q\Rightarrow P.
\]
\end{definicion}

\begin{advertencia}
Una implicación y su recíproca no tienen por qué tener el mismo valor de verdad. Que $P\Rightarrow Q$ sea verdadera no implica que $Q\Rightarrow P$ sea verdadera.
\end{advertencia}

\begin{ejemplo}
La implicación
\[
  n \text{ es múltiplo de }4 \Rightarrow n \text{ es par}
\]
es verdadera. Sin embargo, su recíproca
\[
  n \text{ es par } \Rightarrow n \text{ es múltiplo de }4
\]
es falsa. Por ejemplo, $6$ es par, pero no es múltiplo de $4$.
\end{ejemplo}

\begin{definicion}[Equivalencia]
Si $P$ y $Q$ son proposiciones, la \emph{equivalencia}
\[
  P\Leftrightarrow Q
\]
se lee ``$P$ si y solo si $Q$''. Significa que se cumplen simultáneamente las dos implicaciones
\[
  P\Rightarrow Q
  \qquad \text{y} \qquad
  Q\Rightarrow P.
\]
\end{definicion}

\begin{ejemplo}
Para un número entero $n$, se tiene
\[
  n \text{ es par } \Leftrightarrow \text{ existe } k\in\mathbb{Z} \text{ tal que } n=2k.
\]
Esta equivalencia expresa la definición usual de número par.
\end{ejemplo}

\section{Cuantificadores}

En matemáticas se formulan con frecuencia enunciados que hablan de todos los objetos de una colección o de la existencia de al menos un objeto con cierta propiedad. Para ello se utilizan cuantificadores.

\begin{definicion}[Cuantificador universal]
El símbolo $\forall$ se llama \emph{cuantificador universal} y se lee ``para todo''. Si $A$ es un conjunto y $P(x)$ es una propiedad que depende de un elemento $x\in A$, la expresión
\[
  \forall x\in A,\ P(x)
\]
significa que la propiedad $P(x)$ se cumple para todos los elementos $x$ del conjunto $A$.
\end{definicion}

\begin{ejemplo}
La proposición
\[
  \forall n\in\mathbb{N},\ n+1>n
\]
afirma que todo número natural es menor que su sucesor. Aquí $\mathbb{N}$ denota el conjunto de los números naturales.
\end{ejemplo}

\begin{definicion}[Cuantificador existencial]
El símbolo $\exists$ se llama \emph{cuantificador existencial} y se lee ``existe''. Si $A$ es un conjunto y $P(x)$ es una propiedad que depende de un elemento $x\in A$, la expresión
\[
  \exists x\in A \text{ tal que } P(x)
\]
significa que hay al menos un elemento $x$ de $A$ para el cual la propiedad $P(x)$ se cumple.
\end{definicion}

\begin{ejemplo}
La proposición
\[
  \exists n\in\mathbb{N} \text{ tal que } n^2=9
\]
es verdadera, porque $3\in\mathbb{N}$ y $3^2=9$.
\end{ejemplo}

\begin{definicion}[Cuantificador existencial único]
El símbolo $\exists!$ se lee ``existe un único''. La expresión
\[
  \exists! x\in A \text{ tal que } P(x)
\]
significa que existe un elemento $x\in A$ que cumple $P(x)$ y que no hay ningún otro elemento de $A$ con esa propiedad.
\end{definicion}

\begin{ejemplo}
La proposición
\[
  \exists! n\in\mathbb{N} \text{ tal que } n+1=1
\]
es verdadera si consideramos $0\in\mathbb{N}$, porque el único número natural que la satisface es $n=0$.
\end{ejemplo}

\begin{advertencia}
La notación $\mathbb{N}$ puede variar entre textos. Algunos autores incluyen el $0$ en los números naturales y otros no. En este curso se indicará explícitamente qué convención se adopta antes de utilizarla en resultados donde sea relevante.
\end{advertencia}

\section{Negación de enunciados cuantificados}
\label{sec:negacion-cuantificados}

Negar correctamente enunciados con cuantificadores es una habilidad esencial para comprender demostraciones, formular contraejemplos y aplicar razonamientos por contradicción.

\begin{proposicion}[Negación del cuantificador universal]
Sea $A$ un conjunto y sea $P(x)$ una propiedad definida para los elementos $x\in A$. La negación de
\[
  \forall x\in A,\ P(x)
\]
es
\[
  \exists x\in A \text{ tal que } \neg P(x).
\]
\end{proposicion}

En palabras: negar que una propiedad se cumple para todos los elementos de $A$ equivale a afirmar que existe al menos un elemento de $A$ para el que la propiedad no se cumple.

\begin{ejemplo}
La negación de
\[
  \forall n\in\mathbb{N},\ n \text{ es par}
\]
es
\[
  \exists n\in\mathbb{N} \text{ tal que } n \text{ no es par}.
\]
\end{ejemplo}

\begin{proposicion}[Negación del cuantificador existencial]
Sea $A$ un conjunto y sea $P(x)$ una propiedad definida para los elementos $x\in A$. La negación de
\[
  \exists x\in A \text{ tal que } P(x)
\]
es
\[
  \forall x\in A,\ \neg P(x).
\]
\end{proposicion}

En palabras: negar que exista un elemento con cierta propiedad equivale a afirmar que ningún elemento tiene esa propiedad.

\begin{ejemplo}
La negación de
\[
  \exists x\in\mathbb{R} \text{ tal que } x^2<0
\]
es
\[
  \forall x\in\mathbb{R},\ x^2\geq 0.
\]
\end{ejemplo}

\begin{advertencia}
Al negar un enunciado cuantificado hay que cambiar el cuantificador y negar la propiedad. Un error frecuente es negar solo la propiedad y conservar el cuantificador original.
\end{advertencia}

\begin{ejemplo}
La negación correcta de
\[
  \forall x\in\mathbb{R},\ x^2>0
\]
es
\[
  \exists x\in\mathbb{R} \text{ tal que } x^2\leq 0.
\]
No es correcto escribir
\[
  \forall x\in\mathbb{R},\ x^2\leq 0.
\]
\end{ejemplo}

\section{Condiciones necesarias y suficientes}

Las expresiones ``condición necesaria'' y ``condición suficiente'' aparecen con frecuencia en definiciones, criterios y teoremas. Conviene fijar su significado con precisión.

\begin{definicion}[Condición suficiente]
Sean $P$ y $Q$ dos proposiciones. Si se cumple la implicación
\[
  P\Rightarrow Q,
\]
decimos que $P$ es una \emph{condición suficiente} para $Q$.
\end{definicion}

Es decir, saber que $P$ es verdadera basta para concluir que $Q$ es verdadera.

\begin{definicion}[Condición necesaria]
Sean $P$ y $Q$ dos proposiciones. Si se cumple la implicación
\[
  P\Rightarrow Q,
\]
decimos también que $Q$ es una \emph{condición necesaria} para $P$.
\end{definicion}

Es decir, para que $P$ pueda ser verdadera es necesario que $Q$ sea verdadera.

\begin{ejemplo}
Para un número entero $n$, la afirmación
\[
  n \text{ es múltiplo de }4 \Rightarrow n \text{ es par}
\]
indica que ser múltiplo de $4$ es condición suficiente para ser par. También indica que ser par es condición necesaria para ser múltiplo de $4$.
\end{ejemplo}

\begin{definicion}[Condición necesaria y suficiente]
Sean $P$ y $Q$ dos proposiciones. Decimos que $P$ es una \emph{condición necesaria y suficiente} para $Q$ cuando se cumple
\[
  P\Leftrightarrow Q.
\]
\end{definicion}

\begin{ejemplo}
Para un número entero $n$, ser par es condición necesaria y suficiente para que exista un entero $k$ tal que $n=2k$.
\end{ejemplo}

\section{Métodos de demostración}

Una demostración es un razonamiento que justifica la verdad de un enunciado matemático a partir de definiciones, resultados ya conocidos y reglas lógicas aceptadas.

En esta sección se presentan varios métodos básicos de demostración. No son recetas mecánicas, pero proporcionan estructuras útiles para organizar el razonamiento.

\subsection{Demostración directa}

La demostración directa se utiliza habitualmente para probar una implicación $P\Rightarrow Q$.

\begin{definicion}[Demostración directa]
Una \emph{demostración directa} de $P\Rightarrow Q$ consiste en suponer que $P$ es verdadera y deducir, mediante pasos justificados, que $Q$ también es verdadera.
\end{definicion}

\begin{ejemplo}
Demostremos que si $n$ es un número entero par, entonces $n^2$ es par.

Supongamos que $n$ es par. Por definición de número par, existe $k\in\mathbb{Z}$ tal que
\[
  n=2k.
\]
Entonces
\[
  n^2=(2k)^2=4k^2=2(2k^2).
\]
Como $2k^2\in\mathbb{Z}$, se deduce que $n^2$ es par.
\end{ejemplo}

\subsection{Demostración por contraposición}

Antes de definir el método, introducimos la contrapositiva de una implicación.

\begin{definicion}[Contrapositiva]
Dada una implicación
\[
  P\Rightarrow Q,
\]
su \emph{contrapositiva} es la implicación
\[
  \neg Q\Rightarrow \neg P.
\]
\end{definicion}

Una implicación y su contrapositiva son lógicamente equivalentes: demostrar una es equivalente a demostrar la otra.

\begin{definicion}[Demostración por contraposición]
Una \emph{demostración por contraposición} de $P\Rightarrow Q$ consiste en demostrar la implicación equivalente
\[
  \neg Q\Rightarrow \neg P.
\]
\end{definicion}

\begin{ejemplo}
Demostremos que si $n^2$ es par, entonces $n$ es par.

En lugar de demostrar directamente esta implicación, probamos su contrapositiva: si $n$ no es par, entonces $n^2$ no es par.

Supongamos que $n$ no es par. Entonces $n$ es impar, es decir, existe $k\in\mathbb{Z}$ tal que
\[
  n=2k+1.
\]
Por tanto,
\[
  n^2=(2k+1)^2=4k^2+4k+1=2(2k^2+2k)+1.
\]
Como $2k^2+2k\in\mathbb{Z}$, se deduce que $n^2$ es impar. Luego $n^2$ no es par. Por contraposición, si $n^2$ es par, entonces $n$ es par.
\end{ejemplo}

\subsection{Demostración por contradicción}

\begin{definicion}[Demostración por contradicción]
Una \emph{demostración por contradicción} de una proposición $P$ consiste en suponer que $P$ es falsa, es decir, suponer $\neg P$, y deducir una contradicción. Al obtener una contradicción, se concluye que $P$ debe ser verdadera.
\end{definicion}

Una contradicción puede aparecer, por ejemplo, cuando se deduce simultáneamente una proposición y su negación, o cuando se obtiene una afirmación manifiestamente falsa, como $0=1$.

\begin{ejemplo}
Demostremos que no existe un número entero $n$ tal que $n$ sea par e impar simultáneamente.

Supongamos, por contradicción, que existe $n\in\mathbb{Z}$ que es par e impar. Como $n$ es par, existe $a\in\mathbb{Z}$ tal que $n=2a$. Como $n$ es impar, existe $b\in\mathbb{Z}$ tal que $n=2b+1$. Por tanto,
\[
  2a=2b+1.
\]
De aquí se obtiene
\[
  2(a-b)=1.
\]
El lado izquierdo es un número par, mientras que $1$ no es par. Esto es una contradicción. Por tanto, no existe tal número entero $n$.
\end{ejemplo}

\begin{advertencia}
En una demostración por contradicción no basta con suponer lo contrario y obtener una dificultad. Hay que obtener una contradicción precisa con una definición, una hipótesis, un resultado conocido o una propiedad elemental aceptada.
\end{advertencia}

\subsection{Demostración por casos}

\begin{definicion}[Demostración por casos]
Una \emph{demostración por casos} consiste en dividir la situación en un número finito de posibilidades que cubren todos los casos posibles, y demostrar el resultado en cada una de ellas.
\end{definicion}

\begin{ejemplo}
Demostremos que para todo número entero $n$, el producto $n(n+1)$ es par.

Sea $n\in\mathbb{Z}$. Hay dos casos posibles:

\begin{enumerate}
  \item Si $n$ es par, entonces existe $k\in\mathbb{Z}$ tal que $n=2k$. Por tanto,
  \[
    n(n+1)=2k(n+1),
  \]
  luego $n(n+1)$ es par.

  \item Si $n$ es impar, entonces $n+1$ es par. Por tanto, existe $k\in\mathbb{Z}$ tal que $n+1=2k$. Entonces
  \[
    n(n+1)=n\cdot 2k=2(nk),
  \]
  luego $n(n+1)$ es par.
\end{enumerate}

En ambos casos, $n(n+1)$ es par. Por tanto, para todo $n\in\mathbb{Z}$, el número $n(n+1)$ es par.
\end{ejemplo}

\begin{advertencia}
En una demostración por casos hay que comprobar que los casos considerados cubren todas las posibilidades. Si falta algún caso, la demostración no es completa.
\end{advertencia}

\subsection{Demostración de equivalencias}

Para demostrar una equivalencia $P\Leftrightarrow Q$ hay que demostrar dos implicaciones.

\begin{definicion}[Demostración de una equivalencia]
Una demostración de
\[
  P\Leftrightarrow Q
\]
consiste en demostrar, por separado, las implicaciones
\[
  P\Rightarrow Q
  \qquad \text{y} \qquad
  Q\Rightarrow P.
\]
\end{definicion}

\begin{ejemplo}
Demostremos que, para un número entero $n$, se cumple
\[
  n \text{ es par } \Leftrightarrow n^2 \text{ es par}.
\]

Primero demostramos que si $n$ es par, entonces $n^2$ es par. Esta implicación ya se demostró mediante demostración directa.

Ahora demostramos la implicación recíproca: si $n^2$ es par, entonces $n$ es par. Esta implicación ya se demostró por contraposición.

Como se han probado las dos implicaciones, queda demostrada la equivalencia.
\end{ejemplo}

\begin{erroresfrecuentes}
  \item Confundir una implicación con su recíproca. De $P\Rightarrow Q$ no se deduce automáticamente $Q\Rightarrow P$.
  \item Negar incorrectamente un enunciado con cuantificadores. La negación de ``para todo'' es ``existe un caso en que no'', y la negación de ``existe'' es ``para todo no''.
  \item Usar un símbolo lógico antes de explicar su significado. Por ejemplo, no debe utilizarse $\forall$, $\exists$, $\Rightarrow$ o $\Leftrightarrow$ sin haber indicado previamente cómo se leen y qué significan.
  \item Probar una afirmación universal mediante ejemplos. Los ejemplos pueden orientar, pero no sustituyen una demostración general.
  \item Creer que una demostración por contradicción consiste solo en decir ``supongamos lo contrario''. Es necesario obtener una contradicción explícita.
\end{erroresfrecuentes}

\begin{seminario}
  \item Traducción de frases al lenguaje lógico.
  \item Negación de enunciados con uno o varios cuantificadores.
  \item Demostraciones directas elementales.
  \item Identificación de hipótesis y conclusiones en una implicación.
  \item Comparación entre una implicación, su recíproca y su contrapositiva.
  \item Redacción de demostraciones breves cuidando la definición previa de todos los símbolos relevantes.
\end{seminario}

\begin{ejerciciospropuestos}
  \item Escribe la negación de los siguientes enunciados:
  \begin{enumerate}
    \item Para todo $n\in\mathbb{N}$, $n+1>n$.
    \item Existe $x\in\mathbb{R}$ tal que $x^2=2$.
    \item Para todo $x\in\mathbb{R}$, existe $y\in\mathbb{R}$ tal que $y>x$.
    \item Existe $n\in\mathbb{Z}$ tal que $n$ es par y $n$ es primo.
  \end{enumerate}

  \item Distingue entre condición necesaria, condición suficiente y condición necesaria y suficiente en los siguientes ejemplos:
  \begin{enumerate}
    \item Ser múltiplo de $4$ y ser par.
    \item Ser cuadrado y ser rectángulo.
    \item Ser número entero par y ser de la forma $2k$ para algún $k\in\mathbb{Z}$.
  \end{enumerate}

  \item Para cada una de las siguientes implicaciones, escribe su recíproca y su contrapositiva:
  \begin{enumerate}
    \item Si $n$ es múltiplo de $6$, entonces $n$ es múltiplo de $3$.
    \item Si $x>2$, entonces $x^2>4$.
    \item Si un número entero es múltiplo de $10$, entonces termina en $0$.
  \end{enumerate}

  \item Demuestra directamente que si $n$ es un número entero impar, entonces $n^2$ es impar.

  \item Demuestra por contraposición que si $n^2$ es impar, entonces $n$ es impar.

  \item Demuestra por casos que para todo $n\in\mathbb{Z}$, el número $n^2+n$ es par.

  \item Escribe un ejemplo de una implicación verdadera cuya recíproca sea falsa.

  \item Escribe un ejemplo de una equivalencia matemática y demuestra las dos implicaciones que contiene.
\end{ejerciciospropuestos}

\resumen{El capítulo introduce el lenguaje lógico elemental necesario para formular y demostrar enunciados matemáticos. Se han definido las proposiciones, los conectores lógicos, los cuantificadores y los métodos básicos de demostración. La idea central es que razonar matemáticamente exige identificar con precisión las hipótesis, las conclusiones y el significado de cada símbolo utilizado.}

% ------------------------------------------------------------
\chapter{Conjuntos}

\begin{objetivos}
  \item Introducir el lenguaje elemental de la teoría de conjuntos.
  \item Definir pertenencia, inclusión, unión, intersección, diferencia y complementario.
  \item Relacionar operaciones conjuntistas con conectores lógicos.
  \item Aprender a demostrar igualdades entre conjuntos.
\end{objetivos}

\section{Motivación}

La teoría de conjuntos proporciona uno de los lenguajes básicos de la matemática moderna. Muchas nociones que aparecerán a lo largo del curso, como aplicación, relación, estructura algebraica, número natural, número racional o espacio de soluciones, se describen de forma precisa mediante conjuntos y operaciones entre conjuntos.

En este capítulo no se pretende desarrollar una teoría axiomática de conjuntos. El objetivo es más elemental: introducir el vocabulario y la notación necesarios para trabajar con colecciones de objetos de forma rigurosa.

\begin{convencion}[Símbolos antes de su uso]
Todo símbolo conjuntista que aparezca en un enunciado debe haber sido definido previamente. En particular, antes de escribir expresiones como \(x\in A\), \(A\subseteq B\) o \(A\cup B\), se deberá haber indicado qué objetos representan \(x\), \(A\) y \(B\), y qué significado tienen los símbolos \(\in\), \(\subseteq\) y \(\cup\).
\end{convencion}

\section{Conjuntos y elementos}

De manera intuitiva, un \emph{conjunto} es una colección de objetos bien determinados. Los objetos que forman parte de un conjunto se llaman \emph{elementos} del conjunto.

\begin{definicion}[Conjunto y elemento]
Un \emph{conjunto} es una colección de objetos considerados como una totalidad. Si \(A\) es un conjunto, los objetos que pertenecen a \(A\) se llaman \emph{elementos} de \(A\).
\end{definicion}

La notación habitual consiste en denotar los conjuntos mediante letras mayúsculas, como \(A\), \(B\), \(C\), y sus elementos mediante letras minúsculas, como \(x\), \(y\), \(z\). Esta convención no es obligatoria, pero será utilizada de forma sistemática en estas notas.

\begin{ejemplo}
Sea \(A\) el conjunto formado por los números \(1\), \(2\) y \(3\). Escribimos
\[
  A=\{1,2,3\}.
\]
Los elementos de \(A\) son \(1\), \(2\) y \(3\).
\end{ejemplo}

Un conjunto puede describirse enumerando sus elementos, cuando esto es posible, o mediante una propiedad que caracteriza a sus elementos.

\begin{ejemplo}
El conjunto de los números naturales pares menores que \(10\) puede escribirse como
\[
  \{2,4,6,8\}.
\]
También puede describirse como
\[
  \{n\in\mathbb{N}: n \text{ es par y } n<10\}.
\]
En esta expresión, \(\mathbb{N}\) denota el conjunto de los números naturales y el símbolo ``:'' se lee como ``tal que''.
\end{ejemplo}

\begin{advertencia}
Para que una colección pueda considerarse un conjunto en el sentido elemental que utilizaremos en este curso, debe estar claro si un objeto dado pertenece o no pertenece a ella. Expresiones como ``el conjunto de los números grandes'' son ambiguas si no se precisa qué significa ``grande''.
\end{advertencia}

\section{Pertenencia e inclusión}

El primer símbolo fundamental de la teoría de conjuntos es el símbolo de pertenencia.

\begin{definicion}[Pertenencia]
Sean \(A\) un conjunto y \(x\) un objeto. Escribimos
\[
  x\in A
\]
para indicar que \(x\) es un elemento de \(A\). Escribimos
\[
  x\notin A
\]
para indicar que \(x\) no es un elemento de \(A\).
\end{definicion}

\begin{ejemplo}
Si \(A=\{1,2,3\}\), entonces
\[
  2\in A,
  \qquad
  5\notin A.
\]
\end{ejemplo}

El segundo símbolo fundamental es el símbolo de inclusión.

\begin{definicion}[Inclusión]
Sean \(A\) y \(B\) dos conjuntos. Decimos que \(A\) está \emph{incluido} en \(B\), o que \(A\) es un \emph{subconjunto} de \(B\), si todo elemento de \(A\) es también elemento de \(B\). En ese caso escribimos
\[
  A\subseteq B.
\]
En símbolos lógicos:
\[
  A\subseteq B
  \quad\Longleftrightarrow\quad
  \forall x\,(x\in A \Rightarrow x\in B).
\]
\end{definicion}

Aquí \(\forall\) significa ``para todo'', y \(\Rightarrow\) significa ``implica''. Es importante observar que estos símbolos se introdujeron en el capítulo anterior antes de ser utilizados aquí.

\begin{ejemplo}
Si \(A=\{1,2\}\) y \(B=\{1,2,3\}\), entonces \(A\subseteq B\), porque todo elemento de \(A\) es también elemento de \(B\).
\end{ejemplo}

\begin{advertencia}
No debe confundirse pertenencia con inclusión. La expresión \(x\in A\) afirma que \(x\) es un elemento del conjunto \(A\). La expresión \(A\subseteq B\) afirma que todo elemento de \(A\) es también elemento de \(B\). Por tanto, \(\in\) relaciona un objeto con un conjunto, mientras que \(\subseteq\) relaciona dos conjuntos.
\end{advertencia}

\section{Subconjuntos y conjunto vacío}

La relación de inclusión permite comparar conjuntos. Si \(A\subseteq B\), podemos decir que \(A\) es más pequeño o igual que \(B\) desde el punto de vista de sus elementos, aunque esta interpretación debe manejarse con cuidado en conjuntos infinitos.

\begin{definicion}[Igualdad de conjuntos]
Sean \(A\) y \(B\) dos conjuntos. Decimos que \(A\) y \(B\) son iguales, y escribimos \(A=B\), si tienen exactamente los mismos elementos. Equivalentemente,
\[
  A=B
  \quad\Longleftrightarrow\quad
  (A\subseteq B \text{ y } B\subseteq A).
\]
\end{definicion}

Esta caracterización es la base del método de \emph{doble inclusión}, que se estudiará al final del capítulo.

\begin{definicion}[Conjunto vacío]
El \emph{conjunto vacío} es el conjunto que no tiene ningún elemento. Se denota por
\[
  \varnothing.
\]
\end{definicion}

\begin{proposicion}
Para todo conjunto \(A\), se cumple
\[
  \varnothing\subseteq A.
\]
\end{proposicion}

\begin{proof}
Para probar que \(\varnothing\subseteq A\), debemos demostrar que todo elemento de \(\varnothing\) pertenece a \(A\). Pero \(\varnothing\) no tiene elementos. Por tanto, no existe ningún elemento de \(\varnothing\) que pueda violar la condición. En consecuencia, \(\varnothing\subseteq A\).
\end{proof}

\begin{advertencia}
El conjunto vacío \(\varnothing\) no debe confundirse con el conjunto \(\{\varnothing\}\). El primero no tiene elementos. El segundo tiene un elemento, que es precisamente \(\varnothing\).
\end{advertencia}

\section{Operaciones con conjuntos}

A partir de conjuntos dados se pueden construir nuevos conjuntos mediante operaciones. Las operaciones más importantes en este curso serán la unión, la intersección, la diferencia y el complementario.

En toda esta sección, sean \(A\) y \(B\) conjuntos contenidos en un conjunto de referencia \(U\), llamado \emph{universo} cuando sea necesario hablar de complementarios.

\subsection{Unión}

\begin{definicion}[Unión]
Sean \(A\) y \(B\) dos conjuntos. La \emph{unión} de \(A\) y \(B\) es el conjunto formado por los elementos que pertenecen a \(A\), a \(B\), o a ambos. Se denota por
\[
  A\cup B.
\]
Formalmente,
\[
  A\cup B = \{x : x\in A \text{ o } x\in B\}.
\]
\end{definicion}

La palabra ``o'' se interpreta aquí como una disyunción inclusiva: un elemento de \(A\cup B\) puede pertenecer a \(A\), a \(B\), o a ambos.

\begin{ejemplo}
Si \(A=\{1,2,3\}\) y \(B=\{3,4,5\}\), entonces
\[
  A\cup B=\{1,2,3,4,5\}.
\]
\end{ejemplo}

\subsection{Intersección}

\begin{definicion}[Intersección]
Sean \(A\) y \(B\) dos conjuntos. La \emph{intersección} de \(A\) y \(B\) es el conjunto formado por los elementos que pertenecen simultáneamente a \(A\) y a \(B\). Se denota por
\[
  A\cap B.
\]
Formalmente,
\[
  A\cap B = \{x : x\in A \text{ y } x\in B\}.
\]
\end{definicion}

\begin{ejemplo}
Si \(A=\{1,2,3\}\) y \(B=\{3,4,5\}\), entonces
\[
  A\cap B=\{3\}.
\]
\end{ejemplo}

\begin{definicion}[Conjuntos disjuntos]
Sean \(A\) y \(B\) dos conjuntos. Decimos que \(A\) y \(B\) son \emph{disjuntos} si no tienen elementos comunes, es decir, si
\[
  A\cap B=\varnothing.
\]
\end{definicion}

\subsection{Diferencia}

\begin{definicion}[Diferencia de conjuntos]
Sean \(A\) y \(B\) dos conjuntos. La \emph{diferencia} de \(A\) y \(B\) es el conjunto formado por los elementos que pertenecen a \(A\) y no pertenecen a \(B\). Se denota por
\[
  A\setminus B.
\]
Formalmente,
\[
  A\setminus B = \{x : x\in A \text{ y } x\notin B\}.
\]
\end{definicion}

\begin{ejemplo}
Si \(A=\{1,2,3\}\) y \(B=\{3,4,5\}\), entonces
\[
  A\setminus B=\{1,2\},
  \qquad
  B\setminus A=\{4,5\}.
\]
Por tanto, en general, \(A\setminus B\) no coincide con \(B\setminus A\).
\end{ejemplo}

\subsection{Complementario}

Para definir el complementario de un conjunto es necesario fijar previamente un conjunto de referencia.

\begin{definicion}[Complementario]
Sea \(U\) un conjunto de referencia y sea \(A\subseteq U\). El \emph{complementario} de \(A\) en \(U\) es el conjunto formado por los elementos de \(U\) que no pertenecen a \(A\). Se denota por
\[
  A^c
\]
cuando el conjunto de referencia \(U\) está claro. Formalmente,
\[
  A^c = \{x\in U : x\notin A\}.
\]
\end{definicion}

\begin{ejemplo}
Sea \(U=\{1,2,3,4,5\}\) y sea \(A=\{1,3,5\}\). Entonces
\[
  A^c=\{2,4\}.
\]
\end{ejemplo}

\begin{advertencia}
No tiene sentido hablar del complementario de \(A\) sin saber respecto de qué conjunto de referencia se toma el complementario. El conjunto \(A^c\) depende del universo \(U\).
\end{advertencia}

\section{Relación entre operaciones conjuntistas y conectores lógicos}

Las operaciones con conjuntos están estrechamente relacionadas con los conectores lógicos estudiados en el capítulo anterior.

Sean \(A\) y \(B\) subconjuntos de un conjunto de referencia \(U\), y sea \(x\in U\). Entonces:
\[
\begin{array}{rcl}
  x\in A\cup B &\Longleftrightarrow& x\in A \text{ o } x\in B,\\[2mm]
  x\in A\cap B &\Longleftrightarrow& x\in A \text{ y } x\in B,\\[2mm]
  x\in A\setminus B &\Longleftrightarrow& x\in A \text{ y } x\notin B,\\[2mm]
  x\in A^c &\Longleftrightarrow& x\notin A.
\end{array}
\]

Estas equivalencias permiten transformar problemas conjuntistas en problemas lógicos, y viceversa.

\begin{ejemplo}
La igualdad
\[
  (A\cap B)^c=A^c\cup B^c
\]
corresponde a la ley lógica
\[
  \neg(P\wedge Q) \Longleftrightarrow (\neg P)\vee(\neg Q),
\]
conocida como una de las leyes de De Morgan.
\end{ejemplo}

\section{Producto cartesiano}

Hasta ahora hemos trabajado con elementos individuales. En muchas situaciones matemáticas es necesario considerar pares ordenados.

\begin{definicion}[Par ordenado]
Un \emph{par ordenado} es una expresión de la forma \((a,b)\), donde \(a\) es la primera componente y \(b\) es la segunda componente. Dos pares ordenados \((a,b)\) y \((c,d)\) son iguales si y solo si
\[
  a=c \quad \text{y} \quad b=d.
\]
\end{definicion}

\begin{definicion}[Producto cartesiano]
Sean \(A\) y \(B\) dos conjuntos. El \emph{producto cartesiano} de \(A\) y \(B\) es el conjunto formado por todos los pares ordenados cuya primera componente pertenece a \(A\) y cuya segunda componente pertenece a \(B\). Se denota por
\[
  A\times B.
\]
Formalmente,
\[
  A\times B=\{(a,b): a\in A \text{ y } b\in B\}.
\]
\end{definicion}

\begin{ejemplo}
Si \(A=\{1,2\}\) y \(B=\{a,b\}\), entonces
\[
  A\times B=\{(1,a),(1,b),(2,a),(2,b)\}.
\]
\end{ejemplo}

\begin{advertencia}
En general, \(A\times B\) no coincide con \(B\times A\). Por ejemplo, si \(1\neq a\), entonces \((1,a)\) no es el mismo par ordenado que \((a,1)\).
\end{advertencia}

\section{Conjunto de las partes}

Además de considerar los elementos de un conjunto, también es importante considerar todos sus subconjuntos.

\begin{definicion}[Conjunto de las partes]
Sea \(A\) un conjunto. El \emph{conjunto de las partes} de \(A\), denotado por
\[
  \mathcal{P}(A),
\]
es el conjunto formado por todos los subconjuntos de \(A\). Formalmente,
\[
  \mathcal{P}(A)=\{B: B\subseteq A\}.
\]
\end{definicion}

\begin{ejemplo}
Si \(A=\{1,2\}\), entonces
\[
  \mathcal{P}(A)=\{\varnothing,\{1\},\{2\},\{1,2\}\}.
\]
\end{ejemplo}

\begin{proposicion}
Si \(A\) es un conjunto finito con \(n\) elementos, entonces \(\mathcal{P}(A)\) tiene \(2^n\) elementos.
\end{proposicion}

\begin{proof}
Cada subconjunto de \(A\) se obtiene decidiendo, para cada elemento de \(A\), si se incluye o no se incluye en el subconjunto. Para cada uno de los \(n\) elementos hay dos posibilidades. Por tanto, el número total de subconjuntos es
\[
  2\cdot 2\cdots 2 = 2^n.
\]
\end{proof}

\section{Familias de conjuntos}

En ocasiones se trabaja no solo con dos conjuntos, sino con muchos conjuntos a la vez. Para ello se usa la noción de familia de conjuntos.

\begin{definicion}[Familia de conjuntos]
Sea \(I\) un conjunto, llamado \emph{conjunto de índices}. Una \emph{familia de conjuntos indexada por} \(I\) es una colección de conjuntos
\[
  (A_i)_{i\in I},
\]
donde a cada índice \(i\in I\) le corresponde un conjunto \(A_i\).
\end{definicion}

\begin{ejemplo}
Si \(I=\{1,2,3\}\), una familia de conjuntos indexada por \(I\) puede escribirse como
\[
  (A_1,A_2,A_3).
\]
Por ejemplo,
\[
  A_1=\{1,2\},\qquad A_2=\{2,3\},\qquad A_3=\{3,4\}.
\]
\end{ejemplo}

\begin{definicion}[Unión e intersección de una familia]
Sea \((A_i)_{i\in I}\) una familia de conjuntos. La \emph{unión} de la familia se define por
\[
  \bigcup_{i\in I} A_i = \{x : \text{existe } i\in I \text{ tal que } x\in A_i\}.
\]
La \emph{intersección} de la familia se define por
\[
  \bigcap_{i\in I} A_i = \{x : \text{para todo } i\in I,\ x\in A_i\}.
\]
\end{definicion}

\begin{ejemplo}
Con los conjuntos
\[
  A_1=\{1,2\},\qquad A_2=\{2,3\},\qquad A_3=\{2,4\},
\]
se tiene
\[
  \bigcup_{i=1}^3 A_i = \{1,2,3,4\},
  \qquad
  \bigcap_{i=1}^3 A_i = \{2\}.
\]
\end{ejemplo}

\section{Demostración de identidades conjuntistas}

Una \emph{identidad conjuntista} es una igualdad entre conjuntos que se cumple bajo ciertas hipótesis. Por ejemplo,
\[
  A\cap(B\cup C)=(A\cap B)\cup(A\cap C)
\]
es una identidad conjuntista válida para cualesquiera conjuntos \(A\), \(B\) y \(C\).

El método más importante para demostrar igualdades entre conjuntos es el método de doble inclusión.

\begin{metodo}[Doble inclusión]
Para demostrar que dos conjuntos \(X\) e \(Y\) son iguales, basta demostrar las dos inclusiones:
\[
  X\subseteq Y
  \qquad\text{y}\qquad
  Y\subseteq X.
\]
Es decir, se prueba primero que todo elemento de \(X\) pertenece a \(Y\), y después que todo elemento de \(Y\) pertenece a \(X\).
\end{metodo}

\begin{proposicion}[Distributividad de la intersección respecto de la unión]
Sean \(A\), \(B\) y \(C\) conjuntos. Entonces
\[
  A\cap(B\cup C)=(A\cap B)\cup(A\cap C).
\]
\end{proposicion}

\begin{proof}
Demostraremos la igualdad mediante doble inclusión.

Primero probamos que
\[
  A\cap(B\cup C)\subseteq (A\cap B)\cup(A\cap C).
\]
Sea \(x\in A\cap(B\cup C)\). Entonces \(x\in A\) y \(x\in B\cup C\). Como \(x\in B\cup C\), se tiene \(x\in B\) o \(x\in C\).

Si \(x\in B\), entonces \(x\in A\cap B\), y por tanto
\[
  x\in (A\cap B)\cup(A\cap C).
\]
Si \(x\in C\), entonces \(x\in A\cap C\), y por tanto
\[
  x\in (A\cap B)\cup(A\cap C).
\]
En ambos casos, \(x\in (A\cap B)\cup(A\cap C)\). Por tanto,
\[
  A\cap(B\cup C)\subseteq (A\cap B)\cup(A\cap C).
\]

Ahora probamos la inclusión contraria:
\[
  (A\cap B)\cup(A\cap C)\subseteq A\cap(B\cup C).
\]
Sea \(x\in (A\cap B)\cup(A\cap C)\). Entonces \(x\in A\cap B\) o \(x\in A\cap C\).

Si \(x\in A\cap B\), entonces \(x\in A\) y \(x\in B\). En particular, \(x\in B\cup C\). Por tanto, \(x\in A\cap(B\cup C)\).

Si \(x\in A\cap C\), entonces \(x\in A\) y \(x\in C\). En particular, \(x\in B\cup C\). Por tanto, \(x\in A\cap(B\cup C)\).

En ambos casos, \(x\in A\cap(B\cup C)\). Por tanto,
\[
  (A\cap B)\cup(A\cap C)\subseteq A\cap(B\cup C).
\]

Como se cumplen ambas inclusiones, concluimos que
\[
  A\cap(B\cup C)=(A\cap B)\cup(A\cap C).
\]
\end{proof}

\begin{erroresfrecuentes}
  \item Confundir \(x\in A\) con \(x\subseteq A\).
  \item Confundir \(A\in B\) con \(A\subseteq B\).
  \item Escribir el complementario \(A^c\) sin haber fijado previamente el conjunto de referencia.
  \item Creer que \(A\setminus B = B\setminus A\) en general.
  \item Demostrar solo una inclusión y concluir indebidamente que dos conjuntos son iguales.
\end{erroresfrecuentes}

\begin{seminario}
  \item Demostración de igualdades conjuntistas mediante doble inclusión.
  \item Ejercicios de traducción entre lenguaje lógico y lenguaje conjuntista.
  \item Construcción de ejemplos finitos.
\end{seminario}

\begin{ejerciciospropuestos}
  \item Sean \(A=\{1,2,3\}\), \(B=\{2,3,4\}\) y \(C=\{3,4,5\}\). Calcula \(A\cup B\), \(A\cap B\), \(A\setminus B\), \(B\setminus A\), \(A\cap(B\cup C)\) y \((A\cap B)\cup(A\cap C)\).
  \item Demuestra, usando doble inclusión, la identidad
  \[
    A\cup(B\cap C)=(A\cup B)\cap(A\cup C).
  \]
  \item Demuestra las leyes de De Morgan:
  \[
    (A\cup B)^c=A^c\cap B^c,
    \qquad
    (A\cap B)^c=A^c\cup B^c.
  \]
  \item Calcula el conjunto de las partes de \(\{a,b\}\) y de \(\{a,b,c\}\).
  \item Sea \(A=\{1,2\}\) y \(B=\{x,y,z\}\). Escribe todos los elementos de \(A\times B\) y de \(B\times A\). ¿Son iguales estos dos productos cartesianos?
  \item Sea \((A_i)_{i\in\{1,2,3\}}\) la familia definida por
  \[
    A_1=\{1,2,3\},\qquad A_2=\{2,3,4\},\qquad A_3=\{3,4,5\}.
  \]
  Calcula \(\bigcup_{i=1}^3 A_i\) y \(\bigcap_{i=1}^3 A_i\).
\end{ejerciciospropuestos}

\resumen{El capítulo introduce el lenguaje elemental de la teoría de conjuntos. Las nociones de pertenencia, inclusión, operaciones conjuntistas, producto cartesiano, conjunto de las partes y familias de conjuntos proporcionan una base común para el resto del curso. La técnica central para demostrar igualdades entre conjuntos es la doble inclusión.}

% ------------------------------------------------------------
\chapter{Correspondencias, aplicaciones y relaciones binarias}

\begin{objetivos}
  \item Distinguir entre correspondencia, aplicación y relación binaria.
  \item Definir dominio, codominio, imagen e imagen recíproca.
  \item Estudiar inyectividad, sobreyectividad y biyectividad.
  \item Introducir relaciones de equivalencia y relaciones de orden.
\end{objetivos}

\section{Correspondencias}

En los capítulos anteriores se ha introducido el lenguaje de los conjuntos. Ahora vamos a estudiar distintas formas de relacionar los elementos de un conjunto con los elementos de otro. Esta idea aparece constantemente en matemáticas: asociar a cada número su cuadrado, asociar a cada estudiante su calificación, asociar a cada punto del plano una coordenada, o comparar dos objetos mediante una relación.

Antes de introducir símbolos, fijamos la situación general. Sean \(A\) y \(B\) dos conjuntos. Una forma de relacionar elementos de \(A\) con elementos de \(B\) consiste en indicar qué pares \((a,b)\), con \(a\in A\) y \(b\in B\), están relacionados.

\begin{definicion}[Correspondencia]
Sean \(A\) y \(B\) dos conjuntos. Una \emph{correspondencia} de \(A\) en \(B\) es una regla que asocia a algunos elementos de \(A\) uno o varios elementos de \(B\).
\end{definicion}

Una correspondencia puede verse como una manera flexible de asignar elementos. No exige que todos los elementos de \(A\) tengan asociado algún elemento de \(B\), ni exige que un elemento de \(A\) tenga asociado un único elemento de \(B\).

\begin{ejemplo}
Sea \(A=\{1,2,3\}\) y sea \(B=\{a,b,c\}\). Podemos definir una correspondencia indicando que
\[
1 \longleftrightarrow a,\qquad
1 \longleftrightarrow b,\qquad
2 \longleftrightarrow c.
\]
En esta correspondencia, el elemento \(1\) está relacionado con dos elementos de \(B\), mientras que el elemento \(3\) no está relacionado con ningún elemento de \(B\).
\end{ejemplo}

\begin{advertencia}
Una correspondencia no es necesariamente una aplicación. En una aplicación, cada elemento del conjunto de partida debe tener una única imagen. En una correspondencia general, un elemento puede no tener imagen o puede tener varias.
\end{advertencia}

Una forma rigurosa de describir una correspondencia es mediante un subconjunto del producto cartesiano \(A\times B\). Recordemos que \(A\times B\) denota el conjunto formado por todos los pares ordenados \((a,b)\), donde \(a\in A\) y \(b\in B\).

\begin{definicion}[Grafo de una correspondencia]
Sean \(A\) y \(B\) dos conjuntos. El \emph{grafo} de una correspondencia de \(A\) en \(B\) es un subconjunto \(\Gamma\subset A\times B\). Si \((a,b)\in\Gamma\), diremos que \(a\) está relacionado con \(b\) mediante la correspondencia.
\end{definicion}

\section{Aplicaciones}

Las aplicaciones, también llamadas funciones, son un caso particular especialmente importante de correspondencias. En matemáticas, una aplicación expresa una asignación bien determinada: a cada elemento del conjunto de partida se le asigna exactamente un elemento del conjunto de llegada.

\begin{definicion}[Aplicación]
Sean \(A\) y \(B\) dos conjuntos. Una \emph{aplicación} de \(A\) en \(B\) es una regla \(f\) que asigna a cada elemento \(a\in A\) un único elemento de \(B\), que se denota por \(f(a)\).

Se escribe
\[
f:A\longrightarrow B.
\]
El conjunto \(A\) se llama \emph{dominio} o conjunto de partida de \(f\), y el conjunto \(B\) se llama \emph{codominio} o conjunto de llegada de \(f\).
\end{definicion}

La notación
\[
f:A\longrightarrow B
\]
solo debe utilizarse después de haber definido qué son \(A\), \(B\) y \(f\). Si \(a\in A\), el símbolo \(f(a)\) designa el elemento de \(B\) asignado a \(a\) por la aplicación \(f\).

\begin{ejemplo}
Sea \(A=\{1,2,3\}\) y sea \(B=\{a,b,c\}\). La regla definida por
\[
f(1)=a,\qquad f(2)=a,\qquad f(3)=c
\]
es una aplicación de \(A\) en \(B\). Cada elemento de \(A\) tiene asignado exactamente un elemento de \(B\).
\end{ejemplo}

\begin{contraejemplo}
Sea \(A=\{1,2,3\}\) y sea \(B=\{a,b,c\}\). La regla
\[
1\mapsto a,\qquad 1\mapsto b,\qquad 2\mapsto c
\]
no define una aplicación de \(A\) en \(B\), porque el elemento \(1\) tendría dos imágenes y el elemento \(3\) no tendría ninguna.
\end{contraejemplo}

\begin{convencion}[Notación de imágenes]
Si \(f:A\to B\) es una aplicación y \(a\in A\), escribiremos \(f(a)\) para denotar la imagen de \(a\) por \(f\). La expresión \(f(a)\) solo tiene sentido si previamente se han definido \(f\), \(A\), \(B\) y se sabe que \(a\in A\).
\end{convencion}

\section{Dominio, codominio e imagen}

En una aplicación intervienen varios conjuntos que deben distinguirse cuidadosamente. Esta distinción evita muchos errores conceptuales.

\begin{definicion}[Dominio y codominio]
Sea \(f:A\to B\) una aplicación. El conjunto \(A\) se llama \emph{dominio} de \(f\). El conjunto \(B\) se llama \emph{codominio} de \(f\).
\end{definicion}

El dominio es el conjunto de elementos a los que se puede aplicar \(f\). El codominio es el conjunto en el que están obligadas a vivir las imágenes de los elementos del dominio.

\begin{definicion}[Imagen de un elemento]
Sea \(f:A\to B\) una aplicación y sea \(a\in A\). El elemento \(f(a)\in B\) se llama \emph{imagen de \(a\) por \(f\)}.
\end{definicion}

\begin{definicion}[Imagen de una aplicación]
Sea \(f:A\to B\) una aplicación. La \emph{imagen} de \(f\) es el subconjunto de \(B\) formado por todos los valores que toma la aplicación. Se denota por \(\operatorname{Im}(f)\) y se define por
\[
\operatorname{Im}(f)=\{b\in B:\text{ existe }a\in A\text{ tal que }f(a)=b\}.
\]
\end{definicion}

Equivalentemente, usando cuantificadores,
\[
b\in \operatorname{Im}(f) \quad \Longleftrightarrow \quad \exists a\in A \text{ tal que } f(a)=b.
\]

\begin{ejemplo}
Sea \(f:\{1,2,3\}\to \{a,b,c,d\}\) definida por
\[
f(1)=a,\qquad f(2)=a,\qquad f(3)=c.
\]
Entonces
\[
\operatorname{Im}(f)=\{a,c\}.
\]
Obsérvese que el codominio es \(\{a,b,c,d\}\), pero la imagen es solo \(\{a,c\}\).
\end{ejemplo}

\begin{advertencia}
No debe confundirse el codominio de una aplicación con su imagen. La imagen siempre es un subconjunto del codominio, pero no tiene por qué coincidir con él.
\end{advertencia}

\section{Imagen recíproca}

La imagen recíproca permite transportar subconjuntos del codominio al dominio. Es una construcción fundamental porque se define para cualquier aplicación, sin necesidad de que exista una aplicación inversa.

\begin{definicion}[Imagen recíproca]
Sea \(f:A\to B\) una aplicación y sea \(C\subset B\). La \emph{imagen recíproca} de \(C\) por \(f\) es el subconjunto de \(A\) definido por
\[
f^{-1}(C)=\{a\in A:f(a)
\in C\}.
\]
\end{definicion}

En esta definición, el símbolo \(f^{-1}(C)\) no significa necesariamente que exista una aplicación inversa \(f^{-1}:B\to A\). Es una notación para designar el conjunto de elementos de \(A\) cuya imagen pertenece a \(C\).

\begin{ejemplo}
Sea \(f:\{1,2,3,4\}\to \{a,b,c\}\) definida por
\[
f(1)=a,\qquad f(2)=b,\qquad f(3)=b,\qquad f(4)=c.
\]
Si \(C=\{b,c\}\), entonces
\[
f^{-1}(C)=\{2,3,4\}.
\]
\end{ejemplo}

\begin{advertencia}
La imagen recíproca de un conjunto está siempre contenida en el dominio. Si \(f:A\to B\) y \(C\subset B\), entonces \(f^{-1}(C)\subset A\).
\end{advertencia}

\section{Inyectividad, sobreyectividad y biyectividad}

Las aplicaciones se clasifican según cómo distribuyen los elementos del dominio dentro del codominio.

\begin{definicion}[Aplicación inyectiva]
Sea \(f:A\to B\) una aplicación. Diremos que \(f\) es \emph{inyectiva} si elementos distintos de \(A\) tienen imágenes distintas en \(B\). Es decir,
\[
\forall a_1,a_2\in A,\quad f(a_1)=f(a_2)\Rightarrow a_1=a_2.
\]
\end{definicion}

La definición anterior se puede leer así: si dos elementos tienen la misma imagen, entonces en realidad eran el mismo elemento.

\begin{ejemplo}
La aplicación \(f:\{1,2,3\}\to \{a,b,c,d\}\) dada por
\[
f(1)=a,\qquad f(2)=b,\qquad f(3)=c
\]
es inyectiva.
\end{ejemplo}

\begin{contraejemplo}
La aplicación \(g:\{1,2,3\}\to \{a,b,c\}\) dada por
\[
g(1)=a,\qquad g(2)=a,\qquad g(3)=c
\]
no es inyectiva, porque \(g(1)=g(2)\) y, sin embargo, \(1\ne 2\).
\end{contraejemplo}

\begin{definicion}[Aplicación sobreyectiva]
Sea \(f:A\to B\) una aplicación. Diremos que \(f\) es \emph{sobreyectiva} si todo elemento del codominio es imagen de algún elemento del dominio. Es decir,
\[
\forall b\in B,\quad \exists a\in A\text{ tal que }f(a)=b.
\]
Equivalentemente, \(f\) es sobreyectiva si \(\operatorname{Im}(f)=B\).
\end{definicion}

\begin{ejemplo}
La aplicación \(f:\{1,2,3\}\to \{a,b\}\) definida por
\[
f(1)=a,\qquad f(2)=b,\qquad f(3)=b
\]
es sobreyectiva, porque tanto \(a\) como \(b\) son imágenes de algún elemento del dominio.
\end{ejemplo}

\begin{definicion}[Aplicación biyectiva]
Sea \(f:A\to B\) una aplicación. Diremos que \(f\) es \emph{biyectiva} si es inyectiva y sobreyectiva.
\end{definicion}

Una aplicación biyectiva establece una correspondencia perfecta entre los elementos de \(A\) y los elementos de \(B\): cada elemento de \(A\) tiene una única imagen en \(B\), y cada elemento de \(B\) procede de un único elemento de \(A\).

\begin{ejemplo}
La aplicación \(f:\{1,2,3\}\to \{a,b,c\}\) definida por
\[
f(1)=a,\qquad f(2)=b,\qquad f(3)=c
\]
es biyectiva.
\end{ejemplo}

\section{Composición de aplicaciones}

La composición permite aplicar sucesivamente dos aplicaciones. Para definirla correctamente hay que comprobar que el codominio de la primera aplicación encaja con el dominio de la segunda.

\begin{definicion}[Composición]
Sean \(A\), \(B\) y \(C\) tres conjuntos. Sean \(f:A\to B\) y \(g:B\to C\) dos aplicaciones. La \emph{composición} de \(g\) con \(f\) es la aplicación
\[
g\circ f:A\to C
\]
definida por
\[
(g\circ f)(a)=g(f(a)),\qquad a\in A.
\]
\end{definicion}

El símbolo \(g\circ f\) se lee ``\(g\) compuesta con \(f\)''. El orden es importante: primero se aplica \(f\) y después se aplica \(g\).

\begin{ejemplo}
Sean \(f:\mathbb{R}\to\mathbb{R}\) y \(g:\mathbb{R}\to\mathbb{R}\) definidas por
\[
f(x)=x+1,\qquad g(x)=x^2.
\]
Entonces
\[
(g\circ f)(x)=g(f(x))=g(x+1)=(x+1)^2,
\]
mientras que
\[
(f\circ g)(x)=f(g(x))=f(x^2)=x^2+1.
\]
Por tanto, en general, \(g\circ f\ne f\circ g\).
\end{ejemplo}

\begin{advertencia}
La composición de aplicaciones no es conmutativa en general. Además, \(g\circ f\) solo está definida si las imágenes de \(f\) pertenecen al dominio de \(g\).
\end{advertencia}

\section{Aplicación identidad y aplicación inversa}

La aplicación identidad es la aplicación que no modifica ningún elemento.

\begin{definicion}[Aplicación identidad]
Sea \(A\) un conjunto. La \emph{aplicación identidad} de \(A\) es la aplicación
\[
\operatorname{id}_A:A\to A
\]
definida por
\[
\operatorname{id}_A(a)=a,\qquad a\in A.
\]
\end{definicion}

La identidad actúa como elemento neutro para la composición. Si \(f:A\to B\) es una aplicación, entonces
\[
f\circ \operatorname{id}_A=f,
\qquad
\operatorname{id}_B\circ f=f.
\]

\begin{definicion}[Aplicación inversa]
Sea \(f:A\to B\) una aplicación biyectiva. La \emph{aplicación inversa} de \(f\) es la aplicación
\[
f^{-1}:B\to A
\]
que asigna a cada \(b\in B\) el único elemento \(a\in A\) tal que \(f(a)=b\).
\end{definicion}

La inversa de una aplicación solo existe como aplicación cuando la aplicación original es biyectiva.

\begin{proposicion}
Sea \(f:A\to B\) una aplicación biyectiva. Entonces existe una aplicación \(f^{-1}:B\to A\) tal que
\[
f^{-1}\circ f=\operatorname{id}_A,
\qquad
f\circ f^{-1}=\operatorname{id}_B.
\]
\end{proposicion}

\begin{proof}
Como \(f\) es sobreyectiva, para todo \(b\in B\) existe al menos un \(a\in A\) tal que \(f(a)=b\). Como \(f\) es inyectiva, dicho elemento \(a\) es único. Por tanto, podemos definir \(f^{-1}(b)=a\). Esta regla define una aplicación de \(B\) en \(A\). Las igualdades con las identidades se deducen directamente de la definición.
\end{proof}

\begin{advertencia}
La notación \(f^{-1}(C)\), para la imagen recíproca de un subconjunto \(C\subset B\), puede utilizarse aunque \(f\) no sea biyectiva. En cambio, la notación \(f^{-1}:B\to A\), como aplicación inversa, solo tiene sentido cuando \(f\) es biyectiva.
\end{advertencia}

\section{Relaciones binarias}

Una relación binaria es una forma de comparar o relacionar pares de elementos de un mismo conjunto o de dos conjuntos. En este curso nos centraremos sobre todo en relaciones binarias sobre un conjunto.

\begin{definicion}[Relación binaria]
Sea \(A\) un conjunto. Una \emph{relación binaria} sobre \(A\) es un subconjunto \(R\subset A\times A\).
Si \((a,b)\in R\), diremos que \(a\) está relacionado con \(b\) y escribiremos
\[
a\,R\,b.
\]
\end{definicion}

El símbolo \(aRb\) solo es una abreviatura de \((a,b)\in R\). Por tanto, antes de escribir \(aRb\), deben estar definidos el conjunto \(A\), los elementos \(a,b\in A\) y la relación \(R\subset A\times A\).

\begin{ejemplo}
En el conjunto \(\mathbb{Z}\) de los números enteros, la relación \(\leq\) definida por el orden usual es una relación binaria. Si \(m,n\in\mathbb{Z}\), la expresión \(m\leq n\) significa que el entero \(m\) es menor o igual que el entero \(n\).
\end{ejemplo}

\begin{definicion}[Propiedades de una relación binaria]
Sea \(R\) una relación binaria sobre un conjunto \(A\).
\begin{enumerate}
  \item \(R\) es \emph{reflexiva} si para todo \(a\in A\), se cumple \(aRa\).
  \item \(R\) es \emph{simétrica} si para todo \(a,b\in A\), de \(aRb\) se deduce \(bRa\).
  \item \(R\) es \emph{antisimétrica} si para todo \(a,b\in A\), de \(aRb\) y \(bRa\) se deduce \(a=b\).
  \item \(R\) es \emph{transitiva} si para todo \(a,b,c\in A\), de \(aRb\) y \(bRc\) se deduce \(aRc\).
\end{enumerate}
\end{definicion}

\begin{ejemplo}
En \(\mathbb{Z}\), la relación \(\leq\) es reflexiva, antisimétrica y transitiva. No es simétrica, porque por ejemplo \(2\leq 3\), pero no se cumple \(3\leq 2\).
\end{ejemplo}

\section{Relaciones de equivalencia}

Las relaciones de equivalencia permiten agrupar elementos que se consideran iguales respecto de cierto criterio.

\begin{definicion}[Relación de equivalencia]
Sea \(A\) un conjunto. Una relación binaria \(\sim\) sobre \(A\) se llama \emph{relación de equivalencia} si es reflexiva, simétrica y transitiva.
\end{definicion}

Cuando \(a\sim b\), decimos que \(a\) y \(b\) son equivalentes respecto de la relación \(\sim\).

\begin{ejemplo}
Sea \(A\) el conjunto de todos los estudiantes de una clase. Definimos una relación \(\sim\) sobre \(A\) diciendo que, para \(x,y\in A\),
\[
x\sim y \quad \Longleftrightarrow \quad x\text{ e }y\text{ tienen el mismo grupo de seminario}.
\]
Esta relación es de equivalencia: cada estudiante está en el mismo grupo que sí mismo, si \(x\) está en el mismo grupo que \(y\), entonces \(y\) está en el mismo grupo que \(x\), y si \(x\) está en el mismo grupo que \(y\) y \(y\) está en el mismo grupo que \(z\), entonces \(x\) está en el mismo grupo que \(z\).
\end{ejemplo}

\begin{ejemplo}
En \(\mathbb{Z}\), fijado un entero \(n\geq 2\), se define la relación de congruencia módulo \(n\) por
\[
a\equiv b\pmod n
\quad \Longleftrightarrow \quad
n \text{ divide a } a-b.
\]
Esta relación es una relación de equivalencia sobre \(\mathbb{Z}\).
\end{ejemplo}

\section{Clases de equivalencia}

Una relación de equivalencia divide un conjunto en subconjuntos disjuntos llamados clases de equivalencia.

\begin{definicion}[Clase de equivalencia]
Sea \(\sim\) una relación de equivalencia sobre un conjunto \(A\). Si \(a\in A\), la \emph{clase de equivalencia} de \(a\) es el conjunto
\[
[a]=\{x\in A:x\sim a\}.
\]
\end{definicion}

El símbolo \([a]\) depende de la relación de equivalencia que se esté considerando. Por tanto, no debe utilizarse sin haber definido previamente la relación \(\sim\).

\begin{ejemplo}
Consideremos en \(\mathbb{Z}\) la congruencia módulo \(3\). La clase de equivalencia de \(0\) es
\[
[0]=\{z\in\mathbb{Z}:z\equiv 0\pmod 3\},
\]
es decir, el conjunto de los múltiplos de \(3\). Análogamente, existen tres clases:
\[
[0],\qquad [1],\qquad [2].
\]
Todo entero pertenece a una y solo una de estas clases.
\end{ejemplo}

\begin{proposicion}
Sea \(\sim\) una relación de equivalencia sobre un conjunto \(A\). Si \(a,b\in A\), entonces se cumple una y solo una de las dos alternativas siguientes:
\begin{enumerate}
  \item \([a]=[b]\),
  \item \([a]\cap[b]=\varnothing\).
\end{enumerate}
\end{proposicion}

\begin{proof}
Supongamos que \([a]\cap[b]\ne\varnothing\). Entonces existe \(x\in A\) tal que \(x\in[a]\) y \(x\in[b]\). Por definición de clase de equivalencia, \(x\sim a\) y \(x\sim b\). Como \(\sim\) es simétrica, de \(x\sim a\) se deduce \(a\sim x\). Como \(a\sim x\) y \(x\sim b\), por transitividad se obtiene \(a\sim b\).

Veamos que \([a]=[b]\). Sea \(y\in[a]\). Entonces \(y\sim a\). Como \(a\sim b\), por transitividad se tiene \(y\sim b\). Por tanto, \(y\in[b]\). Hemos probado \([a]\subset[b]\). De forma análoga se prueba \([b]\subset[a]\). Luego \([a]=[b]\).

Por tanto, si dos clases no son disjuntas, son iguales. En consecuencia, dos clases cualesquiera son iguales o disjuntas.
\end{proof}

\section{Relaciones de orden}

Las relaciones de orden permiten comparar elementos. La comparación puede ser total, cuando cualquier par de elementos es comparable, o parcial, cuando algunos elementos no pueden compararse.

\begin{definicion}[Relación de orden]
Sea \(A\) un conjunto. Una relación binaria \(\preceq\) sobre \(A\) se llama \emph{relación de orden parcial} si es reflexiva, antisimétrica y transitiva.
\end{definicion}

\begin{definicion}[Orden total]
Sea \(\preceq\) una relación de orden parcial sobre un conjunto \(A\). Diremos que \(\preceq\) es un \emph{orden total} si para todo \(a,b\in A\), se cumple al menos una de las dos relaciones
\[
a\preceq b
\qquad\text{o}\qquad
b\preceq a.
\]
\end{definicion}

\begin{ejemplo}
El orden usual \(\leq\) sobre \(\mathbb{Z}\) es un orden total. Dados dos enteros \(m,n\in\mathbb{Z}\), siempre se cumple \(m\leq n\) o \(n\leq m\).
\end{ejemplo}

\begin{ejemplo}
Sea \(X\) un conjunto. En el conjunto de las partes \(\mathcal{P}(X)\), la inclusión \(\subset\) define un orden parcial. En general no es un orden total, porque pueden existir subconjuntos \(A,B\subset X\) tales que ni \(A\subset B\) ni \(B\subset A\).
\end{ejemplo}

\begin{advertencia}
No debe confundirse una relación de equivalencia con una relación de orden. Una relación de equivalencia agrupa elementos; una relación de orden compara elementos. Por eso una relación de equivalencia exige simetría, mientras que una relación de orden parcial exige antisimetría.
\end{advertencia}

\begin{erroresfrecuentes}
  \item Confundir una correspondencia con una aplicación.
  \item Usar la notación \(f(a)\) sin haber definido previamente la aplicación \(f\) y su dominio.
  \item Confundir codominio e imagen de una aplicación.
  \item Pensar que toda aplicación tiene inversa.
  \item Confundir la imagen recíproca de un subconjunto con la aplicación inversa.
  \item Confundir simetría con antisimetría en una relación binaria.
  \item Usar la notación \([a]\) sin especificar la relación de equivalencia correspondiente.
\end{erroresfrecuentes}

\begin{seminario}
  \item Clasificación de aplicaciones entre conjuntos finitos.
  \item Verificación de las propiedades de una relación binaria.
  \item Cálculo de clases de equivalencia.
\end{seminario}

\begin{ejerciciospropuestos}
  \item Sea \(A=\{1,2,3\}\) y \(B=\{a,b\}\). Escribe todas las aplicaciones posibles de \(A\) en \(B\). Decide cuáles son sobreyectivas.
  \item Sea \(A=\{1,2,3\}\). Considera la relación \(R\subset A\times A\) dada por
  \[
  R=\{(1,1),(2,2),(3,3),(1,2),(2,1)\}.
  \]
  Decide si \(R\) es reflexiva, simétrica, antisimétrica o transitiva.
  \item Construye una relación de equivalencia no trivial sobre el conjunto \(A=\{1,2,3,4\}\) y calcula sus clases de equivalencia.
  \item En \(\mathbb{Z}\), considera la relación \(a\equiv b\pmod 4\). Calcula las clases \([0]\), \([1]\), \([2]\) y \([3]\).
  \item Sea \(X=\{a,b,c\}\). En \(\mathcal{P}(X)\), decide si los subconjuntos \(\{a\}\) y \(\{b,c\}\) son comparables para la relación de inclusión.
\end{ejerciciospropuestos}

\resumen{El capítulo introduce las nociones de correspondencia, aplicación y relación binaria. Las aplicaciones son correspondencias con una asignación única para cada elemento del dominio. Las relaciones binarias permiten formalizar comparaciones y agrupaciones; entre ellas destacan las relaciones de equivalencia, que producen clases de equivalencia, y las relaciones de orden, que permiten comparar elementos de un conjunto.}

% ------------------------------------------------------------
\chapter{Estructuras algebraicas elementales}

\begin{objetivos}
  \item Introducir la noción de operación interna.
  \item Estudiar las propiedades básicas de las operaciones.
  \item Presentar semigrupos, monoides, grupos, anillos y cuerpos de forma introductoria.
  \item Reconocer estructuras algebraicas en ejemplos elementales.
\end{objetivos}

\section{Motivación}

En los capítulos anteriores se han introducido conjuntos, aplicaciones y relaciones binarias. En este capítulo se añade un nuevo ingrediente: las \emph{operaciones}. Una operación permite combinar elementos de un conjunto para obtener nuevos elementos. Por ejemplo, al sumar dos números naturales se obtiene otro número natural; al componer dos simetrías de una figura se obtiene otra simetría de la misma figura.

La idea fundamental de este capítulo es que muchos objetos matemáticos no se estudian únicamente como conjuntos, sino como conjuntos dotados de operaciones que satisfacen ciertas propiedades. A estos objetos se les llama \emph{estructuras algebraicas}.

\begin{convencion}[Símbolos y operaciones]
Antes de escribir una expresión como \(a*b\), se debe haber definido previamente qué conjunto contiene a los elementos \(a\) y \(b\), qué significa el símbolo \(*\), y en qué conjunto se encuentra el resultado de la operación.
\end{convencion}

\section{Operaciones internas}

Sea \(A\) un conjunto. Una \emph{operación interna} en \(A\) es una aplicación que toma dos elementos de \(A\) y devuelve de nuevo un elemento de \(A\).

\begin{definicion}[Operación interna]
Sea \(A\) un conjunto. Una operación interna en \(A\) es una aplicación
\[
  *:A\times A\longrightarrow A.
\]
Si \(a\) y \(b\) son elementos de \(A\), se escribe
\[
  a*b
\]
para denotar el elemento de \(A\) que resulta de aplicar la operación \(*\) al par ordenado \((a,b)\).
\end{definicion}

La condición de que el resultado pertenezca de nuevo a \(A\) es esencial. Si al operar dos elementos de \(A\) se pudiera obtener un objeto fuera de \(A\), no estaríamos ante una operación interna en \(A\).

\begin{ejemplo}[Suma de números naturales]
Sea \(\mathbb{N}\) el conjunto de los números naturales. La suma usual define una operación interna
\[
  +:\mathbb{N}\times \mathbb{N}\longrightarrow \mathbb{N},
\]
porque si \(m\) y \(n\) son números naturales, entonces \(m+n\) también es un número natural.
\end{ejemplo}

\begin{ejemplo}[Resta en los números naturales]
La resta usual no define una operación interna en \(\mathbb{N}\) si se toma, por ejemplo, \(2-5\), ya que el resultado no es un número natural.
\end{ejemplo}

\begin{advertencia}
Una misma regla puede definir una operación interna en un conjunto y no definirla en otro. Por ejemplo, la resta no es una operación interna en \(\mathbb{N}\), pero sí lo es en el conjunto \(\mathbb{Z}\) de los números enteros.
\end{advertencia}

\section{Asociatividad y conmutatividad}

Una vez definida una operación interna, se estudian sus propiedades. Las dos primeras propiedades fundamentales son la asociatividad y la conmutatividad.

\begin{definicion}[Asociatividad]
Sea \(A\) un conjunto y sea \(*:A\times A\to A\) una operación interna. Se dice que \(*\) es \emph{asociativa} si, para cualesquiera \(a,b,c\in A\), se cumple
\[
  (a*b)*c=a*(b*c).
\]
\end{definicion}

La asociatividad permite escribir expresiones como \(a*b*c\) sin ambigüedad, porque no importa si primero se opera \(a*b\) o si primero se opera \(b*c\).

\begin{definicion}[Conmutatividad]
Sea \(A\) un conjunto y sea \(*:A\times A\to A\) una operación interna. Se dice que \(*\) es \emph{conmutativa} si, para cualesquiera \(a,b\in A\), se cumple
\[
  a*b=b*a.
\]
\end{definicion}

\begin{ejemplo}
La suma usual de números enteros es asociativa y conmutativa. Es decir, para cualesquiera \(a,b,c\in\mathbb{Z}\), se tiene
\[
  (a+b)+c=a+(b+c)
\]
y
\[
  a+b=b+a.
\]
\end{ejemplo}

\begin{ejemplo}
La composición de aplicaciones no es conmutativa en general. Si \(f\) y \(g\) son aplicaciones de un conjunto en sí mismo, normalmente
\[
  f\circ g\neq g\circ f.
\]
Por tanto, el orden en el que se componen las aplicaciones puede cambiar el resultado.
\end{ejemplo}

\section{Elemento neutro}

Algunas operaciones tienen un elemento especial que no modifica a los demás elementos cuando se opera con ellos.

\begin{definicion}[Elemento neutro]
Sea \(A\) un conjunto y sea \(*:A\times A\to A\) una operación interna. Un elemento \(e\in A\) se llama \emph{elemento neutro} para \(*\) si, para todo \(a\in A\), se cumple
\[
  e*a=a
  \qquad\text{y}\qquad
  a*e=a.
\]
\end{definicion}

\begin{ejemplo}
En \(\mathbb{Z}\), el número \(0\) es el elemento neutro de la suma, porque para todo \(a\in\mathbb{Z}\),
\[
  0+a=a
  \qquad\text{y}\qquad
  a+0=a.
\]
\end{ejemplo}

\begin{ejemplo}
En \(\mathbb{Z}\), el número \(1\) es el elemento neutro del producto, porque para todo \(a\in\mathbb{Z}\),
\[
  1\cdot a=a
  \qquad\text{y}\qquad
  a\cdot 1=a.
\]
\end{ejemplo}

\begin{proposicion}[Unicidad del elemento neutro]
Sea \(A\) un conjunto y sea \(*\) una operación interna en \(A\). Si \(*\) tiene elemento neutro, entonces dicho elemento neutro es único.
\end{proposicion}

\begin{proof}
Supongamos que \(e\in A\) y \(e'\in A\) son dos elementos neutros para la operación \(*\). Como \(e\) es neutro, se tiene
\[
  e*e'=e'.
\]
Como \(e'\) también es neutro, se tiene
\[
  e*e'=e.
\]
Por tanto, ambos elementos son iguales:
\[
  e=e'.
\]
Así, el elemento neutro, si existe, es único.
\end{proof}

\section{Elemento inverso}

Cuando una operación tiene elemento neutro, se puede preguntar si un elemento dado puede deshacerse mediante otro elemento.

\begin{definicion}[Elemento inverso]
Sea \(A\) un conjunto, sea \(*:A\times A\to A\) una operación interna y sea \(e\in A\) un elemento neutro para \(*\). Dado \(a\in A\), se dice que un elemento \(b\in A\) es un \emph{inverso} de \(a\) si
\[
  a*b=e
  \qquad\text{y}\qquad
  b*a=e.
\]
\end{definicion}

\begin{ejemplo}
En \(\mathbb{Z}\), con la suma usual, el elemento neutro es \(0\). El inverso de \(a\in\mathbb{Z}\) es \(-a\), porque
\[
  a+(-a)=0
  \qquad\text{y}\qquad
  (-a)+a=0.
\]
\end{ejemplo}

\begin{ejemplo}
En \(\mathbb{Z}\), con el producto usual, el elemento neutro es \(1\). El número \(2\) no tiene inverso multiplicativo en \(\mathbb{Z}\), porque no existe ningún entero \(b\) tal que \(2b=1\).
\end{ejemplo}

\begin{advertencia}
La palabra ``inverso'' depende de la operación. El inverso aditivo de \(a\) es \(-a\), mientras que el inverso multiplicativo de \(a\), cuando existe, es un elemento \(b\) tal que \(ab=1\).
\end{advertencia}

\section{Distributividad}

Algunas estructuras algebraicas tienen dos operaciones. En ese caso, además de estudiar cada operación por separado, interesa estudiar cómo interactúan entre sí.

\begin{definicion}[Distributividad]
Sea \(A\) un conjunto con dos operaciones internas, que denotamos por \(+\) y \(\cdot\). Se dice que el producto \(\cdot\) es \emph{distributivo} respecto de la suma \(+\) si, para cualesquiera \(a,b,c\in A\), se cumplen
\[
  a\cdot (b+c)=a\cdot b+a\cdot c
\]
y
\[
  (a+b)\cdot c=a\cdot c+b\cdot c.
\]
\end{definicion}

La primera igualdad se llama distributividad por la izquierda, y la segunda distributividad por la derecha.

\begin{ejemplo}
En \(\mathbb{Z}\), el producto usual es distributivo respecto de la suma usual. Por ejemplo,
\[
  2\cdot(3+5)=2\cdot 3+2\cdot 5.
\]
Ambos miembros valen \(16\).
\end{ejemplo}

\section{Semigrupos y monoides}

Las primeras estructuras algebraicas se obtienen considerando una sola operación interna.

\begin{definicion}[Semigrupo]
Un \emph{semigrupo} es un par \((A,*)\), donde \(A\) es un conjunto y \(*:A\times A\to A\) es una operación interna asociativa.
\end{definicion}

\begin{ejemplo}
El par \((\mathbb{N},+)\), donde \(+\) es la suma usual, es un semigrupo, porque la suma de números naturales es una operación interna y asociativa.
\end{ejemplo}

\begin{definicion}[Monoide]
Un \emph{monoide} es un triple \((A,*,e)\), donde \(A\) es un conjunto, \(*:A\times A\to A\) es una operación interna asociativa y \(e\in A\) es un elemento neutro para \(*\).
\end{definicion}

A veces se escribe simplemente \((A,*)\) para denotar un monoide, dejando implícito el elemento neutro cuando no hay riesgo de confusión.

\begin{ejemplo}
El par \((\mathbb{Z},+)\) es un monoide. El elemento neutro es \(0\).
\end{ejemplo}

\begin{ejemplo}
Sea \(X\) un conjunto. El conjunto de todas las aplicaciones de \(X\) en \(X\), con la composición de aplicaciones, es un monoide. El elemento neutro es la aplicación identidad \(\operatorname{id}_X:X\to X\), definida por \(\operatorname{id}_X(x)=x\) para todo \(x\in X\).
\end{ejemplo}

\section{Grupos}

Un grupo es una estructura algebraica con una operación interna asociativa, elemento neutro e inversos para todos sus elementos.

\begin{definicion}[Grupo]
Un \emph{grupo} es una terna \((G,*,e)\), donde \(G\) es un conjunto, \(*:G\times G\to G\) es una operación interna y \(e\in G\), tal que se cumplen las siguientes propiedades:
\begin{enumerate}
  \item \textbf{Asociatividad:} para cualesquiera \(a,b,c\in G\),
  \[
    (a*b)*c=a*(b*c).
  \]
  \item \textbf{Elemento neutro:} para todo \(a\in G\),
  \[
    e*a=a
    \qquad\text{y}\qquad
    a*e=a.
  \]
  \item \textbf{Elemento inverso:} para todo \(a\in G\), existe \(b\in G\) tal que
  \[
    a*b=e
    \qquad\text{y}\qquad
    b*a=e.
  \]
\end{enumerate}
\end{definicion}

\begin{definicion}[Grupo abeliano]
Un grupo \((G,*,e)\) se llama \emph{abeliano} o \emph{conmutativo} si, además, para cualesquiera \(a,b\in G\), se cumple
\[
  a*b=b*a.
\]
\end{definicion}

\begin{ejemplo}
El conjunto \(\mathbb{Z}\), con la suma usual, es un grupo abeliano. El elemento neutro es \(0\), y el inverso de \(a\in\mathbb{Z}\) es \(-a\).
\end{ejemplo}

\begin{contraejemplo}
El conjunto \(\mathbb{N}\), con la suma usual, no es un grupo si \(\mathbb{N}\) contiene al \(0\), porque el número \(1\) no tiene inverso aditivo dentro de \(\mathbb{N}\). En efecto, no existe \(n\in\mathbb{N}\) tal que \(1+n=0\).
\end{contraejemplo}

\section{Anillos}

Un anillo es una estructura con dos operaciones, que suelen llamarse suma y producto. La suma proporciona una estructura de grupo abeliano, mientras que el producto es asociativo y distributivo respecto de la suma.

\begin{definicion}[Anillo]
Un \emph{anillo} es una estructura \((A,+,\cdot,0)\), donde \(A\) es un conjunto, \(+\) y \(\cdot\) son operaciones internas en \(A\), y \(0\in A\), tal que:
\begin{enumerate}
  \item \((A,+,0)\) es un grupo abeliano.
  \item El producto \(\cdot\) es asociativo.
  \item El producto \(\cdot\) es distributivo respecto de la suma \(+\).
\end{enumerate}
\end{definicion}

En algunos textos se exige que un anillo tenga elemento neutro para el producto. En estas notas, cuando se necesite dicha propiedad, se dirá explícitamente que se trata de un anillo con unidad.

\begin{definicion}[Anillo con unidad]
Un anillo \((A,+,\cdot,0)\) se llama \emph{anillo con unidad} si existe un elemento \(1\in A\) tal que, para todo \(a\in A\),
\[
  1\cdot a=a
  \qquad\text{y}\qquad
  a\cdot 1=a.
\]
El elemento \(1\) se llama unidad o elemento neutro multiplicativo.
\end{definicion}

\begin{ejemplo}
El conjunto \(\mathbb{Z}\), con la suma y el producto usuales, es un anillo con unidad. La unidad es el número \(1\).
\end{ejemplo}

\begin{ejemplo}
El conjunto de los polinomios con coeficientes reales, denotado por \(\mathbb{R}[x]\), con la suma y el producto usuales de polinomios, es un anillo con unidad.
\end{ejemplo}

\section{Cuerpos}

Un cuerpo es una estructura algebraica en la que se puede sumar, restar, multiplicar y dividir por cualquier elemento no nulo.

\begin{definicion}[Cuerpo]
Un \emph{cuerpo} es una estructura \((K,+,\cdot,0,1)\), donde \(K\) es un conjunto, \(+\) y \(\cdot\) son operaciones internas en \(K\), y \(0,1\in K\) son elementos distintos, tal que:
\begin{enumerate}
  \item \((K,+,0)\) es un grupo abeliano.
  \item \((K\setminus\{0\},\cdot,1)\) es un grupo abeliano.
  \item El producto \(\cdot\) es distributivo respecto de la suma \(+\).
\end{enumerate}
\end{definicion}

La condición \((K\setminus\{0\},\cdot,1)\) significa que todo elemento no nulo de \(K\) tiene inverso multiplicativo.

\begin{ejemplo}
El conjunto \(\mathbb{Q}\) de los números racionales, con la suma y el producto usuales, es un cuerpo.
\end{ejemplo}

\begin{ejemplo}
El conjunto \(\mathbb{R}\) de los números reales y el conjunto \(\mathbb{C}\) de los números complejos son cuerpos con sus operaciones usuales.
\end{ejemplo}

\begin{contraejemplo}
El conjunto \(\mathbb{Z}\) de los números enteros no es un cuerpo, porque el entero \(2\) no tiene inverso multiplicativo en \(\mathbb{Z}\).
\end{contraejemplo}

\section{Ejemplos fundamentales}

Esta sección resume algunos ejemplos que aparecerán de forma recurrente en el curso.

\begin{ejemplo}[Números enteros]
El conjunto \(\mathbb{Z}\), con la suma usual, es un grupo abeliano. Con la suma y el producto usuales, es un anillo con unidad, pero no es un cuerpo.
\end{ejemplo}

\begin{ejemplo}[Números racionales]
El conjunto \(\mathbb{Q}\), con la suma y el producto usuales, es un cuerpo. Lo mismo ocurre con \(\mathbb{R}\) y \(\mathbb{C}\).
\end{ejemplo}

\begin{ejemplo}[Matrices]
Sea \(n\) un número natural positivo. El conjunto de las matrices cuadradas reales de tamaño \(n\times n\), con la suma y el producto usuales de matrices, es un anillo con unidad. Si \(n\geq 2\), el producto de matrices no es conmutativo en general.
\end{ejemplo}

\begin{ejemplo}[Simetrías]
El conjunto de las simetrías de una figura geométrica, con la composición de aplicaciones, forma un grupo. El elemento neutro es la simetría que deja todos los puntos fijos, y el inverso de una simetría es la transformación que deshace su acción.
\end{ejemplo}

\begin{erroresfrecuentes}
  \item Confundir operación interna con una regla que a veces produce elementos fuera del conjunto.
  \item Afirmar que una operación es asociativa comprobando solo algunos ejemplos.
  \item Confundir el elemento neutro aditivo con el elemento neutro multiplicativo.
  \item Hablar de inversos sin indicar respecto de qué operación se consideran.
  \item Suponer que todo anillo es un cuerpo.
  \item Suponer que todo producto algebraico es conmutativo.
\end{erroresfrecuentes}

\begin{seminario}
  \item Verificación de axiomas algebraicos.
  \item Construcción de tablas de operación.
  \item Identificación de neutros e inversos.
\end{seminario}

\begin{ejerciciospropuestos}
  \item Sea \(A=\{0,1\}\) y sea \(*\) la operación definida por la tabla
  \[
  \begin{array}{c|cc}
    * & 0 & 1 \\
    \hline
    0 & 0 & 1 \\
    1 & 1 & 0
  \end{array}
  \]
  Determina si \(*\) es conmutativa y si tiene elemento neutro.

  \item Determina si la operación \(*\) definida en \(\mathbb{Z}\) por
  \[
    a*b=a+b+1
  \]
  es asociativa. Encuentra, si existe, su elemento neutro.

  \item Sea \(A=\{a,b,c\}\). Construye una tabla de operación para una operación interna en \(A\) que sea conmutativa.

  \item Da un ejemplo de una operación interna en un conjunto finito que no sea conmutativa.

  \item Explica por qué \((\mathbb{Q},+)\) es un grupo abeliano.

  \item Explica por qué \((\mathbb{Q}\setminus\{0\},\cdot)\) es un grupo abeliano.

  \item Justifica por qué \(\mathbb{Z}\) es un anillo con unidad pero no es un cuerpo.

  \item Busca un ejemplo de dos matrices reales \(2\times 2\), llamadas \(A\) y \(B\), tales que
  \[
    AB\neq BA.
  \]
\end{ejerciciospropuestos}

\resumen{El capítulo introduce las primeras estructuras algebraicas a partir de la noción de operación interna. Las propiedades de asociatividad, conmutatividad, existencia de elemento neutro, existencia de inversos y distributividad permiten definir semigrupos, monoides, grupos, anillos y cuerpos. Estas estructuras proporcionan un lenguaje común para estudiar sistemas numéricos, matrices, aplicaciones y simetrías.}

% ------------------------------------------------------------
\chapter{Los números naturales y el principio de inducción}

\begin{objetivos}
  \item Introducir el conjunto de los números naturales.
  \item Estudiar el principio de inducción matemática.
  \item Presentar la inducción fuerte.
  \item Trabajar con definiciones recursivas y sucesiones.
\end{objetivos}

\section{El conjunto de los números naturales}

Los números naturales constituyen el primer sistema numérico que aparece de forma sistemática en matemáticas. Se utilizan para contar objetos, ordenar posiciones y formular procesos que avanzan paso a paso.

En este curso denotaremos por \(\mathbb{N}\) el conjunto de los números naturales. Adoptaremos la convención
\[
  \mathbb{N}=\{0,1,2,3,\ldots\}.
\]
Es decir, el número \(0\) pertenece a \(\mathbb{N}\). Cuando se quiera excluir el cero, escribiremos
\[
  \mathbb{N}^{\ast}=\mathbb{N}\setminus\{0\}=\{1,2,3,\ldots\}.
\]

\begin{convencion}[Uso de \(\mathbb{N}\)]
En este manual, salvo indicación explícita en sentido contrario, el símbolo \(\mathbb{N}\) designa el conjunto de los números naturales incluyendo el cero. Por tanto, antes de usar una propiedad que dependa de si \(0\) pertenece o no a \(\mathbb{N}\), se deberá recordar esta convención.
\end{convencion}

Los números naturales pueden pensarse como generados a partir de un primer número, el \(0\), y de una operación de paso al siguiente. Si \(n\in\mathbb{N}\), su sucesor se denota por \(n+1\). Así, el sucesor de \(0\) es \(1\), el sucesor de \(1\) es \(2\), y así sucesivamente.

\begin{definicion}[Sucesor]
Sea \(n\in\mathbb{N}\). Llamamos \emph{sucesor} de \(n\) al número natural \(n+1\).
\end{definicion}

Esta idea de sucesor es fundamental porque permite formular propiedades que se transmiten de un número natural al siguiente. Esa transmisión es la base del principio de inducción matemática.

\section{Orden en los números naturales}

Los números naturales están dotados de un orden natural. Si \(m,n\in\mathbb{N}\), escribimos
\[
  m\leq n
\]
para indicar que \(m\) es menor o igual que \(n\). También escribimos \(m<n\) cuando \(m\leq n\) y \(m\neq n\).

\begin{definicion}[Orden en \(\mathbb{N}\)]
Sean \(m,n\in\mathbb{N}\). La expresión \(m\leq n\) significa que \(m\) no está situado después de \(n\) en la sucesión ordenada
\[
  0,1,2,3,\ldots
\]
de los números naturales.
\end{definicion}

El orden en \(\mathbb{N}\) posee propiedades básicas que se utilizarán constantemente.

\begin{proposicion}[Propiedades elementales del orden]
Sean \(a,b,c\in\mathbb{N}\). Entonces:
\begin{enumerate}
  \item \(a\leq a\).
  \item Si \(a\leq b\) y \(b\leq a\), entonces \(a=b\).
  \item Si \(a\leq b\) y \(b\leq c\), entonces \(a\leq c\).
  \item Dados \(a,b\in\mathbb{N}\), se cumple \(a\leq b\) o \(b\leq a\).
\end{enumerate}
\end{proposicion}

Las tres primeras propiedades expresan que \(\leq\) es una relación de orden. La cuarta indica que dos números naturales cualesquiera son comparables; por eso se dice que el orden de \(\mathbb{N}\) es total.

\begin{advertencia}
El símbolo \(\leq\) no debe leerse únicamente como una instrucción de cálculo. Es una relación entre dos objetos matemáticos. Antes de escribir \(a\leq b\), debe quedar claro qué son \(a\) y \(b\), y en qué conjunto se están considerando.
\end{advertencia}

\section{Principio de inducción}

El principio de inducción matemática es uno de los métodos de demostración más importantes relacionados con los números naturales. Permite demostrar que una propiedad se verifica para todos los números naturales a partir de dos pasos: comprobarla en un caso inicial y demostrar que se transmite de cada número natural al siguiente.

Para formularlo con precisión, sea \(P(n)\) una proposición que depende de un número natural \(n\). Por ejemplo, \(P(n)\) podría ser la afirmación
\[
  0+1+2+\cdots+n=\frac{n(n+1)}{2}.
\]
En este caso, para cada valor de \(n\), la expresión \(P(n)\) representa un enunciado concreto.

\begin{teorema}[Principio de inducción]
Sea \(P(n)\) una proposición dependiente de \(n\in\mathbb{N}\). Supongamos que se cumplen las dos condiciones siguientes:
\begin{enumerate}
  \item \textbf{Caso base:} \(P(0)\) es verdadera.
  \item \textbf{Paso inductivo:} para todo \(n\in\mathbb{N}\), si \(P(n)\) es verdadera, entonces \(P(n+1)\) es verdadera.
\end{enumerate}
Entonces \(P(n)\) es verdadera para todo \(n\in\mathbb{N}\).
\end{teorema}

La proposición \(P(n)\) se denomina \emph{hipótesis inductiva} cuando se usa dentro del paso inductivo. Es importante observar que la hipótesis inductiva no afirma todavía que \(P(n)\) sea verdadera para todo \(n\); se asume provisionalmente para un \(n\) arbitrario con el objetivo de demostrar \(P(n+1)\).

\begin{ejemplo}[Suma de los primeros naturales]
Demostremos que, para todo \(n\in\mathbb{N}\), se cumple
\[
  0+1+2+\cdots+n=\frac{n(n+1)}{2}.
\]

Sea \(P(n)\) la proposición anterior.

\emph{Caso base.} Para \(n=0\), el miembro izquierdo es \(0\), y el miembro derecho es
\[
  \frac{0(0+1)}{2}=0.
\]
Por tanto, \(P(0)\) es verdadera.

\emph{Paso inductivo.} Supongamos que \(P(n)\) es verdadera para un cierto \(n\in\mathbb{N}\), es decir, supongamos que
\[
  0+1+2+\cdots+n=\frac{n(n+1)}{2}.
\]
Queremos demostrar que \(P(n+1)\) es verdadera, es decir, que
\[
  0+1+2+\cdots+n+(n+1)=\frac{(n+1)(n+2)}{2}.
\]
Usando la hipótesis inductiva, obtenemos
\[
\begin{aligned}
  0+1+2+\cdots+n+(n+1)
  &= \frac{n(n+1)}{2}+(n+1) \\
  &= \frac{n(n+1)+2(n+1)}{2} \\
  &= \frac{(n+1)(n+2)}{2}.
\end{aligned}
\]
Por tanto, \(P(n+1)\) es verdadera. Por el principio de inducción, la fórmula se cumple para todo \(n\in\mathbb{N}\).
\end{ejemplo}

\begin{advertencia}
En una demostración por inducción no basta comprobar varios casos particulares. Verificar \(P(0)\), \(P(1)\), \(P(2)\) y \(P(3)\) puede ayudar a descubrir una fórmula, pero no demuestra que \(P(n)\) sea verdadera para todo \(n\in\mathbb{N}\). El paso inductivo es esencial.
\end{advertencia}

\section{Inducción fuerte}

La inducción fuerte es una variante del principio de inducción. Se utiliza cuando, para demostrar \(P(n+1)\), no basta conocer \(P(n)\), sino que conviene conocer todos los casos anteriores.

\begin{teorema}[Principio de inducción fuerte]
Sea \(P(n)\) una proposición dependiente de \(n\in\mathbb{N}\). Supongamos que se cumplen las dos condiciones siguientes:
\begin{enumerate}
  \item \textbf{Caso base:} \(P(0)\) es verdadera.
  \item \textbf{Paso inductivo fuerte:} para todo \(n\in\mathbb{N}\), si \(P(0),P(1),\ldots,P(n)\) son verdaderas, entonces \(P(n+1)\) es verdadera.
\end{enumerate}
Entonces \(P(n)\) es verdadera para todo \(n\in\mathbb{N}\).
\end{teorema}

En el paso inductivo fuerte se supone que la propiedad es cierta para todos los números naturales desde \(0\) hasta \(n\), y con esa información se prueba el caso siguiente.

\begin{ejemplo}[Descomposición en sumas de unos y doses]
Sea \(P(n)\) la proposición: ``el número \(n\) puede escribirse como suma de números iguales a \(1\) o iguales a \(2\)''. Esta afirmación es verdadera para todo \(n\in\mathbb{N}\).

En efecto, \(0\) se escribe como suma vacía. El número \(1\) se escribe como \(1\). Supongamos ahora que la afirmación es verdadera para todos los números naturales menores o iguales que \(n\), con \(n\geq 1\). Entonces \(n+1=(n-1)+2\). Por hipótesis inductiva fuerte, \(n-1\) puede escribirse como suma de unos y doses. Añadiendo un \(2\), obtenemos una escritura de \(n+1\) como suma de unos y doses.
\end{ejemplo}

\begin{advertencia}
La inducción fuerte no es un método ``más verdadero'' que la inducción ordinaria. Ambos principios son equivalentes en \(\mathbb{N}\). La inducción fuerte es, en muchos problemas, una forma más cómoda de organizar la información disponible.
\end{advertencia}

\section{Principio del buen orden}

El conjunto \(\mathbb{N}\) tiene una propiedad fundamental: todo subconjunto no vacío de \(\mathbb{N}\) posee un primer elemento.

\begin{teorema}[Principio del buen orden]
Sea \(A\subseteq\mathbb{N}\). Si \(A\neq\varnothing\), entonces existe un elemento \(m\in A\) tal que
\[
  m\leq a
\]
para todo \(a\in A\).
\end{teorema}

El elemento \(m\) se llama \emph{mínimo} de \(A\), y se denota por
\[
  m=\min A.
\]

\begin{ejemplo}
Sea
\[
  A=\{n\in\mathbb{N}: n\geq 5\}.
\]
Entonces \(A\neq\varnothing\), y su mínimo es \(5\). En efecto, \(5\in A\), y para todo \(a\in A\) se cumple \(5\leq a\).
\end{ejemplo}

El principio del buen orden está estrechamente relacionado con el principio de inducción. Ambos expresan, de formas distintas, que en \(\mathbb{N}\) no es posible descender indefinidamente.

\begin{ejemplo}[Demostración por mínimo contraejemplo]
Supongamos que queremos demostrar que una propiedad \(P(n)\) es verdadera para todo \(n\in\mathbb{N}\). Una estrategia consiste en razonar por contradicción. Se define el conjunto
\[
  A=\{n\in\mathbb{N}: P(n) \text{ es falsa}\}.
\]
Si \(A\) fuera no vacío, tendría un mínimo \(m\). Entonces \(m\) sería el primer contraejemplo. Muchas demostraciones consisten en mostrar que tal primer contraejemplo no puede existir.
\end{ejemplo}

\section{Definiciones recursivas}

Una definición recursiva describe un objeto mediante un caso inicial y una regla que permite construir los casos siguientes.

\begin{definicion}[Definición recursiva]
Una definición recursiva de una sucesión consiste en especificar:
\begin{enumerate}
  \item uno o varios términos iniciales;
  \item una regla que permite calcular un término a partir de términos anteriores.
\end{enumerate}
\end{definicion}

\begin{ejemplo}[Sucesión definida por recurrencia]
Definimos una sucesión \((a_n)_{n\in\mathbb{N}}\) de la siguiente forma:
\[
  a_0=1,
\]
y, para todo \(n\in\mathbb{N}\),
\[
  a_{n+1}=2a_n+1.
\]
Entonces:
\[
  a_0=1,
  \qquad
  a_1=2a_0+1=3,
  \qquad
  a_2=2a_1+1=7,
  \qquad
  a_3=2a_2+1=15.
\]
\end{ejemplo}

En una definición recursiva deben quedar claros todos los símbolos utilizados. En el ejemplo anterior, \((a_n)_{n\in\mathbb{N}}\) denota una sucesión, \(a_0\) es el término inicial, y la fórmula \(a_{n+1}=2a_n+1\) indica cómo obtener el término siguiente a partir del anterior.

\begin{advertencia}
Una regla recursiva no es una fórmula cerrada. La regla \(a_{n+1}=2a_n+1\) permite calcular términos sucesivos, pero no da directamente \(a_n\) en función de \(n\). Encontrar una fórmula cerrada puede requerir una demostración adicional, a menudo por inducción.
\end{advertencia}

\section{Sucesiones y sumas finitas}

Una sucesión es una lista ordenada de objetos indexada por números naturales. En este capítulo consideraremos principalmente sucesiones de números.

\begin{definicion}[Sucesión]
Una sucesión de números reales es una aplicación
\[
  a:\mathbb{N}\longrightarrow\mathbb{R}.
\]
El valor de la sucesión en \(n\in\mathbb{N}\) se denota habitualmente por \(a_n\) en lugar de \(a(n)\). La sucesión completa se escribe como \((a_n)_{n\in\mathbb{N}}\).
\end{definicion}

Aunque el conjunto \(\mathbb{R}\) de los números reales se estudiará con más detalle posteriormente, aquí basta entender que los términos de la sucesión son números.

\begin{definicion}[Suma finita]
Sea \((a_k)_{k\in\mathbb{N}}\) una sucesión y sea \(n\in\mathbb{N}\). La expresión
\[
  \sum_{k=0}^{n} a_k
\]
denota la suma finita
\[
  a_0+a_1+\cdots+a_n.
\]
El símbolo \(k\) se llama índice de suma.
\end{definicion}

\begin{ejemplo}
Si \(a_k=k\) para todo \(k\in\mathbb{N}\), entonces
\[
  \sum_{k=0}^{4} a_k=0+1+2+3+4=10.
\]
También podemos escribir directamente
\[
  \sum_{k=0}^{4} k=0+1+2+3+4=10.
\]
\end{ejemplo}

\begin{proposicion}[Suma de los primeros cuadrados]
Para todo \(n\in\mathbb{N}\), se cumple
\[
  \sum_{k=0}^{n} k^2=\frac{n(n+1)(2n+1)}{6}.
\]
\end{proposicion}

\begin{proof}
La demostración se realiza por inducción sobre \(n\).

Para \(n=0\), tenemos
\[
  \sum_{k=0}^{0} k^2=0^2=0,
\]
y
\[
  \frac{0(0+1)(2\cdot 0+1)}{6}=0.
\]
Por tanto, el caso base es verdadero.

Supongamos que la fórmula es verdadera para un cierto \(n\in\mathbb{N}\), es decir,
\[
  \sum_{k=0}^{n} k^2=\frac{n(n+1)(2n+1)}{6}.
\]
Entonces
\[
\begin{aligned}
  \sum_{k=0}^{n+1} k^2
  &= \sum_{k=0}^{n} k^2+(n+1)^2 \\
  &= \frac{n(n+1)(2n+1)}{6}+(n+1)^2 \\
  &= \frac{(n+1)\bigl(n(2n+1)+6(n+1)\bigr)}{6} \\
  &= \frac{(n+1)(2n^2+7n+6)}{6} \\
  &= \frac{(n+1)(n+2)(2n+3)}{6} \\
  &= \frac{(n+1)((n+1)+1)(2(n+1)+1)}{6}.
\end{aligned}
\]
Esto demuestra el paso inductivo. Por el principio de inducción, la fórmula es verdadera para todo \(n\in\mathbb{N}\).
\end{proof}

\begin{erroresfrecuentes}
  \item Comprobar varios casos particulares y concluir que una fórmula está demostrada.
  \item Escribir el paso inductivo sin declarar explícitamente la hipótesis inductiva.
  \item Confundir \(P(n)\) con \(P(n+1)\).
  \item Usar inducción fuerte sin indicar con precisión qué casos anteriores se están suponiendo verdaderos.
  \item Introducir una sucesión \((a_n)\) sin especificar el conjunto de índices ni los términos iniciales.
\end{erroresfrecuentes}

\begin{seminario}
  \item Demostraciones por inducción básica.
  \item Demostraciones por inducción fuerte.
  \item Fórmulas de sumas finitas.
\end{seminario}

\begin{ejerciciospropuestos}
  \item Demuestra por inducción que, para todo \(n\in\mathbb{N}\),
  \[
    \sum_{k=0}^{n} k=\frac{n(n+1)}{2}.
  \]

  \item Demuestra por inducción que, para todo \(n\in\mathbb{N}\),
  \[
    \sum_{k=0}^{n} (2k+1)=(n+1)^2.
  \]

  \item Define recursivamente una sucesión \((a_n)_{n\in\mathbb{N}}\) mediante
  \[
    a_0=2,
    \qquad
    a_{n+1}=3a_n-1.
  \]
  Calcula \(a_1\), \(a_2\), \(a_3\) y \(a_4\).

  \item Sea \((b_n)_{n\in\mathbb{N}}\) la sucesión definida por
  \[
    b_0=0,
    \qquad
    b_1=1,
    \qquad
    b_{n+2}=b_{n+1}+b_n
  \]
  para todo \(n\in\mathbb{N}\). Calcula sus seis primeros términos.

  \item Demuestra por inducción fuerte que todo número natural \(n\geq 2\) puede escribirse como producto de números primos. Puedes usar la idea de que, si \(n\) no es primo, entonces se descompone como \(n=ab\), con \(1<a<n\) y \(1<b<n\).

  \item Usa el principio del buen orden para justificar que todo subconjunto no vacío de \(\mathbb{N}\) que esté formado por números mayores o iguales que \(10\) tiene un elemento mínimo.
\end{ejerciciospropuestos}

\resumen{El capítulo introduce los números naturales como sistema ordenado y muestra que el principio de inducción, la inducción fuerte y el principio del buen orden son herramientas fundamentales para demostrar propiedades dependientes de un número natural. También se presentan las definiciones recursivas, las sucesiones y las sumas finitas como contextos naturales donde la inducción aparece de forma sistemática.}

% ============================================================
% FdM2
% ============================================================
\part{Fundamentos de Matemáticas II}

% ------------------------------------------------------------
% !TeX root = manual_fundamentos_matematicas.tex
% ------------------------------------------------------------
% Capítulo 7
% ------------------------------------------------------------

\chapter{Los números enteros}

\begin{objetivos}
  \item Introducir el conjunto de los números enteros.
  \item Estudiar sus operaciones y su orden.
  \item Presentar el valor absoluto.
  \item Introducir la división euclídea.
\end{objetivos}

\section{Motivación: la necesidad de ampliar los naturales}

En el capítulo anterior se introdujo el conjunto de los números naturales, que denotamos por
\[
  \mathbb{N}=\{0,1,2,3,\ldots\}
\]
o bien por
\[
  \mathbb{N}^{*}=\{1,2,3,\ldots\},
\]
según se quiera incluir o no el número cero. En este manual adoptaremos la convención
\[
  \mathbb{N}=\{0,1,2,3,\ldots\}.
\]

Los números naturales permiten contar objetos y ordenar cantidades discretas. Sin embargo, presentan una limitación esencial: no siempre es posible resolver ecuaciones de la forma
\[
  a+x=b,
\]
donde \(a\) y \(b\) son números naturales y \(x\) es la incógnita.

Por ejemplo, si queremos resolver
\[
  5+x=3,
\]
no existe ningún número natural \(x\in\mathbb{N}\) que satisfaga la igualdad. Para poder resolver este tipo de ecuaciones es necesario ampliar el sistema numérico e introducir números que representen diferencias negativas.

\begin{ejemplo}
La ecuación
\[
  7+x=10
\]
tiene solución en \(\mathbb{N}\), concretamente \(x=3\). En cambio, la ecuación
\[
  10+x=7
\]
no tiene solución en \(\mathbb{N}\). Para resolverla se introduce el número \(-3\), de manera que
\[
  10+(-3)=7.
\]
\end{ejemplo}

\begin{advertencia}
El símbolo \(-3\) no debe entenderse simplemente como una resta pendiente de realizar. En el conjunto de los números enteros, \(-3\) es un número con entidad propia: el opuesto de \(3\).
\end{advertencia}

\section{El conjunto de los números enteros}

\begin{definicion}[Números enteros]
El conjunto de los números enteros, denotado por \(\mathbb{Z}\), está formado por los números naturales, sus opuestos y el cero:
\[
  \mathbb{Z}=\{\ldots,-3,-2,-1,0,1,2,3,\ldots\}.
\]
\end{definicion}

La letra \(\mathbb{Z}\) procede del término alemán \emph{Zahlen}, que significa ``números''. Los elementos de \(\mathbb{Z}\) se llaman números enteros.

Cada número natural positivo \(n\) tiene asociado un número entero negativo, denotado por \(-n\), que se llama su \emph{opuesto}. El número \(0\) es su propio opuesto, es decir,
\[
  -0=0.
\]

\begin{ejemplo}
Los números
\[
  -5,\quad -1,\quad 0,\quad 2,\quad 17
\]
son enteros. Los números
\[
  \frac{1}{2},\quad \sqrt{2},\quad \pi
\]
no son enteros.
\end{ejemplo}

\begin{convencion}[Inclusión de los naturales en los enteros]
Consideraremos siempre que
\[
  \mathbb{N}\subset \mathbb{Z}.
\]
Por tanto, todo número natural es también un número entero. Sin embargo, no todo número entero es natural: por ejemplo, \(-4\in\mathbb{Z}\), pero \(-4\notin\mathbb{N}\).
\end{convencion}

\section{Suma, producto y orden}

El conjunto \(\mathbb{Z}\) está dotado de dos operaciones fundamentales: la suma y el producto.

La suma de enteros se denota mediante el símbolo \(+\). Si \(a,b\in\mathbb{Z}\), entonces \(a+b\in\mathbb{Z}\). El producto de enteros se denota mediante yuxtaposición o mediante el símbolo \(\cdot\). Si \(a,b\in\mathbb{Z}\), entonces \(ab\in\mathbb{Z}\).

\begin{proposicion}[Propiedades básicas de la suma y del producto]
Para cualesquiera enteros \(a,b,c\in\mathbb{Z}\), se verifican las siguientes propiedades:
\begin{enumerate}
  \item \textbf{Asociatividad de la suma:}
  \[
    (a+b)+c=a+(b+c).
  \]
  \item \textbf{Conmutatividad de la suma:}
  \[
    a+b=b+a.
  \]
  \item \textbf{Elemento neutro de la suma:}
  \[
    a+0=a.
  \]
  \item \textbf{Elemento opuesto:}
  \[
    a+(-a)=0.
  \]
  \item \textbf{Asociatividad del producto:}
  \[
    (ab)c=a(bc).
  \]
  \item \textbf{Conmutatividad del producto:}
  \[
    ab=ba.
  \]
  \item \textbf{Elemento neutro del producto:}
  \[
    a\cdot 1=a.
  \]
  \item \textbf{Distributividad del producto respecto de la suma:}
  \[
    a(b+c)=ab+ac.
  \]
\end{enumerate}
\end{proposicion}

Estas propiedades indican que \(\mathbb{Z}\), con la suma y el producto, posee una estructura algebraica muy importante: es un anillo conmutativo con unidad. Esta terminología se estudiará con más detalle en cursos posteriores de álgebra.

\begin{advertencia}
Aunque todo entero tiene opuesto respecto de la suma, no todo entero distinto de cero tiene inverso respecto del producto dentro de \(\mathbb{Z}\). Por ejemplo, no existe ningún entero \(x\in\mathbb{Z}\) tal que
\[
  2x=1.
\]
Por tanto, \(2\) no tiene inverso multiplicativo en \(\mathbb{Z}\).
\end{advertencia}

El conjunto \(\mathbb{Z}\) también está ordenado. La relación de orden se denota por \(\leq\). Si \(a,b\in\mathbb{Z}\), escribimos
\[
  a\leq b
\]
para indicar que \(a\) es menor o igual que \(b\).

\begin{ejemplo}
Se tienen las desigualdades
\[
  -5<-2<0<3<8.
\]
En particular,
\[
  -5\leq -2,
  \qquad
  0\leq 3,
  \qquad
  -2\leq 8.
\]
\end{ejemplo}

El orden en los enteros es compatible con la suma: si \(a,b,c\in\mathbb{Z}\) y \(a\leq b\), entonces
\[
  a+c\leq b+c.
\]

También es compatible con el producto por enteros positivos: si \(a,b\in\mathbb{Z}\), \(c\in\mathbb{Z}\) y \(c>0\), entonces
\[
  a\leq b \quad \Longrightarrow \quad ac\leq bc.
\]
Si se multiplica por un entero negativo, el sentido de la desigualdad cambia.

\begin{ejemplo}
Como \(2<5\), al multiplicar por \(3>0\), obtenemos
\[
  6<15.
\]
En cambio, al multiplicar por \(-3<0\), obtenemos
\[
  -6>-15.
\]
\end{ejemplo}

\section{Valor absoluto}

El valor absoluto mide la distancia de un número entero al cero en la recta numérica.

\begin{definicion}[Valor absoluto]
Sea \(a\in\mathbb{Z}\). El valor absoluto de \(a\), denotado por \(|a|\), se define por
\[
  |a|=
  \begin{cases}
    a, & \text{si } a\geq 0,\\
    -a, & \text{si } a<0.
  \end{cases}
\]
\end{definicion}

De esta definición se deduce que \(|a|\) es siempre un número natural.

\begin{ejemplo}
Se tiene
\[
  |5|=5,
  \qquad
  |-5|=5,
  \qquad
  |0|=0.
\]
\end{ejemplo}

\begin{proposicion}[Propiedades elementales del valor absoluto]
Para cualesquiera enteros \(a,b\in\mathbb{Z}\), se verifican las siguientes propiedades:
\begin{enumerate}
  \item \(|a|\geq 0\).
  \item \(|a|=0\) si y solo si \(a=0\).
  \item \(|-a|=|a|\).
  \item \(|ab|=|a|\,|b|\).
  \item \(|a+b|\leq |a|+|b|\).
\end{enumerate}
\end{proposicion}

La última desigualdad se conoce como \emph{desigualdad triangular}. Su interpretación es que, si se avanza primero una distancia \(|a|\) y después una distancia \(|b|\), la distancia total al origen no puede ser mayor que la suma de ambas distancias.

\begin{ejemplo}
Si \(a=-3\) y \(b=5\), entonces
\[
  |a+b|=|-3+5|=|2|=2,
\]
mientras que
\[
  |a|+|b|=|-3|+|5|=3+5=8.
\]
Por tanto,
\[
  |a+b|\leq |a|+|b|.
\]
\end{ejemplo}

\section{División euclídea}

La división euclídea formaliza la división entera con resto. Es uno de los resultados fundamentales de la aritmética de los enteros y será esencial para estudiar la divisibilidad, el máximo común divisor y la aritmética modular.

\begin{teorema}[División euclídea]
Sean \(a,b\in\mathbb{Z}\), con \(b>0\). Entonces existen unos únicos enteros \(q,r\in\mathbb{Z}\) tales que
\[
  a=bq+r
\]
y
\[
  0\leq r<b.
\]
El número \(q\) se llama cociente y el número \(r\) se llama resto de la división de \(a\) entre \(b\).
\end{teorema}

En este enunciado, el entero \(a\) es el dividendo y el entero positivo \(b\) es el divisor.

\begin{ejemplo}
Si \(a=17\) y \(b=5\), entonces
\[
  17=5\cdot 3+2,
\]
con
\[
  0\leq 2<5.
\]
Por tanto, el cociente es \(q=3\) y el resto es \(r=2\).
\end{ejemplo}

\begin{ejemplo}
Si \(a=-17\) y \(b=5\), debemos encontrar enteros \(q\) y \(r\) tales que
\[
  -17=5q+r,
  \qquad
  0\leq r<5.
\]
Tomando \(q=-4\) y \(r=3\), se obtiene
\[
  -17=5(-4)+3=-20+3.
\]
Por tanto, el cociente es \(q=-4\) y el resto es \(r=3\).
\end{ejemplo}

\begin{advertencia}
Cuando el dividendo es negativo, el resto de la división euclídea sigue siendo no negativo. Por ejemplo, en la división de \(-17\) entre \(5\), el resto es \(3\), no \(-2\), porque el resto debe satisfacer
\[
  0\leq r<5.
\]
\end{advertencia}

\section*{Errores frecuentes}

\begin{erroresfrecuentes}
  \item Pensar que los números enteros negativos son números naturales. Con la convención de este manual, \(-1\notin\mathbb{N}\), aunque \(-1\in\mathbb{Z}\).
  \item Olvidar que al multiplicar una desigualdad por un número negativo se invierte el sentido de la desigualdad.
  \item Confundir el opuesto de un número, \(-a\), con su valor absoluto, \(|a|\).
  \item Creer que todos los enteros no nulos tienen inverso multiplicativo dentro de \(\mathbb{Z}\).
  \item Calcular incorrectamente el resto en una división euclídea con dividendo negativo.
\end{erroresfrecuentes}

\begin{seminario}
  \item Ejercicios de orden y valor absoluto.
  \item Aplicaciones de la división euclídea.
\end{seminario}

\begin{ejerciciospropuestos}
  \item Calcula el valor absoluto de los enteros \(-12\), \(0\), \(7\), \(-1\) y \(25\).
  \item Ordena de menor a mayor los enteros
  \[
    -3,
    \quad 5,
    \quad -8,
    \quad 0,
    \quad 2.
  \]
  \item Encuentra el cociente y el resto de la división euclídea de \(43\) entre \(6\).
  \item Encuentra el cociente y el resto de la división euclídea de \(-43\) entre \(6\).
  \item Demuestra que, para todo \(a\in\mathbb{Z}\), se tiene \(|-a|=|a|\).
  \item Da un ejemplo de dos enteros \(a,b\in\mathbb{Z}\) tales que
  \[
    |a+b|<|a|+|b|.
  \]
  \item Explica por qué no existe ningún entero \(x\in\mathbb{Z}\) tal que \(3x=1\).
\end{ejerciciospropuestos}

\resumen{Los números enteros amplían los números naturales al incorporar opuestos aditivos. El conjunto \(\mathbb{Z}\) permite resolver ecuaciones de la forma \(a+x=b\), está dotado de suma, producto y orden, y posee una herramienta aritmética fundamental: la división euclídea, que expresa todo entero como múltiplo de un divisor positivo más un resto acotado.}

% ------------------------------------------------------------

% !TEX root = manual_fundamentos_matematicas.tex
% ------------------------------------------------------------
% Capítulo 8. Divisibilidad
% ------------------------------------------------------------

\chapter{Divisibilidad}

\begin{objetivos}
  \item Definir la relación de divisibilidad en los enteros.
  \item Estudiar divisores, múltiplos y números primos.
  \item Introducir máximo común divisor y mínimo común múltiplo.
  \item Presentar el algoritmo de Euclides y la identidad de Bézout.
\end{objetivos}

\section{Motivación}

La divisibilidad estudia cuándo un número entero puede obtenerse multiplicando otro número entero por un tercer entero. Esta idea, aparentemente elemental, es una de las bases de la aritmética. Permite hablar de divisores, múltiplos, números primos, factorizaciones, congruencias y ecuaciones diofánticas.

En este capítulo trabajaremos principalmente con el conjunto de los números enteros, que denotamos por \(\mathbb{Z}\). Recordemos que
\[
\mathbb{Z}=\{\ldots,-3,-2,-1,0,1,2,3,\ldots\}.
\]

\begin{convencion}[Símbolos utilizados en este capítulo]
A lo largo de este capítulo, salvo que se indique lo contrario, las letras \(a,b,c,d,m,n,q,r\) representarán números enteros. Cuando se utilicen números naturales positivos, se indicará explícitamente mediante la expresión \(n\in \mathbb{N}\) con \(n\geq 1\), o bien \(n\in \mathbb{N}^{\ast}\), donde \(\mathbb{N}^{\ast}\) denota el conjunto de los números naturales positivos.
\end{convencion}

\section{Divisores y múltiplos}

La noción central de este capítulo es la relación de divisibilidad.

\begin{definicion}[Divisibilidad]
Sean \(a,b\in\mathbb{Z}\). Diremos que \(a\) divide a \(b\), y escribiremos
\[
a\mid b,
\]
si existe un entero \(q\in\mathbb{Z}\) tal que
\[
b=aq.
\]
En este caso, también decimos que \(a\) es un divisor de \(b\), o que \(b\) es un múltiplo de \(a\).
\end{definicion}

El símbolo \(a\mid b\) no representa una fracción. Es una afirmación lógica: significa que \(b\) puede escribirse como producto de \(a\) por un entero.

\begin{ejemplo}
Tenemos \(3\mid 12\), porque existe \(q=4\in\mathbb{Z}\) tal que
\[
12=3\cdot 4.
\]
También \((-3)\mid 12\), porque
\[
12=(-3)\cdot(-4).
\]
En cambio, \(5\nmid 12\), porque no existe ningún entero \(q\) tal que \(12=5q\).
\end{ejemplo}

\begin{advertencia}
La afirmación \(a\mid b\) no significa que \(a/b\) sea un número entero escrito como cociente, sino que existe un entero \(q\) tal que \(b=aq\). La definición debe expresarse siempre mediante una existencia: \(\exists q\in\mathbb{Z}\) tal que \(b=aq\).
\end{advertencia}

\begin{proposicion}[Propiedades elementales de la divisibilidad]
Sean \(a,b,c\in\mathbb{Z}\). Se verifican las siguientes propiedades:
\begin{enumerate}
  \item Si \(a\mid b\) y \(b\mid c\), entonces \(a\mid c\).
  \item Si \(a\mid b\) y \(a\mid c\), entonces \(a\mid (b+c)\).
  \item Si \(a\mid b\), entonces \(a\mid bc\).
  \item Si \(a\mid b\) y \(a\mid c\), entonces \(a\mid (mb+nc)\) para cualesquiera \(m,n\in\mathbb{Z}\).
\end{enumerate}
\end{proposicion}

\begin{proof}
Demostramos la cuarta propiedad, que contiene a varias de las anteriores como casos particulares. Supongamos que \(a\mid b\) y \(a\mid c\). Por definición de divisibilidad, existen \(q_1,q_2\in\mathbb{Z}\) tales que
\[
b=aq_1,
\qquad
c=aq_2.
\]
Sean ahora \(m,n\in\mathbb{Z}\). Entonces
\[
mb+nc=m(aq_1)+n(aq_2)=a(mq_1+nq_2).
\]
Como \(mq_1+nq_2\in\mathbb{Z}\), se concluye que \(a\mid (mb+nc)\).
\end{proof}

\begin{definicion}[Divisores de un entero]
Sea \(b\in\mathbb{Z}\). El conjunto de los divisores de \(b\) se denota por
\[
\operatorname{Div}(b)=\{a\in\mathbb{Z}:a\mid b\}.
\]
\end{definicion}

\begin{ejemplo}
Los divisores enteros de \(12\) son
\[
\operatorname{Div}(12)=\{-12,-6,-4,-3,-2,-1,1,2,3,4,6,12\}.
\]
Si solo queremos considerar divisores positivos, escribimos
\[
\operatorname{Div}^{+}(12)=\{1,2,3,4,6,12\}.
\]
\end{ejemplo}

\begin{advertencia}
El número \(0\) tiene una situación especial. Si \(a\in\mathbb{Z}\) y \(a\neq 0\), entonces \(a\mid 0\), porque \(0=a\cdot 0\). En cambio, \(0\mid b\) solo ocurre cuando \(b=0\), ya que \(b=0\cdot q=0\). Por esta razón, en muchos enunciados se exige explícitamente que el divisor sea no nulo.
\end{advertencia}

\section{Números primos}

Los números primos son los enteros positivos que no admiten divisores positivos no triviales.

\begin{definicion}[Número primo]
Sea \(p\in\mathbb{N}\) con \(p\geq 2\). Diremos que \(p\) es un número primo si sus únicos divisores positivos son \(1\) y \(p\). Es decir,
\[
\operatorname{Div}^{+}(p)=\{1,p\}.
\]
\end{definicion}

\begin{definicion}[Número compuesto]
Sea \(n\in\mathbb{N}\) con \(n\geq 2\). Diremos que \(n\) es compuesto si no es primo. Equivalentemente, \(n\) es compuesto si existen \(a,b\in\mathbb{N}\) tales que
\[
1<a<n,
\qquad
1<b<n,
\qquad
n=ab.
\]
\end{definicion}

\begin{ejemplo}
Los números \(2,3,5,7,11,13\) son primos. El número \(12\) es compuesto porque
\[
12=3\cdot 4,
\]
con \(1<3<12\) y \(1<4<12\).
\end{ejemplo}

\begin{advertencia}
El número \(1\) no se considera primo. Esta convención es esencial para que la factorización de un entero positivo en factores primos sea única, salvo el orden de los factores.
\end{advertencia}

\begin{proposicion}
Sea \(n\in\mathbb{N}\) con \(n\geq 2\). Si \(n\) no es primo, entonces existe un primo \(p\) tal que \(p\mid n\) y \(p\leq n\).
\end{proposicion}

\begin{proof}
Como \(n\) no es primo, posee un divisor positivo \(d\) tal que \(1<d<n\). Consideremos el conjunto de los divisores positivos de \(n\) mayores que \(1\):
\[
D=\{d\in\mathbb{N}: d\mid n \text{ y } d>1\}.
\]
El conjunto \(D\) no es vacío, porque \(n\in D\). Por el principio del buen orden, \(D\) tiene un menor elemento, que denotamos por \(p\). Veamos que \(p\) es primo. Si no lo fuera, existiría un divisor positivo \(e\) de \(p\) con \(1<e<p\). Como \(e\mid p\) y \(p\mid n\), se tendría \(e\mid n\), lo que contradice que \(p\) sea el menor elemento de \(D\). Por tanto, \(p\) es primo y divide a \(n\).
\end{proof}

\section{Máximo común divisor}

El máximo común divisor mide la parte común de divisibilidad de dos enteros.

\begin{definicion}[Divisor común]
Sean \(a,b\in\mathbb{Z}\), no ambos nulos. Un entero \(d\in\mathbb{Z}\) es un divisor común de \(a\) y \(b\) si
\[
d\mid a
\qquad\text{y}\qquad
d\mid b.
\]
\end{definicion}

\begin{definicion}[Máximo común divisor]
Sean \(a,b\in\mathbb{Z}\), no ambos nulos. El máximo común divisor de \(a\) y \(b\) es el mayor entero positivo que divide simultáneamente a \(a\) y a \(b\). Se denota por
\[
\mcd(a,b)
\]
o, en notación española, por
\[
\operatorname{mcd}(a,b).
\]
En este manual utilizaremos preferentemente la notación \(\operatorname{mcd}(a,b)\).
\end{definicion}

\begin{ejemplo}
Los divisores positivos de \(18\) son
\[
1,2,3,6,9,18,
\]
y los divisores positivos de \(30\) son
\[
1,2,3,5,6,10,15,30.
\]
Los divisores comunes positivos son
\[
1,2,3,6.
\]
Por tanto,
\[
\operatorname{mcd}(18,30)=6.
\]
\end{ejemplo}

\begin{definicion}[Enteros coprimos]
Sean \(a,b\in\mathbb{Z}\), no ambos nulos. Diremos que \(a\) y \(b\) son coprimos, o primos entre sí, si
\[
\operatorname{mcd}(a,b)=1.
\]
\end{definicion}

\begin{ejemplo}
Los números \(8\) y \(15\) son coprimos, porque sus únicos divisores comunes positivos se reducen a \(1\). Es decir,
\[
\operatorname{mcd}(8,15)=1.
\]
\end{ejemplo}

\section{Mínimo común múltiplo}

El mínimo común múltiplo es la noción dual del máximo común divisor en el contexto de los múltiplos positivos.

\begin{definicion}[Múltiplo común]
Sean \(a,b\in\mathbb{Z}\), con \(a\neq 0\) y \(b\neq 0\). Un entero \(m\in\mathbb{Z}\) es un múltiplo común de \(a\) y \(b\) si
\[
a\mid m
\qquad\text{y}\qquad
b\mid m.
\]
\end{definicion}

\begin{definicion}[Mínimo común múltiplo]
Sean \(a,b\in\mathbb{Z}\), con \(a\neq 0\) y \(b\neq 0\). El mínimo común múltiplo de \(a\) y \(b\) es el menor entero positivo que es múltiplo común de \(a\) y de \(b\). Se denota por
\[
\operatorname{mcm}(a,b).
\]
\end{definicion}

\begin{ejemplo}
Los múltiplos positivos de \(6\) son
\[
6,12,18,24,30,36,\ldots,
\]
y los múltiplos positivos de \(8\) son
\[
8,16,24,32,40,\ldots.
\]
El menor múltiplo común positivo es \(24\). Por tanto,
\[
\operatorname{mcm}(6,8)=24.
\]
\end{ejemplo}

\begin{proposicion}[Relación entre \(\operatorname{mcd}\) y \(\operatorname{mcm}\)]
Sean \(a,b\in\mathbb{Z}\), con \(a\neq 0\) y \(b\neq 0\). Entonces
\[
\operatorname{mcd}(a,b)\operatorname{mcm}(a,b)=|ab|.
\]
\end{proposicion}

\begin{advertencia}
La fórmula anterior es válida para dos enteros no nulos. Antes de aplicarla, hay que asegurarse de que ambos números son distintos de cero y de que se está usando el valor absoluto de \(ab\).
\end{advertencia}

\section{Algoritmo de Euclides}

El algoritmo de Euclides permite calcular el máximo común divisor de dos enteros de manera eficiente mediante divisiones euclídeas sucesivas.

Recordemos la división euclídea: si \(a,b\in\mathbb{Z}\) y \(b>0\), existen enteros únicos \(q,r\in\mathbb{Z}\) tales que
\[
a=bq+r,
\qquad
0\leq r<b.
\]
El entero \(q\) se llama cociente y el entero \(r\) se llama resto.

\begin{lema}[Invariancia del máximo común divisor]
Sean \(a,b,q,r\in\mathbb{Z}\), con \(b\neq 0\), tales que
\[
a=bq+r.
\]
Entonces
\[
\operatorname{mcd}(a,b)=\operatorname{mcd}(b,r).
\]
\end{lema}

\begin{proof}
Sea \(d\in\mathbb{Z}\). Supongamos primero que \(d\mid a\) y \(d\mid b\). Como \(r=a-bq\), se deduce que \(d\mid r\). Por tanto, todo divisor común de \(a\) y \(b\) es divisor común de \(b\) y \(r\).

Recíprocamente, si \(d\mid b\) y \(d\mid r\), entonces \(a=bq+r\) implica que \(d\mid a\). Por tanto, los divisores comunes de \(a\) y \(b\) coinciden con los divisores comunes de \(b\) y \(r\). En consecuencia, sus máximos comunes divisores son iguales.
\end{proof}

\begin{ejemplo}[Algoritmo de Euclides]
Calculemos \(\operatorname{mcd}(252,105)\). Aplicamos divisiones euclídeas sucesivas:
\[
252=105\cdot 2+42,
\]
\[
105=42\cdot 2+21,
\]
\[
42=21\cdot 2+0.
\]
El último resto no nulo es \(21\). Por tanto,
\[
\operatorname{mcd}(252,105)=21.
\]
\end{ejemplo}

\begin{advertencia}
En el algoritmo de Euclides, el máximo común divisor es el último resto no nulo, no el primer resto que aparece ni el último cociente.
\end{advertencia}

\section{Identidad de Bézout}

La identidad de Bézout expresa el máximo común divisor como una combinación lineal entera de los dos números iniciales.

\begin{definicion}[Combinación lineal entera]
Sean \(a,b\in\mathbb{Z}\). Una combinación lineal entera de \(a\) y \(b\) es un número de la forma
\[
ax+by,
\]
donde \(x,y\in\mathbb{Z}\).
\end{definicion}

\begin{teorema}[Identidad de Bézout]
Sean \(a,b\in\mathbb{Z}\), no ambos nulos, y sea
\[
d=\operatorname{mcd}(a,b).
\]
Entonces existen enteros \(x,y\in\mathbb{Z}\) tales que
\[
d=ax+by.
\]
\end{teorema}

\begin{ejemplo}[Cálculo de coeficientes de Bézout]
Busquemos enteros \(x,y\) tales que
\[
\operatorname{mcd}(252,105)=252x+105y.
\]
Del ejemplo anterior sabemos que \(\operatorname{mcd}(252,105)=21\). A partir del algoritmo de Euclides:
\[
252=105\cdot 2+42,
\]
\[
105=42\cdot 2+21.
\]
Despejamos hacia atrás:
\[
21=105-42\cdot 2.
\]
Como
\[
42=252-105\cdot 2,
\]
entonces
\[
21=105-2(252-105\cdot 2)=105\cdot 5-252\cdot 2.
\]
Por tanto,
\[
21=252(-2)+105(5).
\]
Así, una solución es \(x=-2\) e \(y=5\).
\end{ejemplo}

\begin{corolario}
Sean \(a,b\in\mathbb{Z}\), no ambos nulos. Entonces \(a\) y \(b\) son coprimos si y solo si existen \(x,y\in\mathbb{Z}\) tales que
\[
ax+by=1.
\]
\end{corolario}

\begin{proof}
Si \(a\) y \(b\) son coprimos, entonces \(\operatorname{mcd}(a,b)=1\). Por la identidad de Bézout, existen \(x,y\in\mathbb{Z}\) tales que \(ax+by=1\).

Recíprocamente, supongamos que existen \(x,y\in\mathbb{Z}\) tales que \(ax+by=1\). Si \(d\) es un divisor común positivo de \(a\) y \(b\), entonces \(d\mid ax+by\). Por tanto, \(d\mid 1\), luego \(d=1\). Así, el máximo común divisor de \(a\) y \(b\) es \(1\).
\end{proof}

\section{Ecuaciones diofánticas lineales}

Una ecuación diofántica es una ecuación en la que se buscan soluciones enteras. En este curso estudiaremos el caso lineal con dos incógnitas.

\begin{definicion}[Ecuación diofántica lineal]
Sean \(a,b,c\in\mathbb{Z}\). Una ecuación diofántica lineal en dos incógnitas es una ecuación de la forma
\[
ax+by=c,
\]
donde se buscan soluciones \((x,y)\in\mathbb{Z}^2\).
\end{definicion}

\begin{teorema}[Criterio de resolubilidad]
Sean \(a,b,c\in\mathbb{Z}\), con \(a\) y \(b\) no ambos nulos. La ecuación
\[
ax+by=c
\]
tiene solución entera si y solo si
\[
\operatorname{mcd}(a,b)\mid c.
\]
\end{teorema}

\begin{proof}
Sea \(d=\operatorname{mcd}(a,b)\). Supongamos primero que existe una solución entera \((x,y)\in\mathbb{Z}^2\) de la ecuación. Como \(d\mid a\) y \(d\mid b\), se tiene \(d\mid ax+by\). Pero \(ax+by=c\), luego \(d\mid c\).

Recíprocamente, supongamos que \(d\mid c\). Entonces existe \(k\in\mathbb{Z}\) tal que \(c=dk\). Por la identidad de Bézout, existen \(u,v\in\mathbb{Z}\) tales que
\[
d=au+bv.
\]
Multiplicando por \(k\), obtenemos
\[
c=dk=a(uk)+b(vk).
\]
Por tanto, \((x,y)=(uk,vk)\) es una solución entera de \(ax+by=c\).
\end{proof}

\begin{ejemplo}
Estudiemos la ecuación
\[
18x+30y=12.
\]
Calculamos
\[
\operatorname{mcd}(18,30)=6.
\]
Como \(6\mid 12\), la ecuación tiene soluciones enteras.

Buscamos una identidad de Bézout para \(18\) y \(30\). Como
\[
30=18\cdot 1+12,
\qquad
18=12\cdot 1+6,
\]
tenemos
\[
6=18-12=18-(30-18)=2\cdot 18-30.
\]
Multiplicando por \(2\), obtenemos
\[
12=4\cdot 18-2\cdot 30.
\]
Por tanto, una solución es
\[
x=4,
\qquad
y=-2.
\]
\end{ejemplo}

\begin{erroresfrecuentes}
  \item Confundir la expresión \(a\mid b\) con el cociente \(a/b\).
  \item Olvidar que, para escribir \(a\mid b\), debe existir un entero \(q\) tal que \(b=aq\).
  \item Considerar erróneamente que \(1\) es un número primo.
  \item Tomar el último cociente del algoritmo de Euclides en lugar del último resto no nulo.
  \item Aplicar el criterio de ecuaciones diofánticas sin comprobar primero si \(\operatorname{mcd}(a,b)\mid c\).
\end{erroresfrecuentes}

\begin{seminario}
  \item Cálculo del máximo común divisor.
  \item Uso del algoritmo de Euclides extendido.
  \item Resolución de ecuaciones diofánticas lineales.
\end{seminario}

\begin{ejerciciospropuestos}
  \item Determina si se verifican las siguientes relaciones de divisibilidad: \(4\mid 28\), \(6\mid 35\), \((-5)\mid 20\), \(7\mid 0\).
  \item Calcula \(\operatorname{mcd}(84,126)\) mediante el algoritmo de Euclides.
  \item Calcula \(\operatorname{mcm}(12,18)\) y comprueba la relación entre \(\operatorname{mcd}\) y \(\operatorname{mcm}\).
  \item Encuentra enteros \(x,y\) tales que \(\operatorname{mcd}(84,126)=84x+126y\).
  \item Decide si la ecuación diofántica \(21x+35y=14\) tiene soluciones enteras. En caso afirmativo, encuentra una solución.
  \item Decide si la ecuación diofántica \(12x+18y=5\) tiene soluciones enteras.
\end{ejerciciospropuestos}

\resumen{La divisibilidad permite estudiar la estructura aritmética de los enteros. La relación \(a\mid b\) significa que \(b\) es un múltiplo entero de \(a\). A partir de esta noción se definen divisores, múltiplos, números primos, máximo común divisor y mínimo común múltiplo. El algoritmo de Euclides permite calcular el máximo común divisor, y la identidad de Bézout lo expresa como combinación lineal entera. Estas herramientas permiten resolver ecuaciones diofánticas lineales de la forma \(ax+by=c\).}



% ------------------------------------------------------------

\chapter{El Teorema Fundamental de la Aritmética}

\begin{objetivos}
  \item Enunciar y demostrar el Teorema Fundamental de la Aritmética.
  \item Comprender la existencia y unicidad de la factorización en primos.
  \item Aplicar la factorización al cálculo de divisores, máximo común divisor y mínimo común múltiplo.
\end{objetivos}

\section{Factorización en primos}

En este capítulo se estudia uno de los resultados centrales de la aritmética elemental: todo entero positivo mayor que uno puede escribirse como producto de números primos, y esta escritura es única salvo el orden de los factores.

Antes de formular el resultado, fijamos la notación que se utilizará durante todo el capítulo.

\begin{convencion}[Notación básica]
Denotaremos por \(\mathbb{N}\) el conjunto de los números naturales y por \(\mathbb{Z}\) el conjunto de los números enteros. Si \(a,b\in\mathbb{Z}\), escribiremos \(a\mid b\) para indicar que \(a\) divide a \(b\), es decir, que existe un entero \(k\in\mathbb{Z}\) tal que \(b=ak\).
\end{convencion}

\begin{definicion}[Número primo]
Sea \(p\in\mathbb{N}\) con \(p>1\). Diremos que \(p\) es un \emph{número primo} si sus únicos divisores positivos son \(1\) y \(p\).
\end{definicion}

\begin{definicion}[Número compuesto]
Sea \(n\in\mathbb{N}\) con \(n>1\). Diremos que \(n\) es \emph{compuesto} si no es primo. Equivalentemente, \(n\) es compuesto si existen enteros positivos \(a,b\in\mathbb{N}\), ambos mayores que \(1\), tales que
\[
  n=ab.
\]
\end{definicion}

\begin{ejemplo}
Los números \(2,3,5,7,11\) son primos. El número \(12\) no es primo porque
\[
  12=3\cdot 4,
\]
y tanto \(3\) como \(4\) son divisores positivos de \(12\) distintos de \(1\) y de \(12\).
\end{ejemplo}

\begin{advertencia}
El número \(1\) no se considera primo. Esta convención es esencial para que la factorización en primos sea única. Si \(1\) fuese primo, podríamos escribir, por ejemplo,
\[
  6=2\cdot 3=1\cdot 2\cdot 3=1^2\cdot 2\cdot 3,
\]
y la unicidad dejaría de tener sentido.
\end{advertencia}

\begin{definicion}[Factorización en primos]
Sea \(n\in\mathbb{N}\) con \(n>1\). Una \emph{factorización en primos} de \(n\) es una escritura de la forma
\[
  n=p_1p_2\cdots p_r,
\]
donde \(r\in\mathbb{N}\), y donde \(p_1,p_2,\ldots,p_r\) son números primos.
\end{definicion}

Por ejemplo,
\[
  60=2\cdot 2\cdot 3\cdot 5.
\]
También podemos escribir esta factorización agrupando los factores iguales:
\[
  60=2^2\cdot 3\cdot 5.
\]

\section{Existencia de la factorización}

El primer paso consiste en demostrar que todo entero positivo mayor que \(1\) admite al menos una factorización en primos.

\begin{teorema}[Existencia de factorización en primos]
Sea \(n\in\mathbb{N}\) con \(n>1\). Entonces \(n\) puede escribirse como producto de números primos.
\end{teorema}

\begin{proof}
Demostraremos el resultado usando el principio del buen orden.

Supongamos, para llegar a una contradicción, que existe algún entero positivo mayor que \(1\) que no puede escribirse como producto de primos. Consideremos el conjunto
\[
  A=\{n\in\mathbb{N}: n>1 \text{ y } n \text{ no puede escribirse como producto de primos}\}.
\]
Por hipótesis, \(A\) no es vacío. Por el principio del buen orden, \(A\) tiene un menor elemento. Lo denotamos por \(m\).

El número \(m\) no puede ser primo, porque si \(m\) fuese primo, entonces ya estaría escrito como producto de un único primo, él mismo. Por tanto, \(m\) es compuesto. Luego existen enteros \(a,b\in\mathbb{N}\) tales que
\[
  m=ab,
\]
con
\[
  1<a<m, \qquad 1<b<m.
\]
Como \(m\) es el menor elemento de \(A\), los números \(a\) y \(b\) no pertenecen a \(A\). Por tanto, ambos pueden escribirse como producto de primos. Al multiplicar las factorizaciones de \(a\) y \(b\), obtenemos una factorización de \(m=ab\) como producto de primos.

Esto contradice que \(m\in A\). Por tanto, el conjunto \(A\) debe ser vacío. En consecuencia, todo entero positivo mayor que \(1\) puede escribirse como producto de primos.
\end{proof}

\begin{ejemplo}
Veamos la existencia en un caso concreto. Como
\[
  84=2\cdot 42=2\cdot 2\cdot 21=2\cdot 2\cdot 3\cdot 7,
\]
obtenemos la factorización
\[
  84=2^2\cdot 3\cdot 7.
\]
\end{ejemplo}

\section{Unicidad de la factorización}

La existencia asegura que podemos descomponer un entero positivo mayor que \(1\) en factores primos. Falta demostrar que esta descomposición es única, salvo el orden de los factores.

Para ello necesitamos un resultado previo: el lema de Euclides.

\begin{lema}[Lema de Euclides]
Sea \(p\) un número primo y sean \(a,b\in\mathbb{Z}\). Si
\[
  p\mid ab,
\]
entonces
\[
  p\mid a \quad \text{o} \quad p\mid b.
\]
\end{lema}

\begin{proof}
Supongamos que \(p\mid ab\). Si \(p\mid a\), no hay nada que demostrar. Supongamos entonces que \(p\nmid a\).

Como \(p\) es primo, sus únicos divisores positivos son \(1\) y \(p\). La hipótesis \(p\nmid a\) implica que \(p\) no es divisor común de \(p\) y \(a\). Por tanto,
\[
  \mcd(p,a)=1.
\]
Por la identidad de Bézout, existen enteros \(u,v\in\mathbb{Z}\) tales que
\[
  up+va=1.
\]
Multiplicando por \(b\), obtenemos
\[
  upb+vab=b.
\]
Como \(p\mid upb\) y, por hipótesis, \(p\mid ab\), también se tiene \(p\mid vab\). Por tanto, \(p\) divide la suma \(upb+vab\), es decir,
\[
  p\mid b.
\]
Así, si \(p\mid ab\), entonces \(p\mid a\) o \(p\mid b\).
\end{proof}

\begin{corolario}
Sea \(p\) un número primo y sean \(a_1,a_2,\ldots,a_r\in\mathbb{Z}\). Si
\[
  p\mid a_1a_2\cdots a_r,
\]
entonces existe algún índice \(i\in\{1,2,\ldots,r\}\) tal que
\[
  p\mid a_i.
\]
\end{corolario}

\begin{proof}
Se demuestra aplicando repetidamente el lema de Euclides. Para \(r=2\) es exactamente el lema anterior. Para \(r>2\), se aplica el lema a
\[
  a_1(a_2\cdots a_r).
\]
Si \(p\mid a_1\), hemos terminado. Si no, entonces \(p\mid a_2\cdots a_r\), y se repite el razonamiento.
\end{proof}

\begin{teorema}[Teorema Fundamental de la Aritmética]
Sea \(n\in\mathbb{N}\) con \(n>1\). Entonces \(n\) puede escribirse como producto de números primos. Además, esta escritura es única salvo el orden de los factores.
\end{teorema}

\begin{proof}
La existencia ya ha sido demostrada. Probemos la unicidad.

Supongamos que \(n\) admite dos factorizaciones en primos:
\[
  n=p_1p_2\cdots p_r=q_1q_2\cdots q_s,
\]
donde todos los \(p_i\) y todos los \(q_j\) son primos.

Como \(p_1\mid n\), tenemos
\[
  p_1\mid q_1q_2\cdots q_s.
\]
Por el corolario del lema de Euclides, existe algún índice \(j\) tal que
\[
  p_1\mid q_j.
\]
Pero \(q_j\) es primo y \(p_1\) también es primo. Como \(p_1\) divide a \(q_j\), necesariamente
\[
  p_1=q_j.
\]
Reordenando los factores \(q_1,\ldots,q_s\), podemos suponer que \(q_j=q_1\). Entonces cancelamos el factor común \(p_1=q_1\) y obtenemos
\[
  p_2\cdots p_r=q_2\cdots q_s.
\]
Repitiendo este argumento, cada primo de la primera factorización coincide con un primo de la segunda. Por tanto, ambas factorizaciones tienen los mismos factores primos, con las mismas repeticiones, salvo el orden.
\end{proof}

\begin{convencion}[Forma canónica]
Gracias al Teorema Fundamental de la Aritmética, todo entero \(n>1\) puede escribirse de manera única en la forma
\[
  n=p_1^{\alpha_1}p_2^{\alpha_2}\cdots p_k^{\alpha_k},
\]
donde \(p_1,p_2,\ldots,p_k\) son primos distintos, normalmente ordenados como
\[
  p_1<p_2<\cdots<p_k,
\]
y donde \(\alpha_1,\alpha_2,\ldots,\alpha_k\in\mathbb{N}\) son exponentes positivos.
\end{convencion}

\section{Aplicaciones}

El Teorema Fundamental de la Aritmética permite transformar problemas de divisibilidad en problemas sobre exponentes de primos.

\subsection*{Cálculo de divisores}

Sea
\[
  n=p_1^{\alpha_1}p_2^{\alpha_2}\cdots p_k^{\alpha_k}
\]
la factorización canónica de \(n\). Entonces todo divisor positivo \(d\) de \(n\) tiene la forma
\[
  d=p_1^{\beta_1}p_2^{\beta_2}\cdots p_k^{\beta_k},
\]
donde
\[
  0\leq \beta_i\leq \alpha_i
\]
para cada índice \(i\in\{1,2,\ldots,k\}\).

\begin{ejemplo}
Como
\[
  60=2^2\cdot 3\cdot 5,
\]
un divisor positivo de \(60\) tiene la forma
\[
  2^a3^b5^c,
\]
donde
\[
  0\leq a\leq 2, \qquad 0\leq b\leq 1, \qquad 0\leq c\leq 1.
\]
Por ejemplo, tomando \(a=1\), \(b=1\), \(c=0\), obtenemos el divisor
\[
  2^1 3^1 5^0=6.
\]
\end{ejemplo}

\subsection*{Número de divisores positivos}

Si
\[
  n=p_1^{\alpha_1}p_2^{\alpha_2}\cdots p_k^{\alpha_k},
\]
entonces el número de divisores positivos de \(n\) es
\[
  (\alpha_1+1)(\alpha_2+1)\cdots(\alpha_k+1).
\]

\begin{ejemplo}
Para \(60=2^2\cdot 3^1\cdot 5^1\), el número de divisores positivos es
\[
  (2+1)(1+1)(1+1)=12.
\]
\end{ejemplo}

\subsection*{Máximo común divisor y mínimo común múltiplo}

Sean \(a,b\in\mathbb{N}\) dos enteros positivos. Supongamos que, usando todos los primos que aparecen en alguno de ellos, podemos escribir
\[
  a=p_1^{\alpha_1}\cdots p_k^{\alpha_k},
  \qquad
  b=p_1^{\beta_1}\cdots p_k^{\beta_k},
\]
donde algunos exponentes pueden ser cero.

Entonces
\[
  \mcd(a,b)=p_1^{\min(\alpha_1,\beta_1)}\cdots p_k^{\min(\alpha_k,\beta_k)},
\]
y
\[
  \operatorname{mcm}(a,b)=p_1^{\max(\alpha_1,\beta_1)}\cdots p_k^{\max(\alpha_k,\beta_k)}.
\]

\begin{ejemplo}
Sean
\[
  72=2^3\cdot 3^2,
  \qquad
  180=2^2\cdot 3^2\cdot 5.
\]
Entonces
\[
  \mcd(72,180)=2^2\cdot 3^2=36,
\]
y
\[
  \operatorname{mcm}(72,180)=2^3\cdot 3^2\cdot 5=360.
\]
\end{ejemplo}

\begin{advertencia}
El máximo común divisor se obtiene tomando los exponentes mínimos de los primos comunes, mientras que el mínimo común múltiplo se obtiene tomando los exponentes máximos de todos los primos que aparecen. Confundir estos dos procedimientos es un error frecuente.
\end{advertencia}

\section{Irracionalidad de ciertas raíces}

La factorización única permite demostrar resultados clásicos sobre números irracionales. Recordemos que un número real \(x\) es racional si existen enteros \(a,b\in\mathbb{Z}\), con \(b\neq 0\), tales que
\[
  x=\frac{a}{b}.
\]
Si un número real no es racional, se dice que es irracional.

\begin{teorema}[Irracionalidad de \(\sqrt{2}\)]
El número \(\sqrt{2}\) es irracional.
\end{teorema}

\begin{proof}
Supongamos, para llegar a una contradicción, que \(\sqrt{2}\) es racional. Entonces existen enteros positivos \(a,b\in\mathbb{N}\), sin factores comunes mayores que \(1\), tales que
\[
  \sqrt{2}=\frac{a}{b}.
\]
Elevando al cuadrado, obtenemos
\[
  2=\frac{a^2}{b^2},
\]
y, por tanto,
\[
  a^2=2b^2.
\]
Esto implica que \(a^2\) es par. Por la factorización única, si \(a^2\) es divisible por \(2\), entonces \(a\) también es divisible por \(2\). Luego existe \(c\in\mathbb{N}\) tal que
\[
  a=2c.
\]
Sustituyendo en la igualdad anterior,
\[
  (2c)^2=2b^2.
\]
Por tanto,
\[
  4c^2=2b^2,
\]
y dividiendo entre \(2\), obtenemos
\[
  b^2=2c^2.
\]
Así, \(b^2\) es par, y de nuevo \(b\) es par.

Hemos demostrado que tanto \(a\) como \(b\) son pares. Esto contradice que \(a\) y \(b\) no tengan factores comunes mayores que \(1\). Por tanto, \(\sqrt{2}\) no puede ser racional.
\end{proof}

\begin{proposicion}
Sea \(p\) un número primo. Entonces \(\sqrt{p}\) es irracional.
\end{proposicion}

\begin{proof}
Supongamos que \(\sqrt{p}=a/b\), donde \(a,b\in\mathbb{N}\) no tienen factores comunes mayores que \(1\). Entonces
\[
  p=\frac{a^2}{b^2},
\]
y por tanto
\[
  a^2=pb^2.
\]
Así, \(p\mid a^2\). Por el Teorema Fundamental de la Aritmética, esto implica que \(p\mid a\). Luego existe \(c\in\mathbb{N}\) tal que \(a=pc\). Sustituyendo,
\[
  p^2c^2=pb^2.
\]
Dividiendo entre \(p\), obtenemos
\[
  pc^2=b^2.
\]
Por tanto, \(p\mid b^2\), y de nuevo \(p\mid b\). Esto contradice que \(a\) y \(b\) no tengan factores comunes mayores que \(1\). Concluimos que \(\sqrt{p}\) es irracional.
\end{proof}

\begin{erroresfrecuentes}
  \item Considerar que \(1\) es primo. Esta convención destruye la unicidad de la factorización.
  \item Confundir existencia con unicidad: saber que una factorización existe no significa todavía que sea única.
  \item Olvidar que la unicidad de la factorización es salvo el orden de los factores.
  \item Calcular el máximo común divisor usando exponentes máximos en lugar de mínimos.
  \item Calcular el mínimo común múltiplo usando solo los primos comunes.
  \item En las demostraciones de irracionalidad, no justificar que la fracción se toma en forma irreducible.
\end{erroresfrecuentes}

\begin{seminario}
  \item Descomposición de enteros positivos en factores primos.
  \item Aplicaciones a problemas de divisibilidad.
  \item Demostraciones clásicas de irracionalidad.
\end{seminario}

\begin{ejerciciospropuestos}
  \item Descompón en factores primos los números \(84\), \(126\), \(360\) y \(2025\).
  \item Calcula el máximo común divisor y el mínimo común múltiplo de \(84\) y \(126\) usando factorización en primos.
  \item Demuestra que si \(p\) es primo y \(p\mid a^2\), entonces \(p\mid a\).
  \item Demuestra que \(\sqrt{3}\) es irracional.
  \item Determina todos los divisores positivos de \(72\).
  \item Si \(n=2^3\cdot 3^2\cdot 5\), calcula cuántos divisores positivos tiene \(n\).
\end{ejerciciospropuestos}

\resumen{El Teorema Fundamental de la Aritmética afirma que todo entero positivo mayor que uno admite una factorización en primos y que dicha factorización es única salvo el orden de los factores. Este resultado permite estudiar divisibilidad, divisores, máximo común divisor, mínimo común múltiplo e irracionalidad de ciertas raíces mediante la estructura de los factores primos.}


% ------------------------------------------------------------

\chapter{Aritmética modular}

\begin{objetivos}
  \item Introducir la congruencia módulo un entero positivo.
  \item Definir las clases de congruencia.
  \item Construir operaciones en el conjunto cociente.
  \item Resolver ecuaciones lineales módulo un entero.
\end{objetivos}

\section{Congruencias}

En este capítulo se estudia una forma de aritmética en la que los números enteros se comparan no por su valor exacto, sino por el resto que dejan al dividirlos por un entero positivo fijo.

\begin{definicion}[Módulo]
Sea \(n\) un número entero positivo. Diremos que \(n\) es el \emph{módulo} de la aritmética considerada.
\end{definicion}

La idea fundamental es identificar enteros que tienen el mismo resto al dividirse por \(n\). Para expresar esta idea se introduce la relación de congruencia.

\begin{definicion}[Congruencia módulo un entero]
Sea \(n\in\mathbb{N}\) con \(n\geq 1\), y sean \(a,b\in\mathbb{Z}\). Diremos que \(a\) es \emph{congruente} con \(b\) módulo \(n\), y escribiremos
\[
  a \equiv b \pmod n,
\]
si \(n\) divide a \(a-b\), es decir, si existe \(k\in\mathbb{Z}\) tal que
\[
  a-b = kn.
\]
\end{definicion}

En esta definición todos los símbolos tienen un papel concreto: \(n\) es el módulo, \(a\) y \(b\) son enteros, y el símbolo \(\equiv\) expresa una relación entre enteros dependiente de \(n\).

\begin{ejemplo}
Se tiene
\[
  17 \equiv 5 \pmod 6,
\]
porque
\[
  17-5=12=2\cdot 6.
\]
También se tiene
\[
  -1 \equiv 4 \pmod 5,
\]
porque
\[
  -1-4=-5=(-1)\cdot 5.
\]
\end{ejemplo}

\begin{proposicion}[Caracterización por restos]
Sean \(n\in\mathbb{N}\), con \(n\geq 1\), y sean \(a,b\in\mathbb{Z}\). Entonces
\[
  a \equiv b \pmod n
\]
si y solo si \(a\) y \(b\) tienen el mismo resto al dividirse por \(n\).
\end{proposicion}

\begin{proof}
Por la división euclídea, existen enteros \(q,r,q',r'\) tales que
\[
  a=qn+r, \qquad b=q'n+r',
\]
con
\[
  0\leq r<n, \qquad 0\leq r'<n.
\]
Entonces
\[
  a-b=(q-q')n+(r-r').
\]
Si \(a\equiv b\pmod n\), entonces \(n\mid(a-b)\). Como \((q-q')n\) es múltiplo de \(n\), también debe serlo \(r-r'\). Pero \(r-r'\) está comprendido entre \(-(n-1)\) y \(n-1\). El único múltiplo de \(n\) en ese intervalo es \(0\). Por tanto, \(r-r'=0\), es decir, \(r=r'\).

Recíprocamente, si \(r=r'\), entonces
\[
  a-b=(q-q')n,
\]
por lo que \(n\mid(a-b)\). Así, \(a\equiv b\pmod n\).
\end{proof}

\begin{advertencia}
La congruencia módulo \(n\) no significa igualdad ordinaria. Por ejemplo, \(17\equiv 5\pmod 6\), pero \(17\neq 5\). La congruencia expresa igualdad de restos, no igualdad de enteros.
\end{advertencia}

\section{Clases de congruencia}

La congruencia módulo \(n\) es una relación de equivalencia sobre \(\mathbb{Z}\). Por tanto, permite agrupar los enteros en clases.

\begin{proposicion}
Sea \(n\in\mathbb{N}\), con \(n\geq 1\). La relación sobre \(\mathbb{Z}\) definida por
\[
  a\sim b \quad \Longleftrightarrow \quad a\equiv b\pmod n
\]
es una relación de equivalencia.
\end{proposicion}

\begin{proof}
Veamos las tres propiedades.

En primer lugar, para todo \(a\in\mathbb{Z}\), se tiene
\[
  a-a=0=0\cdot n,
\]
por lo que \(a\equiv a\pmod n\). La relación es reflexiva.

En segundo lugar, si \(a\equiv b\pmod n\), entonces existe \(k\in\mathbb{Z}\) tal que
\[
  a-b=kn.
\]
Entonces
\[
  b-a=(-k)n,
\]
por lo que \(b\equiv a\pmod n\). La relación es simétrica.

Por último, si \(a\equiv b\pmod n\) y \(b\equiv c\pmod n\), existen \(k,l\in\mathbb{Z}\) tales que
\[
  a-b=kn, \qquad b-c=ln.
\]
Sumando ambas igualdades,
\[
  a-c=(k+l)n.
\]
Así, \(a\equiv c\pmod n\). La relación es transitiva.
\end{proof}

\begin{definicion}[Clase de congruencia]
Sea \(n\in\mathbb{N}\), con \(n\geq 1\), y sea \(a\in\mathbb{Z}\). La \emph{clase de congruencia} de \(a\) módulo \(n\) es el conjunto
\[
  [a]_n = \{x\in\mathbb{Z}: x\equiv a\pmod n\}.
\]
Cuando el módulo \(n\) se sobreentiende, se puede escribir simplemente \([a]\).
\end{definicion}

\begin{ejemplo}
Módulo \(5\), la clase de \(2\) es
\[
  {[2]}_5=\{\ldots,-8,-3,2,7,12,17,\ldots\}.
\]
Todos estos enteros dejan resto \(2\) al dividirse por \(5\), si se toma el resto entre \(0\) y \(4\).
\end{ejemplo}

\begin{proposicion}
Sea \(n\in\mathbb{N}\), con \(n\geq 1\). Todo entero es congruente módulo \(n\) con exactamente uno de los enteros
\[
  0,1,2,\ldots,n-1.
\]
\end{proposicion}

\begin{proof}
Sea \(a\in\mathbb{Z}\). Por la división euclídea, existen \(q,r\in\mathbb{Z}\) tales que
\[
  a=qn+r, \qquad 0\leq r<n.
\]
Entonces
\[
  a-r=qn,
\]
por lo que \(a\equiv r\pmod n\). Esto prueba la existencia.

Para la unicidad, si \(a\equiv r\pmod n\) y \(a\equiv s\pmod n\), con
\[
  0\leq r<n, \qquad 0\leq s<n,
\]
entonces \(r\equiv s\pmod n\). Por tanto, \(n\mid(r-s)\). Como \(r-s\) está entre \(-(n-1)\) y \(n-1\), debe ser \(r-s=0\). Luego \(r=s\).
\end{proof}

\section{El conjunto cociente módulo un entero}

El conjunto formado por todas las clases de congruencia módulo \(n\) se denomina conjunto cociente de \(\mathbb{Z}\) módulo \(n\).

\begin{definicion}[Conjunto cociente módulo \(n\)]
Sea \(n\in\mathbb{N}\), con \(n\geq 1\). El conjunto cociente de \(\mathbb{Z}\) módulo \(n\) se denota por
\[
  \mathbb{Z}/n\mathbb{Z}
\]
y se define como
\[
  \mathbb{Z}/n\mathbb{Z}=\{{[0]}_n,{[1]}_n,\ldots,[n-1]_n\}.
\]
\end{definicion}

\begin{ejemplo}
Para \(n=4\), se tiene
\[
  \mathbb{Z}/4\mathbb{Z}=\{{[0]}_4,{[1]}_4,{[2]}_4,[3]_4\}.
\]
La clase \({[0]}_4\) contiene los múltiplos de \(4\), la clase \({[1]}_4\) contiene los enteros de la forma \(4q+1\), la clase \({[2]}_4\) contiene los enteros de la forma \(4q+2\), y la clase \([3]_4\) contiene los enteros de la forma \(4q+3\), donde \(q\in\mathbb{Z}\).
\end{ejemplo}

\begin{advertencia}
El símbolo \([a]_n\) representa una clase, es decir, un conjunto de enteros. No representa únicamente al entero \(a\). Por ejemplo, \({[2]}_5=[7]_5=[-3]_5\), aunque \(2\), \(7\) y \(-3\) sean enteros distintos.
\end{advertencia}

\section{Operaciones módulo un entero}

Las operaciones de suma y producto de enteros inducen operaciones entre clases de congruencia.

\begin{definicion}[Suma y producto módulo \(n\)]
Sea \(n\in\mathbb{N}\), con \(n\geq 1\). Si \([a]_n,[b]_n\in\mathbb{Z}/n\mathbb{Z}\), definimos
\[
  [a]_n+[b]_n=[a+b]_n
\]
y
\[
  [a]_n\cdot[b]_n=[ab]_n.
\]
\end{definicion}

Para que estas definiciones sean correctas, hay que comprobar que no dependen del representante elegido. Es decir, si \([a]_n=[a']_n\) y \([b]_n=[b']_n\), entonces debe cumplirse
\[
  [a+b]_n=[a'+b']_n
\]
y
\[
  [ab]_n=[a'b']_n.
\]

\begin{proposicion}[Buena definición de las operaciones]
Las operaciones de suma y producto en \(\mathbb{Z}/n\mathbb{Z}\) están bien definidas.
\end{proposicion}

\begin{proof}
Supongamos que \(a\equiv a'\pmod n\) y \(b\equiv b'\pmod n\). Entonces existen \(k,l\in\mathbb{Z}\) tales que
\[
  a-a'=kn, \qquad b-b'=ln.
\]
Sumando,
\[
  (a+b)-(a'+b')=(a-a')+(b-b')=(k+l)n,
\]
por lo que
\[
  a+b\equiv a'+b'\pmod n.
\]
Así, la suma está bien definida.

Para el producto, observamos que
\[
  ab-a'b'=ab-a'b+a'b-a'b'=b(a-a')+a'(b-b').
\]
Sustituyendo las igualdades anteriores,
\[
  ab-a'b'=bkn+a'ln=(bk+a'l)n.
\]
Por tanto,
\[
  ab\equiv a'b'\pmod n.
\]
Así, el producto también está bien definido.
\end{proof}

\begin{ejemplo}
En \(\mathbb{Z}/5\mathbb{Z}\), se tiene
\[
  [3]_5+[4]_5=[7]_5={[2]}_5,
\]
y
\[
  [3]_5\cdot[4]_5=[12]_5={[2]}_5.
\]
\end{ejemplo}

\begin{ejemplo}[Tabla de suma módulo \(3\)]
En \(\mathbb{Z}/3\mathbb{Z}\), la suma queda descrita por la tabla
\[
\begin{array}{c|ccc}
+ & {[0]} & {[1]} & {[2]} \\
\hline
{[0]} & {[0]} & {[1]} & {[2]} \\
{[1]} & {[1]} & {[2]} & {[0]} \\
{[2]} & {[2]} & {[0]} & {[1]}
\end{array}
\]
En la tabla se ha omitido el subíndice \(3\) para aligerar la notación.
\end{ejemplo}

\begin{ejemplo}[Tabla de producto módulo \(3\)]
En \(\mathbb{Z}/3\mathbb{Z}\), el producto queda descrito por la tabla
\[
\begin{array}{c|ccc}
\cdot & {[0]} & {[1]} & {[2]} \\
\hline
{[0]} & {[0]} & {[0]} & {[0]} \\
{[1]} & {[0]} & {[1]} & {[2]} \\
{[2]} & {[0]} & {[2]} & {[1]}
\end{array}
\]
\end{ejemplo}

\section{Inversos modulares}

En aritmética ordinaria, un número entero distinto de \(1\) y \(-1\) no suele tener inverso multiplicativo dentro de \(\mathbb{Z}\). Sin embargo, módulo \(n\), algunas clases sí pueden tener inverso multiplicativo.

\begin{definicion}[Inverso modular]
Sea \(n\in\mathbb{N}\), con \(n\geq 2\), y sea \(a\in\mathbb{Z}\). Diremos que la clase \([a]_n\) tiene \emph{inverso multiplicativo} en \(\mathbb{Z}/n\mathbb{Z}\) si existe \([b]_n\in\mathbb{Z}/n\mathbb{Z}\) tal que
\[
  [a]_n[b]_n={[1]}_n.
\]
En ese caso, \([b]_n\) se llama un \emph{inverso modular} de \([a]_n\).
\end{definicion}

Equivalente, \([b]_n\) es inverso de \([a]_n\) si
\[
  ab\equiv 1\pmod n.
\]

\begin{teorema}[Criterio de invertibilidad modular]
Sea \(n\in\mathbb{N}\), con \(n\geq 2\), y sea \(a\in\mathbb{Z}\). La clase \([a]_n\) tiene inverso multiplicativo en \(\mathbb{Z}/n\mathbb{Z}\) si y solo si
\[
  \mcd(a,n)=1.
\]
\end{teorema}

\begin{proof}
Supongamos primero que \([a]_n\) tiene inverso. Entonces existe \(b\in\mathbb{Z}\) tal que
\[
  ab\equiv 1\pmod n.
\]
Esto significa que existe \(k\in\mathbb{Z}\) tal que
\[
  ab-1=kn.
\]
Por tanto,
\[
  ab-kn=1.
\]
Si un entero \(d\) divide a \(a\) y a \(n\), entonces divide a \(ab\) y a \(kn\), y por tanto divide a \(ab-kn=1\). Así, \(d=1\) o \(d=-1\). En consecuencia, el máximo común divisor positivo de \(a\) y \(n\) es \(1\).

Recíprocamente, si \(\mcd(a,n)=1\), por la identidad de Bézout existen \(u,v\in\mathbb{Z}\) tales que
\[
  au+nv=1.
\]
Entonces
\[
  au-1=-nv,
\]
por lo que
\[
  au\equiv 1\pmod n.
\]
Así, \([u]_n\) es inverso multiplicativo de \([a]_n\).
\end{proof}

\begin{ejemplo}
En \(\mathbb{Z}/7\mathbb{Z}\), la clase \([3]_7\) tiene inverso porque \(\mcd(3,7)=1\). De hecho,
\[
  3\cdot 5=15\equiv 1\pmod 7.
\]
Por tanto, \([5]_7\) es el inverso de \([3]_7\).
\end{ejemplo}

\begin{contraejemplo}
En \(\mathbb{Z}/8\mathbb{Z}\), la clase \({[2]}_8\) no tiene inverso multiplicativo, porque
\[
  \mcd(2,8)=2\neq 1.
\]
En efecto, los productos \(2b\) son siempre pares, y por tanto no pueden ser congruentes con \(1\) módulo \(8\).
\end{contraejemplo}

\section{Ecuaciones lineales módulo un entero}

Una ecuación lineal módulo \(n\) tiene la forma
\[
  ax\equiv b\pmod n,
\]
donde \(a,b\in\mathbb{Z}\), \(n\in\mathbb{N}\) con \(n\geq 2\), y la incógnita es \(x\in\mathbb{Z}\).

\begin{teorema}[Criterio de resolución]
Sean \(a,b\in\mathbb{Z}\) y sea \(n\in\mathbb{N}\), con \(n\geq 2\). La congruencia lineal
\[
  ax\equiv b\pmod n
\]
tiene solución si y solo si
\[
  \mcd(a,n)\mid b.
\]
\end{teorema}

\begin{proof}
La congruencia
\[
  ax\equiv b\pmod n
\]
es equivalente a decir que existe \(y\in\mathbb{Z}\) tal que
\[
  ax-b=ny.
\]
Reordenando,
\[
  ax-ny=b.
\]
Esta es una ecuación diofántica lineal. Por el criterio estudiado en el capítulo de divisibilidad, la ecuación
\[
  ax+nc=b
\]
tiene solución entera si y solo si \(\mcd(a,n)\mid b\). Como cambiar \(c\) por \(-y\) no altera el criterio, se obtiene el resultado.
\end{proof}

\begin{ejemplo}
Resolvamos
\[
  3x\equiv 4\pmod 7.
\]
Como \(\mcd(3,7)=1\), la clase \([3]_7\) tiene inverso. Sabemos que
\[
  3\cdot 5=15\equiv 1\pmod 7.
\]
Multiplicando ambos lados de la congruencia por \(5\), obtenemos
\[
  x\equiv 5\cdot 4=20\equiv 6\pmod 7.
\]
Por tanto, las soluciones son los enteros congruentes con \(6\) módulo \(7\):
\[
  x=6+7k, \qquad k\in\mathbb{Z}.
\]
\end{ejemplo}

\begin{ejemplo}
La congruencia
\[
  6x\equiv 9\pmod{15}
\]
tiene solución porque
\[
  \mcd(6,15)=3
\]
y \(3\mid 9\). Dividiendo toda la congruencia por \(3\), se obtiene
\[
  2x\equiv 3\pmod 5.
\]
Como \(2^{-1}\equiv 3\pmod 5\), multiplicamos por \(3\) y resulta
\[
  x\equiv 9\equiv 4\pmod 5.
\]
Por tanto, las soluciones módulo \(15\) son
\[
  x\equiv 4,9,14\pmod{15}.
\]
\end{ejemplo}

\begin{advertencia}
No siempre se puede dividir directamente una congruencia por un entero. La división solo es automática cuando el factor por el que se quiere dividir es invertible módulo \(n\). En caso contrario, hay que usar el criterio basado en el máximo común divisor.
\end{advertencia}

\section{Teorema Chino del Resto}

El Teorema Chino del Resto permite resolver sistemas de congruencias con módulos coprimos dos a dos.

\begin{definicion}[Módulos coprimos dos a dos]
Sean \(n_1,n_2,\ldots,n_r\in\mathbb{N}\), con \(n_i\geq 2\) para todo \(i\). Diremos que estos módulos son \emph{coprimos dos a dos} si
\[
  \mcd(n_i,n_j)=1
\]
siempre que \(i\neq j\).
\end{definicion}

\begin{teorema}[Teorema Chino del Resto]
Sean \(n_1,n_2,\ldots,n_r\in\mathbb{N}\), con \(n_i\geq 2\), coprimos dos a dos. Sean \(a_1,a_2,\ldots,a_r\in\mathbb{Z}\). Entonces el sistema
\[
\begin{cases}
  x\equiv a_1 \pmod{n_1},\\
  x\equiv a_2 \pmod{n_2},\\
  \quad\vdots\\
  x\equiv a_r \pmod{n_r}
\end{cases}
\]
tiene solución. Además, la solución es única módulo
\[
  N=n_1n_2\cdots n_r.
\]
\end{teorema}

\begin{proof}
Daremos la construcción de una solución. Definimos
\[
  N=n_1n_2\cdots n_r
\]
y, para cada \(i\in\{1,2,\ldots,r\}\), definimos
\[
  N_i=\frac{N}{n_i}.
\]
Como los módulos son coprimos dos a dos, se tiene
\[
  \mcd(N_i,n_i)=1.
\]
Por tanto, \(N_i\) tiene inverso módulo \(n_i\). Sea \(u_i\in\mathbb{Z}\) tal que
\[
  N_i u_i\equiv 1\pmod{n_i}.
\]
Consideramos
\[
  x=a_1N_1u_1+a_2N_2u_2+\cdots+a_rN_ru_r.
\]
Si reducimos módulo \(n_i\), todos los términos con índice \(j\neq i\) son congruentes con \(0\), porque \(n_i\mid N_j\). El término con índice \(i\) satisface
\[
  a_iN_iu_i\equiv a_i\pmod{n_i}.
\]
Luego
\[
  x\equiv a_i\pmod{n_i}
\]
para todo \(i\). Esto prueba la existencia.

Para la unicidad módulo \(N\), si \(x\) e \(y\) son dos soluciones, entonces
\[
  x\equiv y\pmod{n_i}
\]
para todo \(i\). Por tanto, cada \(n_i\) divide a \(x-y\). Como los \(n_i\) son coprimos dos a dos, su producto \(N\) divide a \(x-y\). Así,
\[
  x\equiv y\pmod N.
\]
\end{proof}

\begin{ejemplo}
Resolvamos el sistema
\[
\begin{cases}
  x\equiv 2 \pmod 3,\\
  x\equiv 3 \pmod 5.
\end{cases}
\]
Los módulos \(3\) y \(5\) son coprimos. Buscamos enteros congruentes con \(2\) módulo \(3\):
\[
  2,5,8,11,14,17,\ldots
\]
Entre ellos, \(8\) es congruente con \(3\) módulo \(5\). Por tanto,
\[
  x\equiv 8\pmod{15}.
\]
\end{ejemplo}

\begin{erroresfrecuentes}
  \item Confundir la congruencia \(a\equiv b\pmod n\) con la igualdad ordinaria \(a=b\).
  \item Olvidar que \([a]_n\) es una clase de enteros, no solo el entero \(a\).
  \item Dividir una congruencia por un número que no tiene inverso módulo \(n\).
  \item Suponer que toda clase no nula tiene inverso en \(\mathbb{Z}/n\mathbb{Z}\). Esto solo ocurre cuando \(n\) es primo.
  \item Aplicar el Teorema Chino del Resto sin comprobar que los módulos son coprimos dos a dos.
\end{erroresfrecuentes}

\begin{seminario}
  \item Cálculo con congruencias.
  \item Tablas de suma y producto módulo un entero.
  \item Resolución de congruencias lineales.
\end{seminario}

\begin{ejerciciospropuestos}
  \item Decide si las siguientes congruencias son verdaderas:
  \[
    25\equiv 4\pmod 7,\qquad -3\equiv 9\pmod 4,\qquad 18\equiv 0\pmod 6.
  \]
  \item Calcula las clases de congruencia módulo \(6\) y escribe tres representantes distintos de cada una.
  \item Construye las tablas de suma y producto de \(\mathbb{Z}/4\mathbb{Z}\).
  \item Determina cuáles de las clases \({[0]}_8,{[1]}_8,\ldots,[7]_8\) tienen inverso multiplicativo módulo \(8\).
  \item Resuelve la congruencia
  \[
    5x\equiv 3\pmod{11}.
  \]
  \item Decide si la congruencia
  \[
    6x\equiv 10\pmod{14}
  \]
  tiene solución. En caso afirmativo, encuéntrala.
  \item Resuelve el sistema
  \[
  \begin{cases}
    x\equiv 1\pmod 4,\\
    x\equiv 2\pmod 5.
  \end{cases}
  \]
\end{ejerciciospropuestos}

\resumen{La aritmética modular estudia los enteros agrupados según sus restos al dividir por un módulo fijo. La congruencia módulo \(n\) es una relación de equivalencia sobre \(\mathbb{Z}\), cuyas clases forman el conjunto cociente \(\mathbb{Z}/n\mathbb{Z}\). En este conjunto se pueden definir suma y producto de clases, siempre cuidando que las operaciones estén bien definidas. Los inversos modulares existen exactamente para las clases representadas por enteros coprimos con el módulo. Las ecuaciones lineales módulo \(n\) se resuelven mediante el máximo común divisor, y el Teorema Chino del Resto permite resolver sistemas de congruencias con módulos coprimos dos a dos.}


% ------------------------------------------------------------

\chapter{Los números racionales}

\begin{objetivos}
  \item Construir los números racionales a partir de los enteros.
  \item Interpretar las fracciones como clases de equivalencia.
  \item Estudiar operaciones, orden y densidad de los racionales.
\end{objetivos}

\section{Motivación: división y fracciones}

En el capítulo anterior se estudió el conjunto de los números enteros, que denotamos por \(\mathbb{Z}\). Este conjunto permite sumar, restar y multiplicar sin salir de él. Sin embargo, no siempre permite dividir.

Por ejemplo, la ecuación
\[
  2x=1
\]
no tiene solución en \(\mathbb{Z}\), porque no existe ningún entero \(x\) cuyo doble sea igual a \(1\). Para resolver ecuaciones de este tipo es necesario ampliar el sistema numérico.

La idea fundamental consiste en introducir objetos que representen cocientes de enteros. Estos objetos reciben el nombre de \emph{números racionales}.

\begin{convencion}[Símbolos básicos]
A lo largo de este capítulo, el símbolo \(\mathbb{Z}\) denota el conjunto de los números enteros. El símbolo \(\mathbb{Q}\), que se definirá más adelante, denotará el conjunto de los números racionales. Todo símbolo que aparezca en una definición será introducido antes de su uso.
\end{convencion}

La expresión \(\frac{a}{b}\), donde \(a\) y \(b\) son enteros y \(b\neq 0\), se interpreta inicialmente como una fracción. El entero \(a\) se llama \emph{numerador} y el entero \(b\) se llama \emph{denominador}.

\begin{advertencia}
La expresión \(\frac{a}{b}\) solo tiene sentido como fracción si el denominador \(b\) es distinto de cero. La división por cero no está definida.
\end{advertencia}

\section{Fracciones y clases de equivalencia}

Una dificultad esencial es que una misma cantidad puede representarse mediante fracciones diferentes. Por ejemplo,
\[
  \frac{1}{2},\qquad \frac{2}{4},\qquad \frac{3}{6}
\]
representan el mismo número racional.

Por tanto, no conviene definir un número racional como una fracción concreta, sino como una clase de fracciones equivalentes.

\begin{definicion}[Conjunto de pares admisibles]
Denotamos por
\[
  S=\{(a,b)\in \mathbb{Z}\times \mathbb{Z}: b\neq 0\}
\]
el conjunto de todos los pares ordenados de enteros cuyo segundo componente es no nulo.
\end{definicion}

Un par \((a,b)\in S\) representa informalmente la fracción \(\frac{a}{b}\).

\begin{definicion}[Equivalencia de fracciones]
Sean \((a,b),(c,d)\in S\). Decimos que \((a,b)\) es equivalente a \((c,d)\), y escribimos
\[
  (a,b)\sim(c,d),
\]
si
\[
  ad=bc.
\]
\end{definicion}

La condición \(ad=bc\) expresa la igualdad usual de fracciones:
\[
  \frac{a}{b}=\frac{c}{d}
  \quad \Longleftrightarrow \quad
  ad=bc.
\]

\begin{proposicion}
La relación \(\sim\) definida sobre \(S\) es una relación de equivalencia.
\end{proposicion}

\begin{proof}
Debemos comprobar que \(\sim\) es reflexiva, simétrica y transitiva.

Sea \((a,b)\in S\). Como \(ab=ba\), se tiene \((a,b)\sim(a,b)\). Por tanto, \(\sim\) es reflexiva.

Supongamos ahora que \((a,b)\sim(c,d)\). Entonces \(ad=bc\). La misma igualdad puede escribirse como \(cb=da\), luego \((c,d)\sim(a,b)\). Por tanto, \(\sim\) es simétrica.

Finalmente, supongamos que \((a,b)\sim(c,d)\) y \((c,d)\sim(e,f)\). Entonces
\[
  ad=bc,
  \qquad
  cf=de.
\]
Multiplicando la primera igualdad por \(f\) y la segunda por \(b\), obtenemos
\[
  adf=bcf,
  \qquad
  bcf=bde.
\]
Por tanto, \(adf=bde\). Como \(d\neq 0\), podemos cancelar \(d\) en \(\mathbb{Z}\) y obtenemos
\[
  af=be.
\]
Así, \((a,b)\sim(e,f)\). Por tanto, \(\sim\) es transitiva.
\end{proof}

\begin{definicion}[Clase de una fracción]
Sea \((a,b)\in S\). La clase de equivalencia de \((a,b)\) es el conjunto
\[
  [(a,b)]=\{(c,d)\in S:(c,d)\sim(a,b)\}.
\]
Esta clase contiene todas las fracciones equivalentes a \(\frac{a}{b}\).
\end{definicion}

Por simplicidad de notación, la clase \([(a,b)]\) se suele escribir como
\[
  \frac{a}{b}.
\]
Esta notación debe entenderse como una clase de equivalencia, no como un único par ordenado.

\begin{ejemplo}
Las fracciones \(\frac{1}{2}\), \(\frac{2}{4}\) y \(\frac{-3}{-6}\) representan el mismo número racional, porque
\[
  1\cdot 4=2\cdot 2,
  \qquad
  1\cdot(-6)=2\cdot(-3).
\]
Por tanto,
\[
  \frac{1}{2}=\frac{2}{4}=\frac{-3}{-6}.
\]
\end{ejemplo}

\section{El conjunto de los números racionales}

Ya podemos definir el conjunto de los números racionales.

\begin{definicion}[Números racionales]
El conjunto de los números racionales, denotado por \(\mathbb{Q}\), es el conjunto de clases de equivalencia de pares \((a,b)\in S\) respecto de la relación \(\sim\). Es decir,
\[
  \mathbb{Q}=S/{\sim}.
\]
Sus elementos se escriben en la forma
\[
  \frac{a}{b},
  \qquad a,b\in\mathbb{Z},\quad b\neq 0.
\]
\end{definicion}

Todo número entero puede verse como un número racional. En efecto, si \(n\in\mathbb{Z}\), entonces
\[
  n=\frac{n}{1}.
\]
De este modo, podemos considerar \(\mathbb{Z}\) como subconjunto de \(\mathbb{Q}\).

\begin{convencion}[Enteros dentro de los racionales]
Cuando escribamos \(n\in\mathbb{Z}\) dentro de un contexto racional, identificaremos \(n\) con el racional \(\frac{n}{1}\).
\end{convencion}

\section{Operaciones con racionales}

Para que \(\mathbb{Q}\) sea un sistema numérico útil, debemos definir suma y producto.

\begin{definicion}[Suma de racionales]
Sean \(\frac{a}{b},\frac{c}{d}\in\mathbb{Q}\), con \(b\neq 0\) y \(d\neq 0\). Definimos su suma por
\[
  \frac{a}{b}+\frac{c}{d}
  =
  \frac{ad+bc}{bd}.
\]
\end{definicion}

\begin{definicion}[Producto de racionales]
Sean \(\frac{a}{b},\frac{c}{d}\in\mathbb{Q}\), con \(b\neq 0\) y \(d\neq 0\). Definimos su producto por
\[
  \frac{a}{b}\cdot\frac{c}{d}
  =
  \frac{ac}{bd}.
\]
\end{definicion}

Estas operaciones están bien definidas: si se cambia una fracción por otra equivalente, el resultado obtenido representa la misma clase racional.

\begin{ejemplo}
Calculamos
\[
  \frac{2}{3}+\frac{5}{4}
  =
  \frac{2\cdot 4+3\cdot 5}{3\cdot 4}
  =
  \frac{8+15}{12}
  =
  \frac{23}{12}.
\]
También,
\[
  \frac{2}{3}\cdot\frac{5}{4}
  =
  \frac{10}{12}
  =
  \frac{5}{6}.
\]
\end{ejemplo}

\begin{proposicion}
El conjunto \(\mathbb{Q}\), con la suma y el producto anteriores, es un cuerpo.
\end{proposicion}

\begin{proof}
La verificación completa consiste en comprobar las propiedades asociativa y conmutativa de la suma y del producto, la existencia de neutro aditivo \(0=\frac{0}{1}\), la existencia de neutro multiplicativo \(1=\frac{1}{1}\), la existencia de opuesto para cada racional y la existencia de inverso multiplicativo para cada racional no nulo. También debe comprobarse la distributividad del producto respecto de la suma.

Estas propiedades se deducen de las propiedades correspondientes de los enteros. En este curso se usará este resultado como una propiedad estructural básica de \(\mathbb{Q}\).
\end{proof}

\begin{advertencia}
La igualdad \(\frac{a}{b}=\frac{c}{d}\) no significa necesariamente que \(a=c\) y \(b=d\). Significa que \(ad=bc\). Por ejemplo, \(\frac{1}{2}=\frac{2}{4}\), aunque \(1\neq 2\) y \(2\neq 4\).
\end{advertencia}

\section{Orden en los racionales}

El orden de los enteros se extiende a los racionales.

\begin{definicion}[Racional positivo]
Sea \(\frac{a}{b}\in\mathbb{Q}\), con \(b\neq 0\). Decimos que \(\frac{a}{b}\) es positivo si \(a\) y \(b\) tienen el mismo signo, es decir, si \(ab>0\).
\end{definicion}

\begin{definicion}[Orden en \(\mathbb{Q}\)]
Sean \(r,s\in\mathbb{Q}\). Decimos que \(r<s\) si \(s-r\) es positivo. Decimos que \(r\leq s\) si \(r<s\) o \(r=s\).
\end{definicion}

En términos de fracciones, si \(b>0\) y \(d>0\), entonces
\[
  \frac{a}{b}<\frac{c}{d}
  \quad \Longleftrightarrow \quad
  ad<bc.
\]

\begin{ejemplo}
Comparamos \(\frac{2}{3}\) y \(\frac{3}{4}\). Como los denominadores son positivos, calculamos
\[
  2\cdot 4=8,
  \qquad
  3\cdot 3=9.
\]
Puesto que \(8<9\), se tiene
\[
  \frac{2}{3}<\frac{3}{4}.
\]
\end{ejemplo}

\begin{advertencia}
La regla de comparar productos cruzados debe aplicarse directamente solo cuando los denominadores son positivos. Si aparece un denominador negativo, conviene cambiar primero la fracción por otra equivalente con denominador positivo.
\end{advertencia}

\section{Densidad de los racionales}

Una propiedad esencial de \(\mathbb{Q}\) es que entre dos racionales distintos siempre hay otro racional.

\begin{teorema}[Densidad de \(\mathbb{Q}\)]
Sean \(r,s\in\mathbb{Q}\) tales que \(r<s\). Entonces existe \(q\in\mathbb{Q}\) tal que
\[
  r<q<s.
\]
\end{teorema}

\begin{proof}
Definimos
\[
  q=\frac{r+s}{2}.
\]
Como \(r\) y \(s\) son racionales, también \(r+s\) es racional, y como \(2\in\mathbb{Z}\) con \(2\neq 0\), el número \(q\) es racional.

Además, de \(r<s\) se deduce
\[
  2r<r+s<2s.
\]
Dividiendo entre \(2\), obtenemos
\[
  r<\frac{r+s}{2}<s.
\]
Por tanto, \(q\) es un racional situado estrictamente entre \(r\) y \(s\).
\end{proof}

\begin{corolario}
Entre dos números racionales distintos existen infinitos números racionales.
\end{corolario}

\begin{proof}
Sean \(r,s\in\mathbb{Q}\) con \(r<s\). Por el teorema anterior existe un racional \(q_1\) tal que \(r<q_1<s\). Aplicando de nuevo el teorema a \(r\) y \(q_1\), existe \(q_2\in\mathbb{Q}\) tal que \(r<q_2<q_1\). Repitiendo el proceso se obtienen indefinidamente racionales distintos entre \(r\) y \(s\).
\end{proof}

\section{Expansiones decimales}

Los números racionales también pueden representarse mediante desarrollos decimales.

\begin{definicion}[Expansión decimal]
Una expansión decimal es una escritura de un número mediante potencias de \(10\). Por ejemplo,
\[
  3.25=3+\frac{2}{10}+\frac{5}{100}.
\]
\end{definicion}

Al dividir dos enteros \(a\) y \(b\), con \(b\neq 0\), el algoritmo de la división produce restos. Como solo hay un número finito de posibles restos, o bien el proceso termina, o bien algún resto se repite.

\begin{proposicion}
Todo número racional tiene una expansión decimal finita o periódica.
\end{proposicion}

\begin{proof}
Sea \(\frac{a}{b}\in\mathbb{Q}\), con \(b>0\). Al realizar la división decimal de \(a\) entre \(b\), los restos posibles son
\[
  0,1,2,\ldots,b-1.
\]
Si en algún momento aparece el resto \(0\), la expansión decimal termina. Si no aparece el resto \(0\), como solo hay \(b-1\) restos no nulos posibles, alguno de ellos debe repetirse. Desde ese momento, las cifras decimales se repiten periódicamente.
\end{proof}

\begin{ejemplo}
El racional \(\frac{1}{4}\) tiene expansión decimal finita:
\[
  \frac{1}{4}=0.25.
\]
El racional \(\frac{1}{3}\) tiene expansión decimal periódica:
\[
  \frac{1}{3}=0.333\ldots
\]
\end{ejemplo}

\begin{erroresfrecuentes}
  \item Pensar que una fracción es el par \((a,b)\), en lugar de la clase de todas las fracciones equivalentes a ella.
  \item Olvidar que el denominador de una fracción no puede ser cero.
  \item Comparar fracciones con denominadores negativos sin transformarlas previamente.
  \item Confundir densidad con continuidad: que entre dos racionales haya otro racional no significa que los racionales llenen todos los puntos de la recta real.
  \item Creer que toda expansión decimal infinita representa necesariamente un número racional. Solo las expansiones decimales finitas o periódicas corresponden a racionales.
\end{erroresfrecuentes}

\begin{seminario}
  \item Operaciones con fracciones como clases de equivalencia.
  \item Problemas de orden y densidad.
  \item Comparación de racionales mediante productos cruzados.
  \item Construcción de racionales entre dos racionales dados.
  \item Análisis de expansiones decimales finitas y periódicas.
\end{seminario}

\begin{ejerciciospropuestos}
  \item Comprueba que \(\frac{6}{8}=\frac{3}{4}\) usando la definición de equivalencia de fracciones.
  \item Calcula \(\frac{2}{5}+\frac{3}{7}\) y \(\frac{2}{5}\cdot\frac{3}{7}\).
  \item Ordena de menor a mayor los racionales \(\frac{1}{2}\), \(\frac{3}{5}\), \(\frac{7}{10}\) y \(\frac{2}{3}\).
  \item Encuentra tres números racionales estrictamente comprendidos entre \(\frac{1}{3}\) y \(\frac{1}{2}\).
  \item Demuestra que si \(r,s\in\mathbb{Q}\) y \(r<s\), entonces \(\frac{2r+s}{3}\) es un racional situado entre \(r\) y \(s\).
  \item Decide si las expansiones decimales siguientes corresponden a números racionales: \(0.125\), \(0.272727\ldots\), \(0.10100100010000\ldots\).
\end{ejerciciospropuestos}

\resumen{El capítulo construye los números racionales como clases de equivalencia de fracciones de enteros. Esta construcción permite definir rigurosamente suma, producto y orden en \(\mathbb{Q}\). Los racionales forman un cuerpo ordenado y tienen la propiedad de densidad: entre dos racionales distintos siempre existe otro racional. Además, sus expansiones decimales son exactamente las finitas o periódicas.}


% ------------------------------------------------------------

\chapter{Los n\'umeros reales}

\begin{objetivos}
  \item Explicar por qu\'e los n\'umeros racionales no son suficientes.
  \item Introducir los n\'umeros reales como sistema num\'erico completo.
  \item Estudiar intervalos, valor absoluto y distancia.
  \item Presentar las nociones de supremo e \'infimo.
\end{objetivos}

\section{La insuficiencia de los racionales}

En el cap\'itulo anterior se introdujo el conjunto de los n\'umeros racionales, denotado por
\(\mathbb{Q}\). Sus elementos pueden escribirse en forma de fracci\'on
\[
  \frac{a}{b},
\]
donde \(a\in \mathbb{Z}\), \(b\in \mathbb{Z}\) y \(b\neq 0\). El conjunto \(\mathbb{Q}\) permite sumar, restar, multiplicar y dividir por n\'umeros no nulos. Adem\'as, est\'a ordenado y es denso: entre dos racionales distintos siempre existe otro racional.

Sin embargo, los racionales no bastan para resolver todos los problemas num\'ericos que aparecen de forma natural. Por ejemplo, la longitud de la diagonal de un cuadrado de lado \(1\) no puede expresarse mediante un n\'umero racional.

\begin{ejemplo}[La diagonal del cuadrado unidad]
Sea un cuadrado de lado \(1\). Por el teorema de Pit\'agoras, si \(d\) denota la longitud de su diagonal, entonces
\[
  d^2=1^2+1^2=2.
\]
Por tanto, \(d\) deber\'ia ser un n\'umero positivo cuyo cuadrado sea \(2\). Ese n\'umero se denota por \(\sqrt{2}\). Pero \(\sqrt{2}\notin\mathbb{Q}\).
\end{ejemplo}

La existencia de magnitudes como \(\sqrt{2}\) muestra que \(\mathbb{Q}\) tiene huecos. Es decir, aunque los racionales sean densos, no llenan toda la recta num\'erica.

\begin{advertencia}
La densidad de \(\mathbb{Q}\) no significa que \(\mathbb{Q}\) contenga todos los n\'umeros necesarios para medir longitudes. Significa solamente que entre dos racionales distintos siempre hay otro racional.
\end{advertencia}

\section{N\'umeros irracionales}

Un n\'umero irracional es un n\'umero que no puede escribirse como cociente de dos enteros con denominador no nulo.

\begin{definicion}[N\'umero irracional]
Un n\'umero real \(x\) se llama \emph{irracional} si no pertenece a \(\mathbb{Q}\). Es decir, \(x\) es irracional si no existen enteros \(a,b\in\mathbb{Z}\), con \(b\neq 0\), tales que
\[
  x=\frac{a}{b}.
\]
\end{definicion}

\begin{ejemplo}
Los n\'umeros
\[
  \sqrt{2},\qquad \sqrt{3},\qquad \pi,
\]
son ejemplos de n\'umeros irracionales.
\end{ejemplo}

En este curso no se desarrollar\'a una construcci\'on formal completa de todos los n\'umeros irracionales, pero s\'i se usar\'a la idea fundamental: los n\'umeros reales completan a los racionales a\~nadiendo los puntos que faltan en la recta num\'erica. Desde el punto de vista estructural, trabajaremos con \(\mathbb{R}\) como un cuerpo ordenado completo.

\section{El conjunto de los n\'umeros reales}

El conjunto de los n\'umeros reales se denota por \(\mathbb{R}\). Contiene a los n\'umeros racionales y tambi\'en a los n\'umeros irracionales.

\begin{definicion}[Conjunto de los n\'umeros reales]
El conjunto de los n\'umeros reales, denotado por \(\mathbb{R}\), es el sistema num\'erico que contiene a \(\mathbb{Q}\), incorpora los n\'umeros irracionales y permite representar de forma continua los puntos de la recta num\'erica.
\end{definicion}

Se tiene la inclusi\'on
\[
  \mathbb{Q}\subset \mathbb{R}.
\]
Esto significa que todo n\'umero racional es tambi\'en un n\'umero real. Sin embargo, no todo n\'umero real es racional.

\begin{ejemplo}
El n\'umero \(3/5\) es racional y, por tanto, real. El n\'umero \(\sqrt{2}\) no es racional, pero s\'i es real.
\end{ejemplo}

Los n\'umeros reales se representan geom\'etricamente mediante la recta real. Cada punto de la recta real corresponde a un n\'umero real, y cada n\'umero real corresponde a un punto de la recta.

\section{Orden y completitud}

El conjunto \(\mathbb{R}\) est\'a ordenado. Esto significa que, dados dos n\'umeros reales \(x\) e \(y\), tiene sentido comparar si
\[
  x<y,\qquad x=y,\qquad \text{o}\qquad x>y.
\]

\begin{definicion}[Orden en \(\mathbb{R}\)]
Sean \(x,y\in\mathbb{R}\). Escribimos \(x<y\) si \(x\) es menor que \(y\), y escribimos \(x\leq y\) si \(x<y\) o \(x=y\).
\end{definicion}

El orden de \(\mathbb{R}\) es compatible con la suma y con el producto por n\'umeros positivos.

\begin{proposicion}[Compatibilidad del orden]
Sean \(x,y,z\in\mathbb{R}\).
\begin{enumerate}
  \item Si \(x\leq y\), entonces \(x+z\leq y+z\).
  \item Si \(x\leq y\) y \(z\geq 0\), entonces \(xz\leq yz\).
\end{enumerate}
\end{proposicion}

La propiedad que distingue de forma esencial a los reales de los racionales es la completitud.

\begin{definicion}[Idea de completitud]
Decimos, de manera intuitiva, que \(\mathbb{R}\) es completo porque no tiene huecos: toda colecci\'on no vac\'ia de n\'umeros reales que est\'a acotada superiormente posee una menor cota superior en \(\mathbb{R}\).
\end{definicion}

La formulaci\'on anterior contiene t\'erminos que se precisar\'an m\'as adelante: conjunto no vac\'io, cota superior y menor cota superior. Esta menor cota superior recibe el nombre de supremo.

\begin{advertencia}
La completitud no significa que entre dos reales haya solamente un n\'umero, sino precisamente lo contrario: entre dos reales distintos hay infinitos n\'umeros reales. La completitud significa que la recta real no presenta huecos.
\end{advertencia}

\section{Intervalos}

Los intervalos son subconjuntos de \(\mathbb{R}\) definidos mediante desigualdades. Antes de usarlos, fijamos la notaci\'on.

\begin{definicion}[Intervalos abiertos y cerrados]
Sean \(a,b\in\mathbb{R}\) con \(a<b\).
\begin{enumerate}
  \item El intervalo abierto entre \(a\) y \(b\) es
  \[
    (a,b)=\{x\in\mathbb{R}: a<x<b\}.
  \]
  \item El intervalo cerrado entre \(a\) y \(b\) es
  \[
    [a,b]=\{x\in\mathbb{R}: a\leq x\leq b\}.
  \]
  \item Los intervalos semiabiertos o semicerrados son
  \[
    (a,b]=\{x\in\mathbb{R}: a<x\leq b\},
  \]
  y
  \[
    [a,b)=\{x\in\mathbb{R}: a\leq x<b\}.
  \]
\end{enumerate}
\end{definicion}

\begin{ejemplo}
El intervalo \((1,4)\) contiene todos los reales mayores que \(1\) y menores que \(4\). No contiene ni al \(1\) ni al \(4\).

El intervalo \([1,4]\) contiene todos los reales mayores o iguales que \(1\) y menores o iguales que \(4\). S\'i contiene a sus extremos.
\end{ejemplo}

Tambi\'en se usan intervalos no acotados.

\begin{definicion}[Intervalos no acotados]
Sea \(a\in\mathbb{R}\). Se definen
\[
  (a,+\infty)=\{x\in\mathbb{R}:x>a\},
\]
\[
  [a,+\infty)=\{x\in\mathbb{R}:x\geq a\},
\]
\[
  (-\infty,a)=\{x\in\mathbb{R}:x<a\},
\]
y
\[
  (-\infty,a]=\{x\in\mathbb{R}:x\leq a\}.
\]
\end{definicion}

\begin{advertencia}
Los s\'imbolos \(+\infty\) y \(-\infty\) no representan n\'umeros reales. Se usan para indicar que un intervalo no est\'a acotado hacia la derecha o hacia la izquierda.
\end{advertencia}

\section{Valor absoluto y distancia}

El valor absoluto mide la distancia de un n\'umero real al cero.

\begin{definicion}[Valor absoluto]
Sea \(x\in\mathbb{R}\). El valor absoluto de \(x\), denotado por \(|x|\), se define por
\[
  |x|=
  \begin{cases}
    x, & \text{si } x\geq 0,\\
    -x, & \text{si } x<0.
  \end{cases}
\]
\end{definicion}

\begin{ejemplo}
Se tiene
\[
  |5|=5,\qquad |-5|=5,\qquad |0|=0.
\]
\end{ejemplo}

El valor absoluto permite definir la distancia entre dos n\'umeros reales.

\begin{definicion}[Distancia en la recta real]
Sean \(x,y\in\mathbb{R}\). La distancia entre \(x\) e \(y\) se define por
\[
  d(x,y)=|x-y|.
\]
\end{definicion}

\begin{ejemplo}
La distancia entre \(2\) y \(7\) es
\[
  d(2,7)=|2-7|=|-5|=5.
\]
La distancia entre \(-3\) y \(4\) es
\[
  d(-3,4)=|-3-4|=|-7|=7.
\]
\end{ejemplo}

\begin{proposicion}[Propiedades b\'asicas del valor absoluto]
Para todo \(x,y\in\mathbb{R}\), se cumple:
\begin{enumerate}
  \item \(|x|\geq 0\).
  \item \(|x|=0\) si y solo si \(x=0\).
  \item \(|xy|=|x|\,|y|\).
  \item \(|x+y|\leq |x|+|y|\).
\end{enumerate}
\end{proposicion}

La cuarta propiedad se llama desigualdad triangular.

\section{Supremo e \'infimo}

Las nociones de supremo e \'infimo formalizan la idea de menor cota superior y mayor cota inferior.

\begin{definicion}[Cota superior]
Sea \(A\subset\mathbb{R}\). Un n\'umero \(M\in\mathbb{R}\) se llama cota superior de \(A\) si
\[
  x\leq M
\]
para todo \(x\in A\).
\end{definicion}

\begin{definicion}[Cota inferior]
Sea \(A\subset\mathbb{R}\). Un n\'umero \(m\in\mathbb{R}\) se llama cota inferior de \(A\) si
\[
  m\leq x
\]
para todo \(x\in A\).
\end{definicion}

\begin{definicion}[Supremo]
Sea \(A\subset\mathbb{R}\) un conjunto no vac\'io y acotado superiormente. Un n\'umero \(s\in\mathbb{R}\) se llama supremo de \(A\), y se denota por
\[
  s=\sup A,
\]
si cumple las dos condiciones siguientes:
\begin{enumerate}
  \item \(s\) es una cota superior de \(A\).
  \item Si \(M\) es cualquier otra cota superior de \(A\), entonces \(s\leq M\).
\end{enumerate}
\end{definicion}

\begin{definicion}[\'Infimo]
Sea \(A\subset\mathbb{R}\) un conjunto no vac\'io y acotado inferiormente. Un n\'umero \(i\in\mathbb{R}\) se llama \'infimo de \(A\), y se denota por
\[
  i=\inf A,
\]
si cumple las dos condiciones siguientes:
\begin{enumerate}
  \item \(i\) es una cota inferior de \(A\).
  \item Si \(m\) es cualquier otra cota inferior de \(A\), entonces \(m\leq i\).
\end{enumerate}
\end{definicion}

\begin{ejemplo}
Sea
\[
  A=(0,1)=\{x\in\mathbb{R}:0<x<1\}.
\]
Entonces
\[
  \sup A=1,
  \qquad
  \inf A=0.
\]
Obs\'ervese que ni \(0\) ni \(1\) pertenecen al conjunto \(A\). Por tanto, el supremo o el \'infimo de un conjunto no tienen por qu\'e pertenecer al conjunto.
\end{ejemplo}

\begin{ejemplo}
Sea
\[
  B=[0,1]=\{x\in\mathbb{R}:0\leq x\leq 1\}.
\]
Entonces
\[
  \sup B=1,
  \qquad
  \inf B=0.
\]
En este caso, tanto el supremo como el \'infimo pertenecen a \(B\).
\end{ejemplo}

\begin{proposicion}[Propiedad de completitud de \(\mathbb{R}\)]
Todo subconjunto no vac\'io de \(\mathbb{R}\) que est\'a acotado superiormente tiene supremo en \(\mathbb{R}\).
\end{proposicion}

Esta propiedad no es cierta en \(\mathbb{Q}\) si se exige que el supremo pertenezca a \(\mathbb{Q}\). Por ejemplo, el conjunto de racionales positivos cuyo cuadrado es menor que \(2\) tiene como menor cota superior a \(\sqrt{2}\), que no es racional.

\begin{erroresfrecuentes}
  \item Confundir densidad con completitud.
  \item Creer que todo n\'umero decimal es racional.
  \item Pensar que \(+\infty\) y \(-\infty\) son n\'umeros reales.
  \item Confundir m\'aximo con supremo, o m\'inimo con \'infimo.
  \item Suponer que el supremo de un conjunto debe pertenecer siempre al conjunto.
\end{erroresfrecuentes}

\begin{seminario}
  \item Trabajo con intervalos y desigualdades.
  \item C\'alculo de supremos e \'infimos en ejemplos sencillos.
\end{seminario}

\begin{ejerciciospropuestos}
  \item Decide si los siguientes intervalos contienen sus extremos:
  \[
    (2,5),\qquad [2,5),\qquad (2,5],\qquad [2,5].
  \]
  \item Calcula \(|x|\) para \(x=-7\), \(x=0\) y \(x=\frac{3}{2}\).
  \item Calcula la distancia entre \(-4\) y \(6\).
  \item Halla el supremo y el \'infimo de los conjuntos
  \[
    (1,3),\qquad [1,3),\qquad \left\{\frac{1}{n}:n\in\mathbb{N},\ n\geq 1\right\}.
  \]
  \item Da un ejemplo de un subconjunto de \(\mathbb{R}\) que tenga supremo pero no tenga m\'aximo.
\end{ejerciciospropuestos}

\resumen{El conjunto de los n\'umeros reales ampl\'ia a los racionales incorporando los n\'umeros irracionales y proporcionando un sistema num\'erico completo. La completitud se expresa mediante la existencia de supremos para subconjuntos no vac\'ios y acotados superiormente. Los intervalos, el valor absoluto y la distancia permiten describir la geometr\'ia elemental de la recta real.}


% ------------------------------------------------------------

\chapter{Los números complejos y las ecuaciones algebraicas}

\begin{objetivos}
  \item Motivar la introducción de los números complejos.
  \item Definir la unidad imaginaria y la forma binómica.
  \item Estudiar conjugado, módulo, inverso y representación geométrica.
  \item Resolver ecuaciones algebraicas sencillas en el conjunto de los números complejos.
\end{objetivos}

\section{Motivación: ecuaciones sin solución real}

En los capítulos anteriores se han introducido sucesivamente los conjuntos numéricos
\(\mathbb{N}\), \(\mathbb{Z}\), \(\mathbb{Q}\) y \(\mathbb{R}\). Cada ampliación ha respondido a una necesidad matemática concreta. Los números enteros permiten resolver ecuaciones de la forma
\[
  a+x=b,
\]
cuando \(a\) y \(b\) son números naturales. Los números racionales permiten resolver ecuaciones de la forma
\[
  ax=b,
\]
cuando \(a\) y \(b\) son enteros y \(a\neq 0\). Los números reales permiten incorporar magnitudes que no pueden representarse mediante fracciones, como ocurre con \(\sqrt{2}\).

Sin embargo, incluso en \(\mathbb{R}\) existen ecuaciones algebraicas sencillas que no tienen solución. Por ejemplo,
\[
  x^2+1=0
\]
no tiene solución real, porque el cuadrado de todo número real es mayor o igual que cero. Es decir, si \(x\in\mathbb{R}\), entonces \(x^2\geq 0\), por lo que \(x^2+1>0\).

La introducción de los números complejos responde a la necesidad de ampliar el sistema numérico para que ecuaciones como la anterior tengan solución.

\begin{advertencia}
La introducción de un nuevo conjunto numérico no debe entenderse como un artificio formal arbitrario. En matemáticas, una ampliación numérica se justifica porque permite resolver problemas que no podían tratarse adecuadamente en el sistema anterior.
\end{advertencia}

\section{La unidad imaginaria}

\begin{definicion}[Unidad imaginaria]
Se llama \emph{unidad imaginaria} a un símbolo, denotado por \(i\), que satisface la relación
\[
  i^2=-1.
\]
\end{definicion}

La igualdad \(i^2=-1\) no afirma que exista un número real cuyo cuadrado sea negativo. Afirma que se introduce un nuevo objeto matemático, llamado \(i\), que permite extender el cálculo algebraico ordinario.

A partir de \(i\), la ecuación
\[
  x^2+1=0
\]
tiene las soluciones
\[
  x=i \qquad \text{y} \qquad x=-i,
\]
porque
\[
  i^2+1=-1+1=0,
\]
y también
\[
  (-i)^2+1=i^2+1=-1+1=0.
\]

\begin{convencion}
A lo largo de este capítulo, el símbolo \(i\) denotará siempre la unidad imaginaria, es decir, el objeto que satisface \(i^2=-1\).
\end{convencion}

\section{Forma binómica}

\begin{definicion}[Número complejo]
Un \emph{número complejo} es una expresión de la forma
\[
  z=a+bi,
\]
donde \(a\) y \(b\) son números reales, e \(i\) es la unidad imaginaria.

El número real \(a\) se llama \emph{parte real} de \(z\), y el número real \(b\) se llama \emph{parte imaginaria} de \(z\). Se escribe
\[
  \operatorname{Re}(z)=a,
  \qquad
  \operatorname{Im}(z)=b.
\]
\end{definicion}

El conjunto de todos los números complejos se denota por
\[
  \mathbb{C}=\{a+bi: a,b\in\mathbb{R}\}.
\]
Formalmente, puede construirse \(\mathbb{C}\) a partir de pares de números reales. En este curso utilizaremos la escritura binómica \(a+bi\), que permite trabajar de manera directa con la unidad imaginaria.

La expresión \(a+bi\) se llama \emph{forma binómica} del número complejo.

\begin{ejemplo}
El número
\[
  z=3-2i
\]
es un número complejo. Su parte real es \(3\) y su parte imaginaria es \(-2\). Por tanto,
\[
  \operatorname{Re}(z)=3,
  \qquad
  \operatorname{Im}(z)=-2.
\]
\end{ejemplo}

Todo número real puede considerarse como un número complejo cuya parte imaginaria es cero. En efecto, si \(a\in\mathbb{R}\), entonces
\[
  a=a+0i.
\]
Por esta razón se identifica \(\mathbb{R}\) con un subconjunto de \(\mathbb{C}\).

\begin{advertencia}
La parte imaginaria de \(a+bi\) es \(b\), no \(bi\). Por ejemplo, si \(z=5+7i\), entonces \(\operatorname{Im}(z)=7\), no \(7i\).
\end{advertencia}

\section{Suma y producto de complejos}

Sean \(z,w\in\mathbb{C}\). Supongamos que
\[
  z=a+bi,
  \qquad
  w=c+di,
\]
donde \(a,b,c,d\in\mathbb{R}\).

\begin{definicion}[Suma de números complejos]
La \emph{suma} de \(z=a+bi\) y \(w=c+di\) se define por
\[
  z+w=(a+c)+(b+d)i.
\]
\end{definicion}

Es decir, se suman por separado las partes reales y las partes imaginarias.

\begin{ejemplo}
Si \(z=2+3i\) y \(w=5-i\), entonces
\[
  z+w=(2+5)+(3-1)i=7+2i.
\]
\end{ejemplo}

\begin{definicion}[Producto de números complejos]
El \emph{producto} de \(z=a+bi\) y \(w=c+di\) se define por
\[
  zw=(ac-bd)+(ad+bc)i.
\]
\end{definicion}

Esta fórmula se obtiene multiplicando formalmente y usando la relación \(i^2=-1\):
\[
(a+bi)(c+di)=ac+adi+bci+bdi^2=(ac-bd)+(ad+bc)i.
\]

\begin{ejemplo}
Si \(z=1+2i\) y \(w=3-i\), entonces
\[
  zw=(1+2i)(3-i)=3-i+6i-2i^2=5+5i.
\]
\end{ejemplo}

\section{Conjugado, módulo e inverso}

\begin{definicion}[Conjugado]
Sea \(z=a+bi\in\mathbb{C}\), con \(a,b\in\mathbb{R}\). El \emph{conjugado} de \(z\), denotado por \(\overline{z}\), se define por
\[
  \overline{z}=a-bi.
\]
\end{definicion}

El conjugado conserva la parte real y cambia el signo de la parte imaginaria.

\begin{ejemplo}
Si \(z=4-3i\), entonces
\[
  \overline{z}=4+3i.
\]
\end{ejemplo}

\begin{definicion}[Módulo]
Sea \(z=a+bi\in\mathbb{C}\), con \(a,b\in\mathbb{R}\). El \emph{módulo} de \(z\), denotado por \(|z|\), se define por
\[
  |z|=\sqrt{a^2+b^2}.
\]
\end{definicion}

El módulo de un número complejo es siempre un número real no negativo.

\begin{proposicion}
Para todo número complejo \(z\in\mathbb{C}\), se cumple
\[
  z\overline{z}=|z|^2.
\]
\end{proposicion}

\begin{proof}
Sea \(z=a+bi\), con \(a,b\in\mathbb{R}\). Entonces \(\overline{z}=a-bi\). Por tanto,
\[
  z\overline{z}=(a+bi)(a-bi)=a^2+b^2.
\]
Por definición de módulo,
\[
  |z|=\sqrt{a^2+b^2},
\]
luego
\[
  |z|^2=a^2+b^2.
\]
Por consiguiente,
\[
  z\overline{z}=|z|^2.
\]
\end{proof}

\begin{definicion}[Inverso]
Sea \(z\in\mathbb{C}\) un número complejo no nulo. El \emph{inverso} de \(z\) es el número complejo \(z^{-1}\) que satisface
\[
  zz^{-1}=1.
\]
Si \(z=a+bi\neq 0\), entonces
\[
  z^{-1}=\frac{\overline{z}}{|z|^2}
  =\frac{a-bi}{a^2+b^2}.
\]
\end{definicion}

\begin{ejemplo}
Sea \(z=1+2i\). Entonces
\[
  \overline{z}=1-2i,
  \qquad
  |z|^2=1^2+2^2=5.
\]
Por tanto,
\[
  z^{-1}=\frac{1-2i}{5}=\frac{1}{5}-\frac{2}{5}i.
\]
\end{ejemplo}

\section{El plano complejo}

Todo número complejo \(z=a+bi\) puede representarse mediante el punto \((a,b)\) del plano real \(\mathbb{R}^2\). En esta representación, el eje horizontal corresponde a la parte real y el eje vertical corresponde a la parte imaginaria.

\begin{definicion}[Plano complejo]
Se llama \emph{plano complejo} a la representación geométrica de \(\mathbb{C}\) en la que a cada número complejo \(z=a+bi\) se le asocia el punto \((a,b)\in\mathbb{R}^2\).
\end{definicion}

En el plano complejo, el módulo \(|z|\) representa la distancia desde el origen hasta el punto asociado a \(z\).

\begin{ejemplo}
El número complejo \(z=3+4i\) se representa mediante el punto \((3,4)\). Su módulo es
\[
  |z|=\sqrt{3^2+4^2}=5.
\]
\end{ejemplo}

\section{Forma polar}

La forma binómica describe un número complejo mediante sus coordenadas cartesianas. La forma polar lo describe mediante su distancia al origen y el ángulo que forma con el eje real positivo.

\begin{definicion}[Argumento]
Sea \(z=a+bi\in\mathbb{C}\) un número complejo no nulo. Un \emph{argumento} de \(z\) es un número real \(\theta\) tal que
\[
  a=|z|\cos\theta,
  \qquad
  b=|z|\sin\theta.
\]
\end{definicion}

Si \(r=|z|\), entonces la forma polar de \(z\) es
\[
  z=r(\cos\theta+i\sin\theta).
\]

\begin{ejemplo}
El número \(z=1+i\) tiene módulo
\[
  |z|=\sqrt{1^2+1^2}=\sqrt{2}.
\]
Un argumento de \(z\) es \(\theta=\pi/4\). Por tanto,
\[
  1+i=\sqrt{2}\left(\cos\frac{\pi}{4}+i\sin\frac{\pi}{4}\right).
\]
\end{ejemplo}

\begin{advertencia}
El argumento de un número complejo no nulo no es único. Si \(\theta\) es un argumento de \(z\), entonces también lo es \(\theta+2k\pi\), para todo \(k\in\mathbb{Z}\).
\end{advertencia}

\section{Fórmula de De Moivre}

La forma polar permite describir de manera sencilla las potencias de un número complejo.

\begin{teorema}[Fórmula de De Moivre]
Sea \(\theta\in\mathbb{R}\) y sea \(n\in\mathbb{N}\). Entonces
\[
  (\cos\theta+i\sin\theta)^n=\cos(n\theta)+i\sin(n\theta).
\]
\end{teorema}

\begin{ejemplo}
Para \(n=2\), la fórmula afirma que
\[
  (\cos\theta+i\sin\theta)^2=\cos(2\theta)+i\sin(2\theta).
\]
Al desarrollar el cuadrado se obtiene
\[
  \cos^2\theta-\sin^2\theta+2i\sin\theta\cos\theta,
\]
lo que coincide con las identidades trigonométricas
\[
  \cos(2\theta)=\cos^2\theta-\sin^2\theta,
  \qquad
  \sin(2\theta)=2\sin\theta\cos\theta.
\]
\end{ejemplo}

\section{Raíces de números complejos}

\begin{definicion}[Raíz compleja]
Sea \(w\in\mathbb{C}\) y sea \(n\in\mathbb{N}\), con \(n\geq 1\). Un número complejo \(z\in\mathbb{C}\) se llama \emph{raíz n-ésima} de \(w\) si
\[
  z^n=w.
\]
\end{definicion}

Si \(w\neq 0\) se escribe en forma polar como
\[
  w=r(\cos\theta+i\sin\theta),
\]
con \(r>0\), entonces sus raíces n-ésimas están dadas por
\[
  z_k=\sqrt[n]{r}\left(\cos\frac{\theta+2k\pi}{n}
  +i\sin\frac{\theta+2k\pi}{n}\right),
  \qquad k=0,1,\dots,n-1.
\]

\begin{ejemplo}[Raíces cuadradas de \(-1\)]
El número \(-1\) puede escribirse como
\[
  -1=\cos\pi+i\sin\pi.
\]
Sus raíces cuadradas son
\[
  z_0=\cos\frac{\pi}{2}+i\sin\frac{\pi}{2}=i,
\]
y
\[
  z_1=\cos\frac{3\pi}{2}+i\sin\frac{3\pi}{2}=-i.
\]
Por tanto, las soluciones de \(z^2=-1\) son \(i\) y \(-i\).
\end{ejemplo}

\section{Ecuaciones algebraicas}

Una ecuación algebraica es una ecuación construida a partir de polinomios. En este capítulo consideraremos principalmente ecuaciones de la forma
\[
  p(z)=0,
\]
donde \(p(z)\) es un polinomio con coeficientes reales o complejos.

\begin{ejemplo}
La ecuación
\[
  z^2+1=0
\]
tiene en \(\mathbb{C}\) las soluciones
\[
  z=i
  \qquad \text{y} \qquad
  z=-i.
\]
\end{ejemplo}

\begin{ejemplo}
Consideremos la ecuación
\[
  z^2-2z+2=0.
\]
Aplicando la fórmula general para ecuaciones cuadráticas se obtiene
\[
  z=\frac{2\pm\sqrt{(-2)^2-4\cdot 1\cdot 2}}{2}
  =\frac{2\pm\sqrt{-4}}{2}.
\]
Como \(\sqrt{-4}=2i\), resulta
\[
  z=\frac{2\pm 2i}{2}=1\pm i.
\]
Por tanto, las soluciones son
\[
  z=1+i
  \qquad \text{y} \qquad
  z=1-i.
\]
\end{ejemplo}

\begin{advertencia}
La fórmula cuadrática se puede utilizar en \(\mathbb{C}\), pero hay que interpretar correctamente las raíces cuadradas de números negativos mediante la unidad imaginaria.
\end{advertencia}

\section{Teorema Fundamental del Álgebra}

El resultado central que justifica la importancia de los números complejos en álgebra es el Teorema Fundamental del Álgebra.

\begin{teorema}[Teorema Fundamental del Álgebra]
Todo polinomio no constante con coeficientes complejos tiene al menos una raíz compleja.
\end{teorema}

Una consecuencia importante es que todo polinomio de grado \(n\geq 1\) con coeficientes complejos puede factorizarse como producto de \(n\) factores lineales, contando las raíces con su multiplicidad.

\begin{ejemplo}
El polinomio
\[
  p(z)=z^2+1
\]
se factoriza en \(\mathbb{C}\) como
\[
  z^2+1=(z-i)(z+i).
\]
\end{ejemplo}

\begin{advertencia}
El Teorema Fundamental del Álgebra no afirma que las raíces puedan calcularse siempre mediante una fórmula sencilla. Afirma la existencia de raíces complejas para todo polinomio no constante con coeficientes complejos.
\end{advertencia}

\begin{erroresfrecuentes}
  \item Pensar que la parte imaginaria de \(a+bi\) es \(bi\), cuando en realidad es \(b\).
  \item Olvidar que \(i^2=-1\) al multiplicar números complejos.
  \item Confundir el conjugado \(\overline{z}\) con el opuesto \(-z\).
  \item Pensar que el argumento de un número complejo no nulo es único.
  \item Calcular raíces complejas usando solo una rama del argumento y perder soluciones.
\end{erroresfrecuentes}

\begin{seminario}
  \item Operaciones en forma binómica.
  \item Representación de complejos en el plano.
  \item Cálculo de raíces complejas.
\end{seminario}

\begin{ejerciciospropuestos}
  \item Calcula la suma, el producto, el conjugado y el módulo de los números complejos \(z=2-3i\) y \(w=-1+4i\).
  \item Determina el inverso de \(z=3+i\).
  \item Representa en el plano complejo los números \(1+i\), \(-2+3i\), \(-1-i\) y \(4\).
  \item Escribe en forma polar el número complejo \(z=1-i\).
  \item Calcula todas las raíces cúbicas de \(1\).
  \item Resuelve en \(\mathbb{C}\) la ecuación \(z^2+4=0\).
  \item Resuelve en \(\mathbb{C}\) la ecuación \(z^2-4z+5=0\).
\end{ejerciciospropuestos}

\resumen{Los números complejos amplían los números reales introduciendo la unidad imaginaria \(i\), definida por \(i^2=-1\). Todo complejo se escribe en forma binómica \(a+bi\), puede representarse en el plano y admite operaciones algebraicas compatibles con la relación fundamental de la unidad imaginaria. El conjugado, el módulo y la forma polar permiten estudiar productos, potencias y raíces. El Teorema Fundamental del Álgebra muestra que \(\mathbb{C}\) es el marco natural para resolver ecuaciones polinómicas.}


% ============================================================
% Apéndices
% ============================================================
\appendix
\part{Apéndices}

\chapter{Guía de notación}

Este apéndice recoge la notación principal utilizada en el manual. La regla de uso será siempre la misma: un símbolo solo se incorpora a esta guía después de haber sido definido en el texto principal.

\begin{longtable}{>{\bfseries}p{0.18\textwidth}p{0.72\textwidth}}
\toprule
Símbolo & Significado \\
\midrule
\endhead
\(\in\) & Relación de pertenencia entre un elemento y un conjunto. \\
\(\subseteq\) & Relación de inclusión entre conjuntos. \\
\(\varnothing\) & Conjunto vacío. \\
\(\Pow(A)\) & Conjunto de las partes de un conjunto \(A\). \\
\(A\times B\) & Producto cartesiano de los conjuntos \(A\) y \(B\). \\
\(f:A\to B\) & Aplicación con dominio \(A\) y codominio \(B\). \\
\(\id_A\) & Aplicación identidad del conjunto \(A\). \\
\(f^{-1}(C)\) & Imagen recíproca de \(C\) mediante \(f\). \\
\(\sim\) & Relación de equivalencia. \\
\([a]\) & Clase de equivalencia del elemento \(a\). \\
\(\N\) & Conjunto de los números naturales, incluyendo el cero. \\
\(\N^{\ast}\) & Conjunto de los números naturales positivos. \\
\(\Z\) & Conjunto de los números enteros. \\
\(a\mid b\) & Relación de divisibilidad: \(a\) divide a \(b\). \\
\(\mcd(a,b)\) & Máximo común divisor de \(a\) y \(b\). \\
\(\mcm(a,b)\) & Mínimo común múltiplo de \(a\) y \(b\). \\
\(a\equiv b\pmod n\) & Congruencia de \(a\) y \(b\) módulo \(n\). \\
\(\Z/n\Z\) & Conjunto de clases de congruencia módulo \(n\). \\
\(\Q\) & Conjunto de los números racionales. \\
\(\R\) & Conjunto de los números reales. \\
\(|x|\) & Valor absoluto de un número real o módulo de un número complejo. \\
\(\sup A\) & Supremo del conjunto \(A\), cuando existe. \\
\(\inf A\) & Ínfimo del conjunto \(A\), cuando existe. \\
\(\C\) & Conjunto de los números complejos. \\
\(i\) & Unidad imaginaria, definida por \(i^2=-1\). \\
\(\overline{z}\) & Conjugado del número complejo \(z\). \\
\(\operatorname{Re}(z)\) & Parte real del número complejo \(z\). \\
\(\operatorname{Im}(z)\) & Parte imaginaria del número complejo \(z\). \\
\bottomrule
\end{longtable}

% !TEX root = manual_fundamentos_matematicas.tex
% ------------------------------------------------------------
% Apéndice: Fronteras con otras asignaturas
% ------------------------------------------------------------
% Este fichero está pensado para ser incluido mediante:
% % !TEX root = manual_fundamentos_matematicas.tex
% ------------------------------------------------------------
% Apéndice: Fronteras con otras asignaturas
% ------------------------------------------------------------
% Este fichero está pensado para ser incluido mediante:
% % !TEX root = manual_fundamentos_matematicas.tex
% ------------------------------------------------------------
% Apéndice: Fronteras con otras asignaturas
% ------------------------------------------------------------
% Este fichero está pensado para ser incluido mediante:
% \input{apendice_fronteras_asignaturas.tex}
%
% Requiere que el documento principal tenga definidos los entornos
% convencion, advertencia, seminario y ejerciciospropuestos.

\chapter{Fronteras con otras asignaturas}
\label{ap:fronteras-asignaturas}

\section{Finalidad del apéndice}

Este apéndice tiene como finalidad delimitar el papel de \emph{Fundamentos de Matemáticas I} y \emph{Fundamentos de Matemáticas II} dentro del primer curso del Grado en Ingeniería Matemática. La asignatura no debe entenderse como una repetición simplificada de Análisis Matemático, Álgebra Lineal, Matemática Discreta, Algoritmia o Métodos Numéricos, sino como una asignatura de soporte transversal.

El objetivo principal de \emph{Fundamentos de Matemáticas} es proporcionar al estudiante el lenguaje, la notación, los métodos elementales de demostración y las estructuras numéricas básicas que necesitará en el resto del grado.

\begin{convencion}[Principio de delimitación]
En \emph{Fundamentos de Matemáticas} se introducen el lenguaje, la notación, las definiciones básicas y los métodos elementales de razonamiento. Las asignaturas específicas desarrollan después las teorías particulares, sus técnicas propias y sus aplicaciones.
\end{convencion}

Esta delimitación permite gestionar los solapamientos de forma positiva: un mismo concepto puede aparecer en varias asignaturas, pero con funciones distintas. En \emph{Fundamentos de Matemáticas} aparece como lenguaje o herramienta básica; en la asignatura específica aparece como objeto de estudio técnico.

\section{Frontera con Análisis Matemático}

\begin{convencion}[Frontera con Análisis Matemático]
En \emph{Fundamentos de Matemáticas} se introducen los números reales, los intervalos, el valor absoluto, la distancia, el supremo y el ínfimo como lenguaje básico. El estudio sistemático de sucesiones, límites, continuidad, derivación, integración y técnicas de demostración propias del análisis pertenece a las asignaturas de Análisis Matemático.
\end{convencion}

En particular, en \emph{Fundamentos de Matemáticas} se debe trabajar la comprensión de expresiones como
\[
  |x-a|<\varepsilon,
\]
donde previamente se han definido los símbolos que aparecen: \(x\) y \(a\) son números reales, \(\varepsilon\) es un número real positivo, \(|x-a|\) denota el valor absoluto de la diferencia \(x-a\), y el símbolo \(<\) representa la relación de orden usual en \(\mathbb{R}\).

Sin embargo, no corresponde a \emph{Fundamentos de Matemáticas} desarrollar con profundidad las definiciones \(\varepsilon\)-\(\delta\) de límite o continuidad. Esas definiciones pertenecen a Análisis Matemático, aunque requieran el lenguaje preparado en este curso.

\begin{advertencia}
El hecho de que un concepto aparezca en \emph{Fundamentos de Matemáticas} no significa que se agote en esta asignatura. Por ejemplo, los números reales se introducen aquí como sistema numérico y como lenguaje de orden, distancia y aproximación; su estudio analítico se desarrolla posteriormente en Análisis Matemático.
\end{advertencia}

\section{Frontera con Álgebra Lineal}

\begin{convencion}[Frontera con Álgebra Lineal]
En \emph{Fundamentos de Matemáticas} se introducen conjuntos, aplicaciones, operaciones internas y las primeras estructuras algebraicas. La teoría de espacios vectoriales, matrices, sistemas lineales, bases, dimensión, transformaciones lineales, determinantes, autovalores y diagonalización corresponde a Álgebra Lineal.
\end{convencion}

La asignatura \emph{Fundamentos de Matemáticas} debe preparar al estudiante para entender expresiones del tipo
\[
  f:A\longrightarrow B,
\]
donde \(A\) y \(B\) son conjuntos y \(f\) es una aplicación de \(A\) en \(B\). También debe introducir conceptos como composición, aplicación identidad, aplicación inversa, inyectividad, sobreyectividad y biyectividad.

Álgebra Lineal utilizará después ese lenguaje para estudiar aplicaciones lineales entre espacios vectoriales. Por tanto, \emph{Fundamentos de Matemáticas} debe evitar desarrollar la teoría específica de espacios vectoriales más allá de mencionarlos como ejemplo de estructura algebraica.

\begin{advertencia}
En \emph{Fundamentos de Matemáticas} puede citarse que un espacio vectorial es una estructura algebraica, pero no deben estudiarse en detalle bases, dimensión, rango, determinantes o diagonalización, porque esos contenidos pertenecen a Álgebra Lineal.
\end{advertencia}

\section{Frontera con Matemática Discreta}

\begin{convencion}[Frontera con Matemática Discreta]
En \emph{Fundamentos de Matemáticas} se estudian los números naturales, la inducción, las definiciones recursivas, conjuntos y relaciones como lenguaje general. Las aplicaciones a combinatoria, probabilidad discreta y grafos pertenecen a Matemática Discreta.
\end{convencion}

La inducción matemática debe aparecer en \emph{Fundamentos de Matemáticas} como método de demostración. Por ejemplo, se puede demostrar que, para todo número natural \(n\geq 1\),
\[
  1+2+\cdots+n=\frac{n(n+1)}{2}.
\]

En cambio, el uso de la inducción para justificar fórmulas combinatorias, propiedades de grafos o algoritmos discretos puede desarrollarse en Matemática Discreta.

\begin{advertencia}
La asignatura \emph{Fundamentos de Matemáticas} no debe absorber contenidos propios de Matemática Discreta. En particular, no debe convertirse en un curso de combinatoria, grafos o probabilidad discreta.
\end{advertencia}

\section{Frontera con Algoritmia}

\begin{convencion}[Frontera con Algoritmia]
En \emph{Fundamentos de Matemáticas II} se estudian procedimientos aritméticos como el algoritmo de Euclides desde el punto de vista matemático: definición, justificación y demostración de corrección. La implementación mediante variables, funciones, estructuras de control, iteración, recursividad y programación corresponde a Algoritmia.
\end{convencion}

Por ejemplo, en \emph{Fundamentos de Matemáticas II} se puede estudiar que, dados dos enteros \(a\) y \(b\), no ambos nulos, el algoritmo de Euclides permite calcular su máximo común divisor. También se puede demostrar que existen enteros \(u\) y \(v\) tales que
\[
  \operatorname{mcd}(a,b)=ua+vb.
\]

Algoritmia podrá después traducir este procedimiento a pseudocódigo o a un lenguaje de programación concreto.

\begin{advertencia}
En \emph{Fundamentos de Matemáticas} se demuestra por qué un procedimiento funciona. En Algoritmia se estudia cómo implementarlo de forma correcta y eficiente.
\end{advertencia}

\section{Frontera con Métodos Numéricos}

\begin{convencion}[Frontera con Métodos Numéricos]
En \emph{Fundamentos de Matemáticas II} se introducen el valor absoluto, la distancia, los intervalos, las desigualdades y la noción elemental de aproximación. El estudio de errores, ceros de funciones, interpolación, diferenciación e integración numérica, así como la resolución numérica de sistemas, pertenece a Métodos Numéricos.
\end{convencion}

El estudiante debe llegar a Métodos Numéricos con soltura en expresiones como
\[
  |x-y|, \qquad [a,b], \qquad x\in (a,b), \qquad |x-a|<\varepsilon.
\]

En \emph{Fundamentos de Matemáticas} estas expresiones se interpretan desde el punto de vista del lenguaje matemático. En Métodos Numéricos se convierten en herramientas para estimar errores, controlar aproximaciones y analizar algoritmos.

\begin{advertencia}
No corresponde a \emph{Fundamentos de Matemáticas} desarrollar métodos de bisección, Newton, interpolación, cuadratura o análisis de error. Sí corresponde preparar el lenguaje de desigualdades, distancia y aproximación que estos métodos necesitan.
\end{advertencia}

\section{Frontera con Estructuras Algebraicas}

\begin{convencion}[Frontera con Estructuras Algebraicas]
En \emph{Fundamentos de Matemáticas} los conceptos de semigrupo, monoide, grupo, anillo y cuerpo se presentan de forma introductoria, como ejemplos de estructuras algebraicas definidas mediante operaciones y axiomas. La teoría sistemática de grupos, subgrupos, homomorfismos, cocientes, acciones, anillos y cuerpos pertenece a la asignatura Estructuras Algebraicas.
\end{convencion}

En \emph{Fundamentos de Matemáticas} basta con que el estudiante comprenda que una estructura algebraica consiste en un conjunto dotado de una o varias operaciones que satisfacen ciertas propiedades. Por ejemplo, puede estudiarse que \((\mathbb{Z},+)\) es un grupo abeliano, donde \(\mathbb{Z}\) es el conjunto de los números enteros y \(+\) es la suma usual.

No es necesario desarrollar de forma extensa subgrupos, morfismos, grupos cociente, acciones de grupo o clasificación de estructuras. Estos temas requieren un tratamiento propio y corresponden a una asignatura posterior.

\begin{advertencia}
La función del capítulo de estructuras algebraicas en \emph{Fundamentos de Matemáticas} es preparar el lenguaje axiomático, no adelantar la asignatura completa de Estructuras Algebraicas.
\end{advertencia}

\section{Frontera con Criptografía}

\begin{convencion}[Frontera con Criptografía]
En \emph{Fundamentos de Matemáticas II} se introducen divisibilidad, números primos, máximo común divisor, identidad de Bézout, congruencias, inversos modulares y ecuaciones lineales módulo un entero. El uso de estas herramientas en sistemas criptográficos, protocolos de clave pública, factorización computacional, logaritmo discreto o funciones hash pertenece a Criptografía.
\end{convencion}

La aritmética modular debe estudiarse en \emph{Fundamentos de Matemáticas II} como parte del lenguaje aritmético básico. Por ejemplo, si \(n\) es un entero positivo y \(a,b\in\mathbb{Z}\), se define
\[
  a\equiv b \pmod n
\]
cuando \(n\) divide a \(a-b\).

Criptografía utilizará después esta base para estudiar algoritmos y protocolos concretos. Por tanto, en \emph{Fundamentos de Matemáticas} no es necesario presentar RSA, Diffie--Hellman, ElGamal o criptografía postcuántica, salvo como motivación cultural breve.

\section{Frontera con Análisis Complejo y de Fourier}

\begin{convencion}[Frontera con Análisis Complejo y de Fourier]
En \emph{Fundamentos de Matemáticas II} se estudian los números complejos como sistema numérico: forma binómica, operaciones, conjugado, módulo, inverso, plano complejo, forma polar y raíces. El estudio de funciones complejas, series de potencias, teoría de Cauchy, residuos y series de Fourier pertenece a Análisis Complejo y de Fourier.
\end{convencion}

En \emph{Fundamentos de Matemáticas II} el objetivo es que el estudiante comprenda expresiones como
\[
  z=a+bi,
\]
donde \(a,b\in\mathbb{R}\), \(i\) es la unidad imaginaria y satisface \(i^2=-1\). También debe saber calcular conjugados, módulos, inversos y raíces sencillas.

El estudio de funciones \(f:\mathbb{C}\to\mathbb{C}\), derivabilidad compleja, integrales complejas o residuos queda fuera del alcance de esta asignatura.

\begin{advertencia}
\emph{Fundamentos de Matemáticas II} estudia números complejos. \emph{Análisis Complejo y de Fourier} estudia funciones de variable compleja y herramientas analíticas avanzadas.
\end{advertencia}

\section{Criterios prácticos de coordinación}

Para que estas fronteras sean efectivas, se recomienda aplicar los siguientes criterios de coordinación docente.

\subsection*{Criterio 1: prioridad del lenguaje}

Si un contenido aparece en \emph{Fundamentos de Matemáticas}, debe tratarse en primer lugar desde el punto de vista del lenguaje: símbolos, definiciones, ejemplos, contraejemplos y demostraciones elementales.

\subsection*{Criterio 2: profundidad controlada}

Un concepto que pertenezca de manera natural a una asignatura posterior puede introducirse en \emph{Fundamentos de Matemáticas}, pero solo hasta el nivel necesario para que el estudiante reconozca su significado y pueda utilizarlo en ejemplos básicos.

\subsection*{Criterio 3: evaluación transversal}

La evaluación de \emph{Fundamentos de Matemáticas} debe centrarse en la lectura y escritura de textos matemáticos, la precisión simbólica, la construcción de ejemplos y contraejemplos, y las demostraciones elementales. No debe evaluar contenidos técnicos propios de otras asignaturas.

\subsection*{Criterio 4: coordinación horizontal}

En primer curso, los contenidos de \emph{Fundamentos de Matemáticas} deben ordenarse de manera que sirvan de apoyo inmediato a Análisis Matemático, Álgebra Lineal, Matemática Discreta, Algoritmia y Métodos Numéricos.

\subsection*{Criterio 5: coordinación vertical}

Los capítulos de estructuras algebraicas, aritmética modular, números reales y números complejos deben presentarse como preparación para asignaturas posteriores: Estructuras Algebraicas, Criptografía, Análisis Matemático, Métodos Numéricos, Topología, Análisis Funcional y Análisis Complejo.

\section{Resumen operativo}

\begin{center}
\begin{tabular}{p{0.28\textwidth}p{0.32\textwidth}p{0.32\textwidth}}
\toprule
\textbf{Tema} & \textbf{En Fundamentos de Matemáticas} & \textbf{En la asignatura específica} \\
\midrule
Inducción & Método de demostración & Aplicaciones en combinatoria, grafos y algoritmos \\
Conjuntos y relaciones & Lenguaje básico y ejemplos finitos & Uso sistemático en estructuras discretas y topológicas \\
Aplicaciones & Dominio, codominio, imagen, composición & Aplicaciones lineales, funciones continuas, funciones computables \\
Estructuras algebraicas & Axiomas básicos y ejemplos & Grupos, anillos, cuerpos, cocientes y homomorfismos \\
Números reales & Orden, intervalos, valor absoluto, distancia & Límites, continuidad, derivación e integración \\
Aritmética modular & Congruencias e inversos modulares & Criptografía y algoritmos aritméticos \\
Complejos & Sistema numérico y raíces & Funciones complejas, Cauchy, residuos y Fourier \\
\bottomrule
\end{tabular}
\end{center}

\resumen{Este apéndice establece las fronteras entre \emph{Fundamentos de Matemáticas} y otras asignaturas del grado. La regla general es que Fundamentos introduce lenguaje, notación, razonamiento y estructuras básicas, mientras que las asignaturas específicas desarrollan las teorías técnicas y sus aplicaciones.}

%
% Requiere que el documento principal tenga definidos los entornos
% convencion, advertencia, seminario y ejerciciospropuestos.

\chapter{Fronteras con otras asignaturas}
\label{ap:fronteras-asignaturas}

\section{Finalidad del apéndice}

Este apéndice tiene como finalidad delimitar el papel de \emph{Fundamentos de Matemáticas I} y \emph{Fundamentos de Matemáticas II} dentro del primer curso del Grado en Ingeniería Matemática. La asignatura no debe entenderse como una repetición simplificada de Análisis Matemático, Álgebra Lineal, Matemática Discreta, Algoritmia o Métodos Numéricos, sino como una asignatura de soporte transversal.

El objetivo principal de \emph{Fundamentos de Matemáticas} es proporcionar al estudiante el lenguaje, la notación, los métodos elementales de demostración y las estructuras numéricas básicas que necesitará en el resto del grado.

\begin{convencion}[Principio de delimitación]
En \emph{Fundamentos de Matemáticas} se introducen el lenguaje, la notación, las definiciones básicas y los métodos elementales de razonamiento. Las asignaturas específicas desarrollan después las teorías particulares, sus técnicas propias y sus aplicaciones.
\end{convencion}

Esta delimitación permite gestionar los solapamientos de forma positiva: un mismo concepto puede aparecer en varias asignaturas, pero con funciones distintas. En \emph{Fundamentos de Matemáticas} aparece como lenguaje o herramienta básica; en la asignatura específica aparece como objeto de estudio técnico.

\section{Frontera con Análisis Matemático}

\begin{convencion}[Frontera con Análisis Matemático]
En \emph{Fundamentos de Matemáticas} se introducen los números reales, los intervalos, el valor absoluto, la distancia, el supremo y el ínfimo como lenguaje básico. El estudio sistemático de sucesiones, límites, continuidad, derivación, integración y técnicas de demostración propias del análisis pertenece a las asignaturas de Análisis Matemático.
\end{convencion}

En particular, en \emph{Fundamentos de Matemáticas} se debe trabajar la comprensión de expresiones como
\[
  |x-a|<\varepsilon,
\]
donde previamente se han definido los símbolos que aparecen: \(x\) y \(a\) son números reales, \(\varepsilon\) es un número real positivo, \(|x-a|\) denota el valor absoluto de la diferencia \(x-a\), y el símbolo \(<\) representa la relación de orden usual en \(\mathbb{R}\).

Sin embargo, no corresponde a \emph{Fundamentos de Matemáticas} desarrollar con profundidad las definiciones \(\varepsilon\)-\(\delta\) de límite o continuidad. Esas definiciones pertenecen a Análisis Matemático, aunque requieran el lenguaje preparado en este curso.

\begin{advertencia}
El hecho de que un concepto aparezca en \emph{Fundamentos de Matemáticas} no significa que se agote en esta asignatura. Por ejemplo, los números reales se introducen aquí como sistema numérico y como lenguaje de orden, distancia y aproximación; su estudio analítico se desarrolla posteriormente en Análisis Matemático.
\end{advertencia}

\section{Frontera con Álgebra Lineal}

\begin{convencion}[Frontera con Álgebra Lineal]
En \emph{Fundamentos de Matemáticas} se introducen conjuntos, aplicaciones, operaciones internas y las primeras estructuras algebraicas. La teoría de espacios vectoriales, matrices, sistemas lineales, bases, dimensión, transformaciones lineales, determinantes, autovalores y diagonalización corresponde a Álgebra Lineal.
\end{convencion}

La asignatura \emph{Fundamentos de Matemáticas} debe preparar al estudiante para entender expresiones del tipo
\[
  f:A\longrightarrow B,
\]
donde \(A\) y \(B\) son conjuntos y \(f\) es una aplicación de \(A\) en \(B\). También debe introducir conceptos como composición, aplicación identidad, aplicación inversa, inyectividad, sobreyectividad y biyectividad.

Álgebra Lineal utilizará después ese lenguaje para estudiar aplicaciones lineales entre espacios vectoriales. Por tanto, \emph{Fundamentos de Matemáticas} debe evitar desarrollar la teoría específica de espacios vectoriales más allá de mencionarlos como ejemplo de estructura algebraica.

\begin{advertencia}
En \emph{Fundamentos de Matemáticas} puede citarse que un espacio vectorial es una estructura algebraica, pero no deben estudiarse en detalle bases, dimensión, rango, determinantes o diagonalización, porque esos contenidos pertenecen a Álgebra Lineal.
\end{advertencia}

\section{Frontera con Matemática Discreta}

\begin{convencion}[Frontera con Matemática Discreta]
En \emph{Fundamentos de Matemáticas} se estudian los números naturales, la inducción, las definiciones recursivas, conjuntos y relaciones como lenguaje general. Las aplicaciones a combinatoria, probabilidad discreta y grafos pertenecen a Matemática Discreta.
\end{convencion}

La inducción matemática debe aparecer en \emph{Fundamentos de Matemáticas} como método de demostración. Por ejemplo, se puede demostrar que, para todo número natural \(n\geq 1\),
\[
  1+2+\cdots+n=\frac{n(n+1)}{2}.
\]

En cambio, el uso de la inducción para justificar fórmulas combinatorias, propiedades de grafos o algoritmos discretos puede desarrollarse en Matemática Discreta.

\begin{advertencia}
La asignatura \emph{Fundamentos de Matemáticas} no debe absorber contenidos propios de Matemática Discreta. En particular, no debe convertirse en un curso de combinatoria, grafos o probabilidad discreta.
\end{advertencia}

\section{Frontera con Algoritmia}

\begin{convencion}[Frontera con Algoritmia]
En \emph{Fundamentos de Matemáticas II} se estudian procedimientos aritméticos como el algoritmo de Euclides desde el punto de vista matemático: definición, justificación y demostración de corrección. La implementación mediante variables, funciones, estructuras de control, iteración, recursividad y programación corresponde a Algoritmia.
\end{convencion}

Por ejemplo, en \emph{Fundamentos de Matemáticas II} se puede estudiar que, dados dos enteros \(a\) y \(b\), no ambos nulos, el algoritmo de Euclides permite calcular su máximo común divisor. También se puede demostrar que existen enteros \(u\) y \(v\) tales que
\[
  \operatorname{mcd}(a,b)=ua+vb.
\]

Algoritmia podrá después traducir este procedimiento a pseudocódigo o a un lenguaje de programación concreto.

\begin{advertencia}
En \emph{Fundamentos de Matemáticas} se demuestra por qué un procedimiento funciona. En Algoritmia se estudia cómo implementarlo de forma correcta y eficiente.
\end{advertencia}

\section{Frontera con Métodos Numéricos}

\begin{convencion}[Frontera con Métodos Numéricos]
En \emph{Fundamentos de Matemáticas II} se introducen el valor absoluto, la distancia, los intervalos, las desigualdades y la noción elemental de aproximación. El estudio de errores, ceros de funciones, interpolación, diferenciación e integración numérica, así como la resolución numérica de sistemas, pertenece a Métodos Numéricos.
\end{convencion}

El estudiante debe llegar a Métodos Numéricos con soltura en expresiones como
\[
  |x-y|, \qquad [a,b], \qquad x\in (a,b), \qquad |x-a|<\varepsilon.
\]

En \emph{Fundamentos de Matemáticas} estas expresiones se interpretan desde el punto de vista del lenguaje matemático. En Métodos Numéricos se convierten en herramientas para estimar errores, controlar aproximaciones y analizar algoritmos.

\begin{advertencia}
No corresponde a \emph{Fundamentos de Matemáticas} desarrollar métodos de bisección, Newton, interpolación, cuadratura o análisis de error. Sí corresponde preparar el lenguaje de desigualdades, distancia y aproximación que estos métodos necesitan.
\end{advertencia}

\section{Frontera con Estructuras Algebraicas}

\begin{convencion}[Frontera con Estructuras Algebraicas]
En \emph{Fundamentos de Matemáticas} los conceptos de semigrupo, monoide, grupo, anillo y cuerpo se presentan de forma introductoria, como ejemplos de estructuras algebraicas definidas mediante operaciones y axiomas. La teoría sistemática de grupos, subgrupos, homomorfismos, cocientes, acciones, anillos y cuerpos pertenece a la asignatura Estructuras Algebraicas.
\end{convencion}

En \emph{Fundamentos de Matemáticas} basta con que el estudiante comprenda que una estructura algebraica consiste en un conjunto dotado de una o varias operaciones que satisfacen ciertas propiedades. Por ejemplo, puede estudiarse que \((\mathbb{Z},+)\) es un grupo abeliano, donde \(\mathbb{Z}\) es el conjunto de los números enteros y \(+\) es la suma usual.

No es necesario desarrollar de forma extensa subgrupos, morfismos, grupos cociente, acciones de grupo o clasificación de estructuras. Estos temas requieren un tratamiento propio y corresponden a una asignatura posterior.

\begin{advertencia}
La función del capítulo de estructuras algebraicas en \emph{Fundamentos de Matemáticas} es preparar el lenguaje axiomático, no adelantar la asignatura completa de Estructuras Algebraicas.
\end{advertencia}

\section{Frontera con Criptografía}

\begin{convencion}[Frontera con Criptografía]
En \emph{Fundamentos de Matemáticas II} se introducen divisibilidad, números primos, máximo común divisor, identidad de Bézout, congruencias, inversos modulares y ecuaciones lineales módulo un entero. El uso de estas herramientas en sistemas criptográficos, protocolos de clave pública, factorización computacional, logaritmo discreto o funciones hash pertenece a Criptografía.
\end{convencion}

La aritmética modular debe estudiarse en \emph{Fundamentos de Matemáticas II} como parte del lenguaje aritmético básico. Por ejemplo, si \(n\) es un entero positivo y \(a,b\in\mathbb{Z}\), se define
\[
  a\equiv b \pmod n
\]
cuando \(n\) divide a \(a-b\).

Criptografía utilizará después esta base para estudiar algoritmos y protocolos concretos. Por tanto, en \emph{Fundamentos de Matemáticas} no es necesario presentar RSA, Diffie--Hellman, ElGamal o criptografía postcuántica, salvo como motivación cultural breve.

\section{Frontera con Análisis Complejo y de Fourier}

\begin{convencion}[Frontera con Análisis Complejo y de Fourier]
En \emph{Fundamentos de Matemáticas II} se estudian los números complejos como sistema numérico: forma binómica, operaciones, conjugado, módulo, inverso, plano complejo, forma polar y raíces. El estudio de funciones complejas, series de potencias, teoría de Cauchy, residuos y series de Fourier pertenece a Análisis Complejo y de Fourier.
\end{convencion}

En \emph{Fundamentos de Matemáticas II} el objetivo es que el estudiante comprenda expresiones como
\[
  z=a+bi,
\]
donde \(a,b\in\mathbb{R}\), \(i\) es la unidad imaginaria y satisface \(i^2=-1\). También debe saber calcular conjugados, módulos, inversos y raíces sencillas.

El estudio de funciones \(f:\mathbb{C}\to\mathbb{C}\), derivabilidad compleja, integrales complejas o residuos queda fuera del alcance de esta asignatura.

\begin{advertencia}
\emph{Fundamentos de Matemáticas II} estudia números complejos. \emph{Análisis Complejo y de Fourier} estudia funciones de variable compleja y herramientas analíticas avanzadas.
\end{advertencia}

\section{Criterios prácticos de coordinación}

Para que estas fronteras sean efectivas, se recomienda aplicar los siguientes criterios de coordinación docente.

\subsection*{Criterio 1: prioridad del lenguaje}

Si un contenido aparece en \emph{Fundamentos de Matemáticas}, debe tratarse en primer lugar desde el punto de vista del lenguaje: símbolos, definiciones, ejemplos, contraejemplos y demostraciones elementales.

\subsection*{Criterio 2: profundidad controlada}

Un concepto que pertenezca de manera natural a una asignatura posterior puede introducirse en \emph{Fundamentos de Matemáticas}, pero solo hasta el nivel necesario para que el estudiante reconozca su significado y pueda utilizarlo en ejemplos básicos.

\subsection*{Criterio 3: evaluación transversal}

La evaluación de \emph{Fundamentos de Matemáticas} debe centrarse en la lectura y escritura de textos matemáticos, la precisión simbólica, la construcción de ejemplos y contraejemplos, y las demostraciones elementales. No debe evaluar contenidos técnicos propios de otras asignaturas.

\subsection*{Criterio 4: coordinación horizontal}

En primer curso, los contenidos de \emph{Fundamentos de Matemáticas} deben ordenarse de manera que sirvan de apoyo inmediato a Análisis Matemático, Álgebra Lineal, Matemática Discreta, Algoritmia y Métodos Numéricos.

\subsection*{Criterio 5: coordinación vertical}

Los capítulos de estructuras algebraicas, aritmética modular, números reales y números complejos deben presentarse como preparación para asignaturas posteriores: Estructuras Algebraicas, Criptografía, Análisis Matemático, Métodos Numéricos, Topología, Análisis Funcional y Análisis Complejo.

\section{Resumen operativo}

\begin{center}
\begin{tabular}{p{0.28\textwidth}p{0.32\textwidth}p{0.32\textwidth}}
\toprule
\textbf{Tema} & \textbf{En Fundamentos de Matemáticas} & \textbf{En la asignatura específica} \\
\midrule
Inducción & Método de demostración & Aplicaciones en combinatoria, grafos y algoritmos \\
Conjuntos y relaciones & Lenguaje básico y ejemplos finitos & Uso sistemático en estructuras discretas y topológicas \\
Aplicaciones & Dominio, codominio, imagen, composición & Aplicaciones lineales, funciones continuas, funciones computables \\
Estructuras algebraicas & Axiomas básicos y ejemplos & Grupos, anillos, cuerpos, cocientes y homomorfismos \\
Números reales & Orden, intervalos, valor absoluto, distancia & Límites, continuidad, derivación e integración \\
Aritmética modular & Congruencias e inversos modulares & Criptografía y algoritmos aritméticos \\
Complejos & Sistema numérico y raíces & Funciones complejas, Cauchy, residuos y Fourier \\
\bottomrule
\end{tabular}
\end{center}

\resumen{Este apéndice establece las fronteras entre \emph{Fundamentos de Matemáticas} y otras asignaturas del grado. La regla general es que Fundamentos introduce lenguaje, notación, razonamiento y estructuras básicas, mientras que las asignaturas específicas desarrollan las teorías técnicas y sus aplicaciones.}

%
% Requiere que el documento principal tenga definidos los entornos
% convencion, advertencia, seminario y ejerciciospropuestos.

\chapter{Fronteras con otras asignaturas}
\label{ap:fronteras-asignaturas}

\section{Finalidad del apéndice}

Este apéndice tiene como finalidad delimitar el papel de \emph{Fundamentos de Matemáticas I} y \emph{Fundamentos de Matemáticas II} dentro del primer curso del Grado en Ingeniería Matemática. La asignatura no debe entenderse como una repetición simplificada de Análisis Matemático, Álgebra Lineal, Matemática Discreta, Algoritmia o Métodos Numéricos, sino como una asignatura de soporte transversal.

El objetivo principal de \emph{Fundamentos de Matemáticas} es proporcionar al estudiante el lenguaje, la notación, los métodos elementales de demostración y las estructuras numéricas básicas que necesitará en el resto del grado.

\begin{convencion}[Principio de delimitación]
En \emph{Fundamentos de Matemáticas} se introducen el lenguaje, la notación, las definiciones básicas y los métodos elementales de razonamiento. Las asignaturas específicas desarrollan después las teorías particulares, sus técnicas propias y sus aplicaciones.
\end{convencion}

Esta delimitación permite gestionar los solapamientos de forma positiva: un mismo concepto puede aparecer en varias asignaturas, pero con funciones distintas. En \emph{Fundamentos de Matemáticas} aparece como lenguaje o herramienta básica; en la asignatura específica aparece como objeto de estudio técnico.

\section{Frontera con Análisis Matemático}

\begin{convencion}[Frontera con Análisis Matemático]
En \emph{Fundamentos de Matemáticas} se introducen los números reales, los intervalos, el valor absoluto, la distancia, el supremo y el ínfimo como lenguaje básico. El estudio sistemático de sucesiones, límites, continuidad, derivación, integración y técnicas de demostración propias del análisis pertenece a las asignaturas de Análisis Matemático.
\end{convencion}

En particular, en \emph{Fundamentos de Matemáticas} se debe trabajar la comprensión de expresiones como
\[
  |x-a|<\varepsilon,
\]
donde previamente se han definido los símbolos que aparecen: \(x\) y \(a\) son números reales, \(\varepsilon\) es un número real positivo, \(|x-a|\) denota el valor absoluto de la diferencia \(x-a\), y el símbolo \(<\) representa la relación de orden usual en \(\mathbb{R}\).

Sin embargo, no corresponde a \emph{Fundamentos de Matemáticas} desarrollar con profundidad las definiciones \(\varepsilon\)-\(\delta\) de límite o continuidad. Esas definiciones pertenecen a Análisis Matemático, aunque requieran el lenguaje preparado en este curso.

\begin{advertencia}
El hecho de que un concepto aparezca en \emph{Fundamentos de Matemáticas} no significa que se agote en esta asignatura. Por ejemplo, los números reales se introducen aquí como sistema numérico y como lenguaje de orden, distancia y aproximación; su estudio analítico se desarrolla posteriormente en Análisis Matemático.
\end{advertencia}

\section{Frontera con Álgebra Lineal}

\begin{convencion}[Frontera con Álgebra Lineal]
En \emph{Fundamentos de Matemáticas} se introducen conjuntos, aplicaciones, operaciones internas y las primeras estructuras algebraicas. La teoría de espacios vectoriales, matrices, sistemas lineales, bases, dimensión, transformaciones lineales, determinantes, autovalores y diagonalización corresponde a Álgebra Lineal.
\end{convencion}

La asignatura \emph{Fundamentos de Matemáticas} debe preparar al estudiante para entender expresiones del tipo
\[
  f:A\longrightarrow B,
\]
donde \(A\) y \(B\) son conjuntos y \(f\) es una aplicación de \(A\) en \(B\). También debe introducir conceptos como composición, aplicación identidad, aplicación inversa, inyectividad, sobreyectividad y biyectividad.

Álgebra Lineal utilizará después ese lenguaje para estudiar aplicaciones lineales entre espacios vectoriales. Por tanto, \emph{Fundamentos de Matemáticas} debe evitar desarrollar la teoría específica de espacios vectoriales más allá de mencionarlos como ejemplo de estructura algebraica.

\begin{advertencia}
En \emph{Fundamentos de Matemáticas} puede citarse que un espacio vectorial es una estructura algebraica, pero no deben estudiarse en detalle bases, dimensión, rango, determinantes o diagonalización, porque esos contenidos pertenecen a Álgebra Lineal.
\end{advertencia}

\section{Frontera con Matemática Discreta}

\begin{convencion}[Frontera con Matemática Discreta]
En \emph{Fundamentos de Matemáticas} se estudian los números naturales, la inducción, las definiciones recursivas, conjuntos y relaciones como lenguaje general. Las aplicaciones a combinatoria, probabilidad discreta y grafos pertenecen a Matemática Discreta.
\end{convencion}

La inducción matemática debe aparecer en \emph{Fundamentos de Matemáticas} como método de demostración. Por ejemplo, se puede demostrar que, para todo número natural \(n\geq 1\),
\[
  1+2+\cdots+n=\frac{n(n+1)}{2}.
\]

En cambio, el uso de la inducción para justificar fórmulas combinatorias, propiedades de grafos o algoritmos discretos puede desarrollarse en Matemática Discreta.

\begin{advertencia}
La asignatura \emph{Fundamentos de Matemáticas} no debe absorber contenidos propios de Matemática Discreta. En particular, no debe convertirse en un curso de combinatoria, grafos o probabilidad discreta.
\end{advertencia}

\section{Frontera con Algoritmia}

\begin{convencion}[Frontera con Algoritmia]
En \emph{Fundamentos de Matemáticas II} se estudian procedimientos aritméticos como el algoritmo de Euclides desde el punto de vista matemático: definición, justificación y demostración de corrección. La implementación mediante variables, funciones, estructuras de control, iteración, recursividad y programación corresponde a Algoritmia.
\end{convencion}

Por ejemplo, en \emph{Fundamentos de Matemáticas II} se puede estudiar que, dados dos enteros \(a\) y \(b\), no ambos nulos, el algoritmo de Euclides permite calcular su máximo común divisor. También se puede demostrar que existen enteros \(u\) y \(v\) tales que
\[
  \operatorname{mcd}(a,b)=ua+vb.
\]

Algoritmia podrá después traducir este procedimiento a pseudocódigo o a un lenguaje de programación concreto.

\begin{advertencia}
En \emph{Fundamentos de Matemáticas} se demuestra por qué un procedimiento funciona. En Algoritmia se estudia cómo implementarlo de forma correcta y eficiente.
\end{advertencia}

\section{Frontera con Métodos Numéricos}

\begin{convencion}[Frontera con Métodos Numéricos]
En \emph{Fundamentos de Matemáticas II} se introducen el valor absoluto, la distancia, los intervalos, las desigualdades y la noción elemental de aproximación. El estudio de errores, ceros de funciones, interpolación, diferenciación e integración numérica, así como la resolución numérica de sistemas, pertenece a Métodos Numéricos.
\end{convencion}

El estudiante debe llegar a Métodos Numéricos con soltura en expresiones como
\[
  |x-y|, \qquad [a,b], \qquad x\in (a,b), \qquad |x-a|<\varepsilon.
\]

En \emph{Fundamentos de Matemáticas} estas expresiones se interpretan desde el punto de vista del lenguaje matemático. En Métodos Numéricos se convierten en herramientas para estimar errores, controlar aproximaciones y analizar algoritmos.

\begin{advertencia}
No corresponde a \emph{Fundamentos de Matemáticas} desarrollar métodos de bisección, Newton, interpolación, cuadratura o análisis de error. Sí corresponde preparar el lenguaje de desigualdades, distancia y aproximación que estos métodos necesitan.
\end{advertencia}

\section{Frontera con Estructuras Algebraicas}

\begin{convencion}[Frontera con Estructuras Algebraicas]
En \emph{Fundamentos de Matemáticas} los conceptos de semigrupo, monoide, grupo, anillo y cuerpo se presentan de forma introductoria, como ejemplos de estructuras algebraicas definidas mediante operaciones y axiomas. La teoría sistemática de grupos, subgrupos, homomorfismos, cocientes, acciones, anillos y cuerpos pertenece a la asignatura Estructuras Algebraicas.
\end{convencion}

En \emph{Fundamentos de Matemáticas} basta con que el estudiante comprenda que una estructura algebraica consiste en un conjunto dotado de una o varias operaciones que satisfacen ciertas propiedades. Por ejemplo, puede estudiarse que \((\mathbb{Z},+)\) es un grupo abeliano, donde \(\mathbb{Z}\) es el conjunto de los números enteros y \(+\) es la suma usual.

No es necesario desarrollar de forma extensa subgrupos, morfismos, grupos cociente, acciones de grupo o clasificación de estructuras. Estos temas requieren un tratamiento propio y corresponden a una asignatura posterior.

\begin{advertencia}
La función del capítulo de estructuras algebraicas en \emph{Fundamentos de Matemáticas} es preparar el lenguaje axiomático, no adelantar la asignatura completa de Estructuras Algebraicas.
\end{advertencia}

\section{Frontera con Criptografía}

\begin{convencion}[Frontera con Criptografía]
En \emph{Fundamentos de Matemáticas II} se introducen divisibilidad, números primos, máximo común divisor, identidad de Bézout, congruencias, inversos modulares y ecuaciones lineales módulo un entero. El uso de estas herramientas en sistemas criptográficos, protocolos de clave pública, factorización computacional, logaritmo discreto o funciones hash pertenece a Criptografía.
\end{convencion}

La aritmética modular debe estudiarse en \emph{Fundamentos de Matemáticas II} como parte del lenguaje aritmético básico. Por ejemplo, si \(n\) es un entero positivo y \(a,b\in\mathbb{Z}\), se define
\[
  a\equiv b \pmod n
\]
cuando \(n\) divide a \(a-b\).

Criptografía utilizará después esta base para estudiar algoritmos y protocolos concretos. Por tanto, en \emph{Fundamentos de Matemáticas} no es necesario presentar RSA, Diffie--Hellman, ElGamal o criptografía postcuántica, salvo como motivación cultural breve.

\section{Frontera con Análisis Complejo y de Fourier}

\begin{convencion}[Frontera con Análisis Complejo y de Fourier]
En \emph{Fundamentos de Matemáticas II} se estudian los números complejos como sistema numérico: forma binómica, operaciones, conjugado, módulo, inverso, plano complejo, forma polar y raíces. El estudio de funciones complejas, series de potencias, teoría de Cauchy, residuos y series de Fourier pertenece a Análisis Complejo y de Fourier.
\end{convencion}

En \emph{Fundamentos de Matemáticas II} el objetivo es que el estudiante comprenda expresiones como
\[
  z=a+bi,
\]
donde \(a,b\in\mathbb{R}\), \(i\) es la unidad imaginaria y satisface \(i^2=-1\). También debe saber calcular conjugados, módulos, inversos y raíces sencillas.

El estudio de funciones \(f:\mathbb{C}\to\mathbb{C}\), derivabilidad compleja, integrales complejas o residuos queda fuera del alcance de esta asignatura.

\begin{advertencia}
\emph{Fundamentos de Matemáticas II} estudia números complejos. \emph{Análisis Complejo y de Fourier} estudia funciones de variable compleja y herramientas analíticas avanzadas.
\end{advertencia}

\section{Criterios prácticos de coordinación}

Para que estas fronteras sean efectivas, se recomienda aplicar los siguientes criterios de coordinación docente.

\subsection*{Criterio 1: prioridad del lenguaje}

Si un contenido aparece en \emph{Fundamentos de Matemáticas}, debe tratarse en primer lugar desde el punto de vista del lenguaje: símbolos, definiciones, ejemplos, contraejemplos y demostraciones elementales.

\subsection*{Criterio 2: profundidad controlada}

Un concepto que pertenezca de manera natural a una asignatura posterior puede introducirse en \emph{Fundamentos de Matemáticas}, pero solo hasta el nivel necesario para que el estudiante reconozca su significado y pueda utilizarlo en ejemplos básicos.

\subsection*{Criterio 3: evaluación transversal}

La evaluación de \emph{Fundamentos de Matemáticas} debe centrarse en la lectura y escritura de textos matemáticos, la precisión simbólica, la construcción de ejemplos y contraejemplos, y las demostraciones elementales. No debe evaluar contenidos técnicos propios de otras asignaturas.

\subsection*{Criterio 4: coordinación horizontal}

En primer curso, los contenidos de \emph{Fundamentos de Matemáticas} deben ordenarse de manera que sirvan de apoyo inmediato a Análisis Matemático, Álgebra Lineal, Matemática Discreta, Algoritmia y Métodos Numéricos.

\subsection*{Criterio 5: coordinación vertical}

Los capítulos de estructuras algebraicas, aritmética modular, números reales y números complejos deben presentarse como preparación para asignaturas posteriores: Estructuras Algebraicas, Criptografía, Análisis Matemático, Métodos Numéricos, Topología, Análisis Funcional y Análisis Complejo.

\section{Resumen operativo}

\begin{center}
\begin{tabular}{p{0.28\textwidth}p{0.32\textwidth}p{0.32\textwidth}}
\toprule
\textbf{Tema} & \textbf{En Fundamentos de Matemáticas} & \textbf{En la asignatura específica} \\
\midrule
Inducción & Método de demostración & Aplicaciones en combinatoria, grafos y algoritmos \\
Conjuntos y relaciones & Lenguaje básico y ejemplos finitos & Uso sistemático en estructuras discretas y topológicas \\
Aplicaciones & Dominio, codominio, imagen, composición & Aplicaciones lineales, funciones continuas, funciones computables \\
Estructuras algebraicas & Axiomas básicos y ejemplos & Grupos, anillos, cuerpos, cocientes y homomorfismos \\
Números reales & Orden, intervalos, valor absoluto, distancia & Límites, continuidad, derivación e integración \\
Aritmética modular & Congruencias e inversos modulares & Criptografía y algoritmos aritméticos \\
Complejos & Sistema numérico y raíces & Funciones complejas, Cauchy, residuos y Fourier \\
\bottomrule
\end{tabular}
\end{center}

\resumen{Este apéndice establece las fronteras entre \emph{Fundamentos de Matemáticas} y otras asignaturas del grado. La regla general es que Fundamentos introduce lenguaje, notación, razonamiento y estructuras básicas, mientras que las asignaturas específicas desarrollan las teorías técnicas y sus aplicaciones.}



\iffalse
\chapter{Plantilla de capítulo}

Cada capítulo nuevo podrá desarrollarse siguiendo la siguiente plantilla:

\begin{verbatim}
\chapter{Título del capítulo}

\begin{objetivos}
  \item Objetivo 1.
  \item Objetivo 2.
\end{objetivos}

\section{Motivación}

\section{Definiciones fundamentales}

\section{Ejemplos}

\section{Resultados principales}

\section{Demostraciones modelo}

\begin{erroresfrecuentes}
  \item Error frecuente 1.
\end{erroresfrecuentes}

\begin{seminario}
  \item Actividad 1.
\end{seminario}

\begin{ejerciciospropuestos}
  \item Ejercicio 1.
\end{ejerciciospropuestos}

\resumen{Resumen del capítulo.}
\end{verbatim}
\fi

\backmatter

\renewcommand{\bibname}{Bibliografía inicial recomendada}

\begin{thebibliography}{9}
\addcontentsline{toc}{chapter}{Bibliografía inicial recomendada}

  \bibitem{bourbaki1970}
  N. Bourbaki.
  \emph{Éléments de mathématique}.
  Springer, 1970.
  % 2ª edición revisada (1970), reimpresa por
  % Springer‑Verlag, Berlin–Heidelberg, 2006.
  ISBN 3-540-34034-3.

  \bibitem{vorkurs}
Materiales introductorios de matemáticas universitarias y cursos cero de razonamiento matemático.

\bibitem{hammack}
R. Hammack, \emph{Book of Proof}.

\bibitem{velleman}
D. J. Velleman, \emph{How to Prove It: A Structured Approach}.

\bibitem{rosen}
K. H. Rosen, \emph{Discrete Mathematics and Its Applications}.

\bibitem{niven}
I. Niven, H. S. Zuckerman, H. L. Montgomery, \emph{An Introduction to the Theory of Numbers}.

\end{thebibliography}

\end{document}
